\chapter*{空间直线与平面}
\section*{平面及其基本性质}
\subsection*{平面及其基本性质 (1)}

A 组：
\begin{enumerate}
    \item 几何中平面的特征:平面没有\blkda{薄厚}，是\blkda{无限延展}的.
    
    \item 判断下列命题是否正确，正确的打 “ \ding{51} ”，错误的打 “ \ding{55} ”
    \begin{enumerate}
        \item 平行四边形就是平面. \blkda{\ding{55}}
        
        \item 一条直线把它所在的平面分成两部分，一个平面把空间分成两部分. \blkda{\ding{51}}
        
        \item 若平面 $a$ 的面积是 ${20}{m}^{2}$ ，平面 $\beta$ 的面积是 ${10}{m}^{2}$ ，则平面 $a$ 的面积是平面 $\beta$ 的面积的 2 倍 \blkda{\ding{55}}
        
        \item 五个平面重叠在一起比两个平面重叠在一起要厚. \blkda{\ding{55}}
        
        \item 如果一条直线上所有的点都在某一个面上, 那么这个面一定是平面. \blkda{\ding{55}}
        
        \item 经过面内任意两点的直线，若直线上各点都在这个面内，那么这个面是平面. \blkda{\ding{51}}
    \end{enumerate}

    \item 如图所示，用集合的语言表述下列语句:
    
    \begin{center}
    \includegraphics[width=0.2\textwidth]{bodyxb/src/bo_d4g0j5k601uc73fkjrb0_0_1078_1165_323_175_0.jpg}
    \end{center} 
    \begin{dingautolist}{172}
        \item 点 $A\text{ 、 }B$ 在直线 $a$ 上:\blkda[4]{$A\in a , B \in a$}
        
        \item 直线 $a$ 在平面 $\alpha$ 内:\blkda[4]{$a \subset \alpha$}
        
        \item 点 $C$ 在平面 $\alpha$ 内:\blkda[4]{$C \in \alpha$}

        \item 点 $D$ 不在平面 $\alpha$ 内:\blkda[4]{$D \notin \alpha$}
        
        \item 直线 $b$ 不在平面 $\alpha$ 内:\blkda[4]{$b \not \subset \alpha$}.
    \end{dingautolist}
\end{enumerate}

B 组：
\begin{enumerate}
    \item 四个平面将空间最少分成 $m$ 部分，最多分成 $n$ 部分，则 $m + n =$ \blkda{20}.
    \begin{jieda}
        $m = 15, n = 5$
    \end{jieda}
    
    \item 不都在一条直线上的四点可以确定\blkda{1或4}个平面.
    \begin{jieda}
        不存在三个点共线时有4个面（正四面体）

        存在三个点共线时根据公理2的推论1可知有1个面
    \end{jieda}

    \item 一条直线和直线外三点可以确定\blkda[2]{1或3或4}个平面.
    \begin{jieda}
        三点共线时可以确定1个面或者3个面

        三个点不共线时根据公理2可知这三个点确定一个平面$\alpha$，若直线在$\alpha$内则只确定一个平面，否则，直线和三个点会确定$\beta_1, \beta_2, \beta_3$三个平面，因此最多确定4个平面，若直线和某两点的连线有交点，则会使得$\beta_i$与$\beta_j$重合，退化为3个平面.
    \end{jieda}

    \item 两两相交于一点的三条直线可以确定\blkda{1或3}个平面.
    \begin{jieda}
        两条相交直线确定一个平面，若第三个直线不在该面内，则确定三个平面，否则确定一个平面
    \end{jieda}

    \item 两两平行的四条直线可以确定\blkda[2]{1或4或6}个平面.
    \begin{jieda}
        四条直线共面则可以确定一个平面

        三条直线共面，另一条不在面内则确定四个平面

        任意三条直线都不共面则确定六个平面
    \end{jieda}

    \item 由正方体的一条棱和一个不在此棱上的顶点确定的平面有\blkda{12}个.

    \item 下列判断中: 
    \begin{dingautolist}{172}
        \item 三点确定一个平面
        \item 一条直线和一点确定一个平面
        \item 两条直线确定一个平面
        \item 三角形和梯形一定是平面图形
        \item 四边形一定是平面图形
        \item 六边形一定是平面图形
        \item 两两相交的三条直线确定一个平面
    \end{dingautolist}
    其中正确的是\blkda{\ding{175}}.
    
    \item 空间三个平面可以把空间分成\dotfill{(\kh{B})}
    \xx{4 或 6 或 7 个部分}{4 或 6 或 7 或 8 个部分}{4 或 7 或 8 个部分}{6或 7 或 8 个部分}
    
    \item 若直线上有两个点在平面外, 则\dotfill{(\kh{D})}
    \xx{直线上至少有一个点在平面内}{直线上有无穷多个点在平面内}{直线上所有点都在平面外}{直线上至多有一个点在平面内}
    
    \item 已知命题 “直线 $l$ 上两点 $A\text{ 、 }B$ 在平面 $a$ 上”,则下列命题中与其不等价的是\dotfill{(\kh{C})}
    \xx{$l \subset  a$}{平面 $a$ 过直线 $l$}{直线 $l$ 上只有两点在平面 $a$ 上}{直线 $l$ 上所有点都在平面 $a$ 上}
    
    \item 若点 $M$ 在直线 $b$ 上， $b$ 在平面 $\beta$ 上，则 $M$ 、 $b$ 、 $\beta$ 之间的关系可记作\dotfill{(\kh{C})}
    \xx{$M \in  b \in  \beta$}{$M \subset  b \subset  \beta$}{$M \in  b \subset  \beta$}{$M \subset  b \in  \beta$}
    
    \item 请将以下四图中，看得见的部分用实线描出.
    \fourtu{\includegraphics[width=0.6\textwidth]{bodyxb/src/bo_d4g0j5k601uc73fkjrb0_1_258_1352_241_197_0.jpg}}{\includegraphics[width=0.6\textwidth]{bodyxb/src/bo_d4g0j5k601uc73fkjrb0_1_557_1363_304_177_0.jpg}}{\includegraphics[width=0.6\textwidth]{bodyxb/src/bo_d4g0j5k601uc73fkjrb0_1_916_1355_247_199_0.jpg}}{\includegraphics[width=0.6\textwidth]{bodyxb/src/bo_d4g0j5k601uc73fkjrb0_1_1227_1357_225_194_0.jpg}}
    
    \item 画图表示下列集合语言所描述的关系:
    \begin{enumerate}
        \item $A \in  \alpha ,a \subset  \alpha ,A \notin  a$
        \item $a \cap  \beta  = 1,A \in  \alpha$ ，且 $A \in  \beta$
        \item $a,b,c \subset  \alpha$ ,且 $a \cap  b = A,b \cap  c = B,c \cap  a = C$
        \item $A \in  \alpha, A \in  l,l \cap  \beta  = B, \alpha \cap  \beta  = m,B \in  m$
        \item $a \cap  \beta  = \varnothing ,l \cap  a = A,l \cap  \beta  = B$ 
        \item ${l}_{1} \subset  \alpha,{l}_{2} \cap  \alpha = M,M \notin  {l}_{1}$
    \end{enumerate}
    
    \item 已知直线 $a$ 和直线 $b$ 相交于点 $M$ ，点 $M$ 在平面 $\alpha$ 上，试用集合语言描述上述位置关系, 并画出图形.
    \begin{jieda}
        $a \cap b = M, M \in \alpha$
    \end{jieda}
    
    \item 已知 $I$ 为 $\bigtriangleup  {ABC}$ 的内心. 若 $A$ 、 $B$ 、 $C \in  \alpha$ ，利用公理 1 证明: $I \in  \alpha$ .
    \begin{jieda}
        连接$AI$交$BC$于$M$，根据公理1可知$M \in \alpha$，故$AM \subset \alpha$，因此$I \in \alpha$
    \end{jieda}
    
    \item 用反证法证明: 一个平面 $\alpha$ 与不在 $\alpha$ 上的一条直线 $l$ 最多只能有一个交点.
    \begin{jieda}
        假设$\alpha$与$l$存在两个及以上的公共点，设其中两个为$A,B$，则$A, B \in \alpha$，根据公理1可知$l \subset \alpha$，矛盾，因此一个平面 $\alpha$ 与不在 $\alpha$ 上的一条直线 $l$ 最多只能有一个交点.
    \end{jieda}
\end{enumerate}

\clearpage
\subsection*{平面及其基本性质 (2)}

A 组：
\begin{enumerate}
    \item 判断下列命题是否正确，正确的打 “ \ding{51} ”，错误的打 “ \ding{55} ”
    \begin{dingautolist}{172}
        \item 若 $\bigtriangleup  {ABC}$ 的两边 ${AB}$ 、 ${AC}$ 在平面 $\alpha$ 上，则它的第三边 ${BC}$ 也在平面 $a$ 上\blkda{\ding{51}}
        \item 若两个平面有不同的三个公共点，则这两个平面重合\blkda{\ding{55}}
        \item 若四点不共面，则每三个点一定不共线\blkda{\ding{51}}
        \item 已知三个平面两两相交，且有三条交线. 如果这三条交线中有两条交线平行，那么第三条交线也与另两条直线平行\blkda{\ding{51}}
        \item 四边相等的四边形是平面图形\blkda{\ding{55}}
    \end{dingautolist}
    
    \item 已知 $a \cap  \beta  = l,m \subset  a,n \subset  \beta ,m \cap  n = P$ ，则点 $P$ 与直线 $l$ 的位置关系用相应的符号表示为\blkda{$P \in l$}.
    \begin{jieda}
        两平面的公共点都在交线上.
    \end{jieda}
\end{enumerate}

B 组：
\begin{enumerate}
    \item 下列命题
    \begin{dingautolist}{172}
        \item 若空间四点共面，则其中必有三点共线
        \item 若空间四点中有三点共线，则此四点必共面
        \item 若空间四点中任何三点不共线，则此四点不共面
        \item 若空间四点不共面，则其中任意三点不共线
    \end{dingautolist}
    其中真命题的个数是\dotfill{(\kh{B})}
    \xx{1 个}{2 个}{3 个}{4 个}
    \begin{jieda}
        \begin{dingautolist}{172}
            \item 四点共面不一定得到三点共线，比如平面四边形，故错误
            \item 若空间四点中有三点共线，则此四点必共面，故正确
            \item 若空间四点中任何三点不共线，则此四点可能共面，比如平面四边形，故错误
            \item 若空间四点不共面，则其中任意三点不共线；假设其中三个点共线，则第四个点要么与这三点在一条直线上，要么在直线外，根据直线和直线外一点可确定一个平面可知，这四点共面，矛盾，故任意三点不共线；故正确
        \end{dingautolist}
    \end{jieda}

    \item 在正方体中， ${E}$ 、 ${F}$ 、 ${G}$ 、 ${H}$ 分别是该点所在棱的中点，则下列图形中 ${E}$ 、 ${F}$ 、 ${G}$ 、 ${H}$ 四点共面的是\dotfill{(\kh{B})}
    \xx{\includegraphics[width=0.3\textwidth]{bodyxb/src/bo_d4g0jos601uc73fkjrl0_0_291_1650_208_227_0.jpg}}{\includegraphics[width=0.3\textwidth]{bodyxb/src/bo_d4g0jos601uc73fkjrl0_0_583_1665_236_214_0.jpg}}{\includegraphics[width=0.3\textwidth]{bodyxb/src/bo_d4g0jos601uc73fkjrl0_0_847_1665_239_211_0.jpg}}{\includegraphics[width=0.3\textwidth]{bodyxb/src/bo_d4g0jos601uc73fkjrl0_0_1165_1633_213_242_0.jpg}}
    \begin{jieda}
        对于选项$A$，如下图，点$E$、$F$、$H$、$M$确定一个平面，该平面与底面交于$FM$，而点$G$不在平面$EHMF$上，故$E$、$F$、$G$、$H$四点不共面
        \begin{center}                  \includegraphics[width=0.2\linewidth]{bodyxb//src/pingmian(2)B3A.png}
        \end{center}

        对于选项$B$，连结底面对角线$AC$，由中位线定理得$FG//AC$，又$EH//AC$，则$EH//FG$，故$E$、$F$、$G$、$H$四点共面
        \begin{center}                  \includegraphics[width=0.2\linewidth]{bodyxb//src/pingmian(2)B3B.png}
        \end{center}

        对于选项$C$，显然$E$、$F$、$H$所确定的平面为正方体的底面，而点$G$不在该平面内，故$E$、$F$、$G$、$H$四点不共面

        对于选项$D$，如图，取部分棱的中点，顺次连接，得一个正六边形，即点$E$、$G$、$H$确定的平面，该平面与正方体正面的交线为$PQ$，而点$F$不在直线$PQ$上，故$E$、$F$、$G$、$H$四点不共面．
        \begin{center}                  \includegraphics[width=0.2\linewidth]{bodyxb//src/pingmian(2)B3D.png}
        \end{center}
        \end{jieda}
        
    \item 下列说法正确的是\dotfill{(\kh{D})}
    \xx{线段 AB 在平面 $\alpha$ 内,直线 AB 不在 $\alpha$ 内}{平面 $\alpha$ 和 $\beta$ 有时只有一个公共点}{空间三点可以确定一个平面}{若两个平面有不共线的三个公共点, 则两个平面就重合}
    
    \item 在空间四点中，无三点共线是四点共面的\dotfill{(\kh{D})}
    \xx{充分不必要条件}{必要不充分条件}{充分必要条件}{既不充分又不必要条件}
    
    \item 如图，已知 $A\text{ 、 }B\text{ 、 }C$ 在正方体的顶点或棱上，画出过 $A\text{ 、 }B\text{ 、 }C$ 的正方体的截面(保留作图痕迹).
    
    \begin{center}
    \includegraphics[width=0.9\textwidth]{bodyxb/src/bo_d4g0jos601uc73fkjrl0_1_248_784_1146_1046_0.jpg}
    \end{center}

    \begin{jieda}
        \begin{center}
    
            \includegraphics[width=1\linewidth]{asfdbajsfa.png}
        \end{center}
    \end{jieda}
    
    \item 如图， $A\text{ 、 }P$ 是边长为 1 的正方体的顶点， $B\text{ 、 }C$ 在此正方体的棱上，且 $B$ 是所在棱的中点，根据 ${PC}$ 的长度分类讨论过 $A\text{ 、 }B\text{ 、 }C$ 的正方体截面的形状(三角形、四边形、 五边形或者六边形).
    
    \begin{center}
    \includegraphics[width=0.2\textwidth]{bodyxb/src/bo_d4g0jos601uc73fkjrl0_2_209_441_332_313_0.jpg}
    \end{center}

    \begin{jieda}
        ${PC} \in  \left\lbrack  {0,\dfrac{1}{2}}\right\rbrack   \cup   1\}$ : 四边形
        
        ${PC} \in  \left( {\dfrac{1}{2},1}\right)$ : 五边形
    \end{jieda}
    
    \item 如图，已知 $a \cap  \beta  = {EF},A \in  a,{C}$ 、 $B \in  \beta ,{BC}$ 与 ${EF}$ 相交，在图中画出平面 ${ABC}$ 分别与 $\alpha$ 、 $\beta$ 的交线.
    
    \begin{center}
    \includegraphics[width=0.2\textwidth]{bodyxb/src/bo_d4g0jos601uc73fkjrl0_2_242_989_344_217_0.jpg}
    \end{center}
    
    \begin{jieda}
        \begin{center}
            \includegraphics[width=0.25\linewidth]{bodyxb//src/pingmian(2)7.png}
        \end{center}
    \end{jieda}

    \item 已知: 直线 $a\text{ 、 }b\text{ 、 }c$ 相交于 ${P}$ 点,直线 $d$ 与 $a\text{ 、 }b\text{ 、 }c$ 分别相交于 ${A}\text{ 、 }{B}\text{ 、 }{C}$ 三点. 求证: 直线 $a$ 、 $b$ 、 $c$ 与直线 $d$ 共面.
    
    \begin{center}
    \includegraphics[width=0.3\textwidth]{bodyxb/src/bo_d4g0jos601uc73fkjrl0_2_257_1432_381_250_0.jpg}
    \end{center}

    \begin{jieda}
        根据公理2的推论可知$d$和$P$确定一个平面$\alpha$，则$A,B,C,P$均在$\alpha$内，根据公理1可知$a, b, c$均在$\alpha$内.
    \end{jieda}


    \item 如图， $\bigtriangleup  {{ABC}}$ 和 $\bigtriangleup  {{A}}_{1}{{B}}_{1}{{C}}_{1}$ 不在同一个平面内，如果三直线 ${A}{{A}}_{1}\text{ 、 }{B}{{B}}_{1}\text{ 、 }{{{CC}}}_{1}$ 两两相交，求证: 三直线 $A{A}_{1}\text{ 、 }B{B}_{1}\text{ 、 }C{C}_{1}$ 交于一点.

    \begin{center}
    \includegraphics[width=0.3\textwidth]{bodyxb/src/bo_d4g0jos601uc73fkjrl0_3_246_337_358_238_0.jpg}
    \end{center}

    \begin{jieda}
        设 $B{B}_{1}$ 与 $C{C}_{1},C{C}_{1}$ 与 $A{A}_{1},A{A}_{1}$ 与 $B{B}_{1}$ 分别确定平面 $\alpha ,\beta ,\gamma ,A{A}_{1}$ 与 $B{B}_{1}$ 的交点为 $P$
        
        因为 $P \in  A{A}_{1},P \in  B{B}_{1}$，$A{A}_{1} \subset  \beta ,B{B}_{1} \subset  \alpha$ ,所以 $P \in  \alpha ,P \in  \beta$ ,即 $P \in  \alpha  \cap  \beta$ .
        
        又 $\alpha  \cap  \beta  = C{C}_{1}$ ,所以 $P \in  C{C}_{1}$ .所以三条直线 $A{A}_{1},B{B}_{1},C{C}_{1}$ 交于一点
    \end{jieda}

\end{enumerate}

C 组:
\begin{enumerate}
    \item 到空间不共面的 4 个点距离都相等的平面有\dotfill{(\kh{A})}
    \xx{7 个}{6 个}{4 个}{3 个}
    \begin{jieda}
        将空间不共面的四个点看作四面体的顶点,
        
        第一类:平行于底面，且过测棱的中点，此时由 4 种情况符合要求，如图:
        
        \begin{center}
        \includegraphics[width=0.2\textwidth]{bodyxb/src/bo_d4hptm3ef24c739bif30_0_243_609_265_301_0.jpg}
        \end{center}
        
        第二类:一条棱的中点于底面的一条中位线所成的平面，此时有 3 种情况符合要求，如图
        
        \begin{center}
        \includegraphics[width=0.2\textwidth]{bodyxb/src/bo_d4hptm3ef24c739bif30_0_243_978_265_302_0.jpg}
        \end{center}
        
        故总共有 7 种情况符合要求.
    \end{jieda}
    
    \item 空间有 6 个点，其中任意三点不共线，且有五个点共面，则这 6 个点最多可以确定\blkda{11} 个平面.

    \begin{jieda}
        记 5 个共面的点分别为 $a,b,c,d,e$ ，不共面的第 6 个点为 $O$ ，则从 $a,b,c,d,e$ 中随机选取两个点，连同第 6 个点，即可构成一个平面，此时共有
        
        $\left( {abO}\right) ,\left( {acO}\right) ,\left( {adO}\right) ,\left( {aeO}\right) ,\left( {bcO}\right) ,\left( {bdO}\right) ,\left( {beO}\right) ,\left( {cdO}\right) ,\left( {ceO}\right)$ 共有 10 个平面，(abcde) 这 5 个点确定一个平面，故最多可以确定 10+1=11 个平面
    \end{jieda}
    
    \item 金字塔在埃及和美洲等地均有分布，现在的尼罗河下游，散布着约 80 座金字塔遗迹，大小不一, 其中最高大的是胡夫金字塔, 如图, 胡夫金字塔可以近似看做一个正四棱锥, 则该正四棱锥的 5 个面所在的平面将空间分成\blkda{23}部分 (用数字作答).
    
    \begin{center}
    \includegraphics[width=0.3\textwidth]{bodyxb/src/bo_d4g0jos601uc73fkjrl0_3_244_1253_400_272_0.jpg}
    \end{center}
    \begin{jieda}
        假想一个没有上顶的正方体，该正方体会把空间分割成 18 块，把四面进行极限倾斜相交，如图所示，
        
        \begin{center}
        \includegraphics[width=0.2\textwidth]{bodyxb/src/bo_d4hptm3ef24c739bif30_1_243_889_281_260_0.jpg}
        \end{center}
        
        在倾斜的过程中，在不管底面的情况下，4 个侧面在顶点以下的“水平范围”内最多可以切割出 9 个空间，与没有倾斜极限的情况一样，多出来的空间是交叉的切割出来的空间，
        
        在空间上是对称的，四个倾斜的侧面在空间中的延伸还是这样的倾斜侧面，
        
        如图所示的对称的锥面同样会切割出 $g$ 个空间，即顶点之上的 4 个延伸的倾斜的面同样会切割出 9 个空间.
        
        但是四个空间和下面的四个倾斜的侧面切出的是同一个，即标记“ $\times$ ”的位置，
        
        \begin{center}
        \includegraphics[width=0.3\textwidth]{bodyxb/src/bo_d4hptm3ef24c739bif30_1_239_1681_359_281_0.jpg}
        \end{center}
        
        所以在 18 的基础上加 9 减 4 ，即结果是 23.
    \end{jieda}

\end{enumerate}

\clearpage
\subsection*{平面及其基本性质 (3)}
\begin{enumerate}
    \item 两个相等的角在直观图中\blkda{不一定}相等；两条相等的线段在直观图中\blkda{不一定} 相等. (填 “一定” “不一定”)

    \item 如果正三角形 ${ABC}$ 的边长为 $a$ ，那么 $\bigtriangleup  {ABC}$ 的水平放置的直观图 $\bigtriangleup  {A}^{\prime }{B}^{\prime }{C}^{\prime }$ 的面积为\blkda{$\dfrac{\sqrt{6}}{16}a^2$}，已知 $\bigtriangleup  {ABC}$ 的平面直观图 $\bigtriangleup  {A}^{\prime }{B}^{\prime }C$ 是边长为 $a$ 的正三角形，那么原$\bigtriangleup  {ABC}$ 的面积为\blkda{$\dfrac{\sqrt{6}}{2}a^2$}.

    \item 若右图表示水平放置的 $\bigtriangleup  {ABC}$ 的直观图(其中 $D$ 为 ${BC}$ 的中点)，则在原来的 ${\Delta ABC}$ 中，四条线段中最长的是\blkda{AC}.
    
    \begin{center}
    \includegraphics[width=0.2\textwidth]{bodyxb/src/bo_d4g0jtbef24c73ek1ku0_0_393_724_276_204_0.jpg}
    \end{center}
    
    \item 一梯形的直观图是一个如图所示的等腰梯形，且该梯形面积为 $\sqrt{2}$ ，则原梯形的面积为\blkda{4}.
    
    \begin{center}
    \includegraphics[width=0.2\textwidth]{bodyxb/src/bo_d4g0jtbef24c73ek1ku0_0_996_770_245_152_0.jpg}
    \end{center}

    \item 利用斜二测画法得到的:
    \begin{dingautolist}{172}
        \item 三角形的直观图一定是三角形； \item 正方形的直观图一定是菱形；
        \item 等腰梯形的直观图可以是平行四边形； \item 菱形的直观图一定是菱形.
    \end{dingautolist}
    以上结论正确的是\blkda{\ding{172}}.
    
    \item 如图，模块 \ding{172}-\ding{176} 均由 4 个棱长为 1 的小正方体构成，模块 \ding{177} 由 15 个棱长为 1 的小正方体构成. 现从模块 \ding{172}-\ding{176} 中选出三个放到模块 \ding{177} 上，使得模块 \ding{177} 成为一个棱长为 3 的大正方体. 则十种可能的选择方案中，能够完成任务的为\blkda[4]{\ding{172}\ding{173}\ding{176}或\ding{172}\ding{175}\ding{176}或\ding{173}\ding{174}\ding{175}}.
    
    \begin{center}
    \includegraphics[width=0.8\textwidth]{bodyxb/src/bo_d4g0jtbef24c73ek1ku0_0_241_1690_981_223_0.jpg}
    \end{center}

    \item 若长方体的全面积为 22，棱长之和为 24，则其体对角线的长为\dotfill{(\kh{D})}
    \xx{$\sqrt{11}$}{$\sqrt{12}$}{$\sqrt{13}$}{$\sqrt{14}$}
    
    \item 在正方体 ${ABCD} - {A}_{1}{B}_{1}{C}_{1}{D}_{1}$ 中，点 $P$ 在侧面 ${BC}{C}_{1}{B}_{1}$ 及其边界上运动，并且总是保持 ${AP} \bot  {B{D}_{1}}$ ，则动点 $P$ 的轨迹是\dotfill{(\kh{B})}
    \xx{一个点}{线段}{正三角形}{圆}
    
    \item 在从长方体同一顶点出发的三条棱上各取一点，过它们作截面，所得截面的形状一定是\dotfill{(\kh{C})}
    \xx{等边三角形}{直角三角形}{锐角三角形}{钝角三角形}
    
    \item 在水平放置的平面 $a$ 上有一个四边形,当用斜二测画法画出它的直观图时, 它被画成边长为 $2{{cm}}$ 的正方形，并且底边在水平位置，如图所示. 画出这个四边形的真实形状，并求它的面积.
    
    \begin{center}
    \includegraphics[width=0.4\textwidth]{bodyxb/src/bo_d4g0jtbef24c73ek1ku0_1_872_990_521_194_0.jpg}
    \end{center}

    \begin{jieda}
        $S = 8\sqrt{2} cm^2$
    \end{jieda}
    
    \item 如图，已知四边形 ${OABC}$ ，画出它的直观图.
    
    \begin{center}
    \includegraphics[width=0.2\textwidth]{bodyxb/src/bo_d4g0jtbef24c73ek1ku0_1_209_1338_343_337_0.jpg}
    \end{center}

    \item 在水平放置的平面 $a$ 上，画一个相邻两边长分别为 $2{{cm}}$ 、 $3{{cm}}$ ，且它们的夹角为 ${60}^{ \circ  }$ 的平行四边形的直观图.

    \item 右图是某个水平放置的平面图形的直观图，画出原来的平面图形.
    
    \begin{center}
    \includegraphics[width=0.3\textwidth]{bodyxb/src/bo_d4g0jtbef24c73ek1ku0_2_948_325_459_269_0.jpg}
    \end{center}
\end{enumerate}


\section*{直线与直线的位置关系}
\subsection*{直线与直线的位置关系（1）}

A 组：
\begin{enumerate}

    \item 如果一个角的两边和另一个角的两边分别平行, 那么这两个角的
关系是\blkda[3]{相等或互补}.
    
    
    \item 已知空间四边形 $ {ABCD} $ (即 $ {A}\text{、}{B}\text{、}{C}\text{、}{D} $ 四点不共面) 各边中点
分别为 $ E $ 、 $ F $ 、 $ G $ 、 $ H $ ,则四边形 $ {EFGH} $ 是\blkda[2]{平行四边}形.

\begin{center}
    \centering
    \includegraphics[width=0.25\linewidth]{imagiasfygiuahgfe.png}
\end{center}

    \item 判断下列命题的真假,真的打“√”,假的打“×”：
    \begin{enumerate}
        \item 平行于同一直线的两条直线平行 \blkda{\ding{51}}
        \item 过直线外一点,有且只有一条直线与已知直线平行 \blkda{\ding{51}}
        \item 与已知直线平行且距离等于定长的直线只有两条 \blkda{\ding{55}}
        \item 若一个角的两边分别与另一个角的两边平行,那么这两个角相等\blkda{\ding{55}}
        \item 若两条相交直线和另两条相交直线分别平行,那么这两组直线所成的锐角(或直角)相等 \blkda{\ding{51}}

    \end{enumerate}

\end{enumerate}

B组：
\begin{enumerate}
    \item 若 $ {AB}//{A}^{\prime }{B}^{\prime },{BC}//{B}^{\prime }{C}^{\prime } $ ,且 $ \angle {ABC} = {45}^{ \circ  } $ ,则 $ \angle {A}^{\prime }{B}^{\prime }{C}^{\prime } $ 等于\blkda[2]{$45^{\circ}$或$135^{\circ}$}.
    
    \item 下列命题中,正确的是\blkda{(2)(4)}. (在横线上填写所有正确说法的编号)
    \begin{enumerate}
        \item 有三个角是直角的四边形是矩形；
        \item 有四个角是直角的四边形是矩形；
        \item 四边相等且四个角也相等的四边形一定是正方形；
        \item 两条对角线互相平分的四边形是平行四边形.
    \end{enumerate}
    

    \item 已知空间中两个平行四边形 $ {ABCD} $ 和 $ {EFGH} $ ,若 $ {AB}\overset{//}{ = }{EF},{BC}\overset{//}{ = }{FG} $ ,则下列说法中一定正确的是\blkda{(1)(4)}.
    \begin{enumerate}
        \item  $ {S}_{ABCD} = {S}_{EFGH} $ ；
        \item $ {AC} = {EG} $ ；
        \item 直线 $ {BD} $ 与直线 $ {FH} $ 可以确定一个平面；
        \item 直线 $ {CD} $ 和直线 $ {GH} $ 可以确定一个平面.
    \end{enumerate}

    \item 若空间中四条直线 $ {l}_{1} $ 、 $ {l}_{2} $ 、 $ {l}_{3} $ 、 $ {l}_{4} $ 满足 $ {l}_{1}\bot {l}_{2} $ , $ {l}_{2}//{l}_{3} $ , $ {l}_{3}\bot {l}_{4} $ ,则下列结论一定正确的是\blkda{4,6}.\\
    (1) $ {l}_{1}//{l}_{3} $ 、(2) $ {l}_{1}//{l}_{4} $ 、(3) $ {l}_{2}//{l}_{4} $ 、(4) $ {l}_{1}\bot {l}_{3} $ 、(5) $ {l}_{1}\bot {l}_{4} $ 、(6) $ {l}_{2}\bot {l}_{4} $ .
    
    \item 如图所示,在空间四边形 $ {{ABCD}} $ 中,E、F、G、H 分别是 $ {{AB}} $ 、 $ {{BC}} $ 、 $ {{CD}} $ 、 $ {{DA}} $的中点.
    \begin{enumerate}
        \item 求证：EFGH 是平行四边形；
        \begin{center}
    
            \includegraphics[width=0.25\linewidth]{fhyauhfa'uoeffa.png}
        \end{center}
        \begin{jieda}
            因为$ EH \overset{/\!/}{=\!\!=} \dfrac{1}{2}BD,FG \overset{/\!/}{=\!\!=} \dfrac{1}{2}BD$,所以$EH \overset{/\!/}{=\!\!=} FG$,
            故为平行四边形.
        \end{jieda}

        \item 若 $ {EFGH} $ 为菱形,则 $ {AC} $ 、 $ {BD} $ 满足的条件为\blkda{$AC=BD$}
        
        \item 若 $ {{EFGH}} $ 为矩形,则 $ {AC} $ 、 $ {BD} $ 满足的条件为\blkda{$AC// BD$}
        
        \item 若 $ {EFGH} $ 为正方形,则 $ {AC} $ 、 $ {BD} $ 满足条件为\blkda[4]{$AC// BD$且$AC=BD$}
        
        \item 若 $ {{AC}} + {{BD}} = a $ , $ {{AC}} \cdot  {{BD}} = b $ ,则 $ E{G}^{2} + H{F}^{2} = $ \blkda{$\dfrac{1}{2}a^2-b$}
    \end{enumerate}
    
    \item 如图,在空间四边形 $ {{ABCD}} $ 中,点 $ {E} $ 、 $ {H} $ 分别是边 $ {{AB}} $ 、$ {{AD}} $ 上的点, $ {F} $ 、 $ {G} $ 分别是边 $ {{BC}} $ 、 $ {{CD}} $ 上的点. 若点 $ {E} $ 、 $ {H} $ 分别是边 $ {{AB}} $ 、 $ {{AD}} $ 的中点,且 $ \dfrac{CF}{CB} = \dfrac{CG}{CD} = \dfrac{2}{3} $ ,则
        \dotfill{(\kh{D})}
        \begin{center}
    
            \includegraphics[width=0.25\linewidth]{imaafha'oisfge.png}
        \end{center}
    \xx{EF与 $ {GH} $ 互相平行;}
    {EF与GH 异面;}
    {EF与 $ {GH} $ 的交点 $ {M} $ 可能在直线 $ {{AC}} $ 上,也可能不在直线 $ {{AC}} $ 上；}
    {EF与 $ {GH} $ 的交点 $ M $ 一定在直线 $ {AC} $ 上.}
    \begin{jieda}
        因为$EF\subset 平面ABC，HG\subset 平面ACD$，由公理3可知交点$M$一定在两个面的交线$AC$上.
    \end{jieda}

   
    \item 下列命题中,正确结论有\blkda{2,3}.
    \begin{enumerate}
        \item 如果一个角的两边与另一个角的两边分别平行,那么这两个角相等；
        \item 如果两条相交直线和另两条相交直线分别平行,那么这两组直线所成的锐角 (或直角) 相等;
        \item 如果两条直线同平行于第三条直线,那么这两条直线互相平行.     
    \end{enumerate}

    \item 如果空间四边形的对角线相等,那么顺次联结它的各边中点所成的四边形一定是\dotfill{(\kh{C})}
    \xx{梯形}{长方形}{菱形}{正方形}

    
    \item 在空间四边形 $ {ABCD} $ 中, $ P $ 、 $ Q $ 分别是边 $ {AB} $ 、 $ {CD} $ 的中点,设 $ {BD} + {AC} = {2a} $ ,则\dotfill{(\kh{C})}
    \xx{$ {PQ} > a $}
    {$ {PQ} = a $}
    {$ {PQ} < a $}
    {$ {PQ} $ 与 $ a $ 的大小不确定} 
    
    \item 在空间四边形 $ {{ABCD}} $ 的边 $ {{AB}} $ 、 $ {{BC}} $ 、 $ {{CD}} $ 、 $ {{DA}} $ 上分别取 $ {E} $ 、 $ {F} $ 、 $ {G} $ 、 $ {H} $ 四点,如果 $ {{EF}} $ 与 $ {{HG}} $ 交于点 $ {M} $ ,则 \dotfill{(\kh{A})}
    \xx{M一定在直线 AC 上}
    {M一定在直线 $ {{BD}} $ 上}
    {$ {M} $ 可能在 $ {{AC}} $ 上,也可能在 $ {{BD}} $ 上}
    {$ {M} $ 不在 $ {{AC}} $ 上,也不在 $ {{BD}} $ 上} 

    \item 如图,在长方体 $ {ABCD} - {A}_{1}{B}_{1}{C}_{1}{D}_{1} $ 中, $ M $ 为 $ {B}_{1}B $ 上一点,试过点 $ M $ 作平面 $ {D}_{1}{DC}{C}_{1} $的对角线 $ D{C}_{1} $ 的平行线,并说明理由.

    \begin{center}

        \includegraphics[width=0.25\linewidth]{fah'lushdf'afas.png}
    \end{center}

    \begin{jieda}
        在$AB$上找作点$N$，使得$MN// AB_1$，根据直线平行的传递性可知$MN// DC_1$
    \end{jieda}


    \item 如图, $ \bigtriangleup {{ABC}} $ 和 $ \bigtriangleup {{A}}_{1}{{B}}_{1}{{C}}_{1} $ 的对应顶点连线, $ {{{AA}}}_{1} $ 、 $ {{{BB}}}_{1} $ 、$ {{{CC}}}_{1} $ 交于点 $ {O},{O} $ 在平面 $ {{ABC}} $ 和平面 $ {{A}}_{1}{{B}}_{1}{{C}}_{1} $ 之间,且$ \dfrac{{{AO}}}{{{A}}^{\prime }{O}} = \dfrac{{{BO}}}{{{B}}^{\prime }{O}} = \dfrac{{{CO}}}{{{C}}^{\prime }{O}} = \dfrac{2}{3}. $
    \begin{enumerate}
        \item 求证： $ \angle {CAB} = \angle {C}_{1}{A}_{1}{B}_{1} $ ；
        \item 求: $ {S}_{\bigtriangleup {ABC}} : {S}_{\bigtriangleup {A}_{1}{B}_{1}{C}_{1}} $ 的值.
    \end{enumerate}
    \begin{center}

        \includegraphics[width=0.25\linewidth]{f'ahsufa'fad.png}
    \end{center}
\end{enumerate}

    \begin{jieda}
        \begin{enumerate}
            \item 因为$\triangle ACO \sim \triangle A_1C_1O$，所以$AC// A_1C_1$,因为$\triangle ABO \sim A_1B_1O$，所以$AB// A_1B_1$.由等角定理可得，$\angle{CAB}=\angle{C_1A_1B_1}$

            \item $\triangle ABC \sim \triangle A_1B_1C_1$,因此面积之比是边长之比的平方，故4:9.
            
        \end{enumerate}
    \end{jieda}
\subsection*{直线与直线的位置关系（2）}

A 组：
\begin{enumerate}
    \item 分别在两个平面内的两条直线的位置关系是\blkda[4]{异面、相交、平行}；不平行的两条直线的位置关系是\blkda[3]{相交或异面}；没有公共点的两条直线的位置关系\blkda[3]{平行或异面}.

    \item 若两异面直线所成的角为 $ \theta $ ,则角 $ \theta $ 的范围是\blkda[3]{$(0,\dfrac{\pi}{2}]$}.
    
    \item 两条异面直线是指
    \dotfill{(\kh{D})}
    \xx{在空间不相交的两条直线}
    {分别位于两个不同平面内的两条直线}
    {某平面内的一条直线和这个平面外的一条直线}
    {不在同一平面内的两条直线} 
    
    \item 分别在两个相交平面内的两条直线的位置关系是
    \dotfill{(\kh{D})}
    \xx{异面}{平行}{相交}{异面、平行或相交}
    
    \item 分别在两个平行平面内的两条直线的位置关系是
    \dotfill{(\kh{D})}
    \xx{异面}{平行}{相交}{异面或平行}
    
    \item 分别与两条异面直线都相交的两条直线的位置关系是
    \dotfill{(\kh{B})}
    \xx{平行或相交}{相交或异面}{平行或异面}{均有可能}    

    \item 如下图,正方体 $ {ABCD} - {A}_{1}{B}_{1}{C}_{1}{D}_{1} $ 中, $M\text{、}N\text{、}P\text{、}Q\text{、}R $ 、$ S $ 分别是 $ {AB}\text{、}{BC}\text{、}{C}_{1}{D}_{1}\text{、}{C}_{1}C\text{、}{A}_{1}{B}_{1}\text{、}{B}_{1}B $ 的中点,则下列判断：(1) $ {PQ} $ 与 $ {RS} $ 共面；(2) $ {MN} $ 与 $ {RS} $ 共面；(3) $ {PQ} $与 $ {MN} $ 共面；则正确的结论是\blkda{(1)(3)}.
    
    \begin{center}

        \includegraphics[width=0.25\linewidth]{i;ugyf;iyffsa'i.png}
    \end{center} 
    
\end{enumerate}


B组：
\begin{enumerate}
    \item 已知 $ a $ 、 $ b $ 是异面直线,直线 $ c//a $ ,那么直线 $ c $ 与 $ b $ 
    \dotfill{(\kh{D})}
    \xx{一定是异面直线}{一定是相交直线}{不可能是相交直线}{不可能是平行直线}

    \item 已知 $ {a} $ 、 $ {b} $ 、 $ {c} $ 是空间三条直线,有下列四个命题,其中正确命题的个数是\blkda{1个}.
    \begin{enumerate}
        \item 若 $ a\bot b,b\bot c $ ,则 $ a\bot c $
        \item 若 $ a $ 、 $ b $ 异面, $ b $ 、 $ c $ 异面,则 $ a $ 、 $ c $ 异面
        \item 若 $ a $ 、 $ b $ 共面, $ b $ 、 $ c $ 共面,则 $ a $ 、 $ c $ 共面
        \item 若 $ a $ 、 $ b $ 平行, $ b $ 、 $ c $ 平行,则 $ a $ 、 $ c $ 平行
    \end{enumerate}

    \item 正方体 $ {{ABCD}} - {{A}}_{1}{{B}}_{1}{{C}}_{1}{{D}}_{1} $ 中, $ {{A}}_{1}{A} $ 和 $ {{D}}_{1}{{C}}_{1} $ 所成的角是\blkda{ $\dfrac{\pi}{2}$}, $ {{AC}} $ 和$ {{A}}_{1}{{B}}_{1} $ 所成的角是\blkda{ $\dfrac{\pi}{4}$}, $ {{{DA}}}_{1} $ 和 $ {{AC}} $ 所成的角是\blkda{ $\dfrac{\pi}{3}$}, $ {{AC}} $ 与 $ {{D}}_{1}{{B}}_{1} $所成的角是\blkda{ $\dfrac{\pi}{2}$}, $ {{{AD}}}_{1} $ 和 $ {{{DC}}}_{1} $ 所成的角是\blkda{$\dfrac{\pi}{3}$ }.
    
    \item 在长方体 $ {ABCD} - {A}_{1}{B}_{1}{C}_{1}{D}_{1} $ 中, $ {AB} = 5,{BC} = 4,{C{C}_{1}} = 3 $ ,则 $ {B}_{1}C $ 与 $ D{D}_{1} $ 所成的角的大小为\blkda{$arctan \dfrac{4}{3}$}, $ {BC} $ 与 $ {A}_{1}{C}_{1} $ 所成角的大小为\blkda{$arctan\dfrac{5}{4}$}, $ {B}_{1}C $ 与 $ A{D}_{1} $所成角的大小为\blkda{$arccos\dfrac{7}{25}$}

    \item 右下图是一个正方体的展开图,在原正方体中,有下列命题：
    \begin{enumerate}
        \item AB 与 CD 所在直线垂直；
        \item CD 与 EF 所在直线平行;
        \item $ {{AB}} $ 与 $ {{MN}} $ 所在直线成 $ {60}^{ \circ  } $ 角；
        \item MN 与 EF 所在直线异面
    \end{enumerate}
    其中正确命题的序号是\blkda{(3)(4)}.
    \begin{center}

        \includegraphics[width=0.25\linewidth]{imasdufha;iuge.png}
    \end{center}

    \item 在正方体 $ {ABCD} - {A}_{1}{B}_{1}{C}_{1}{D}_{1} $ 中,与 $ A{D}_{1} $ 成 $ {60}^{ \circ  } $ 角的面对角线共有\blkda{8}条.
    
    \item 已知正方体 $ {{ABCD}} - {{A}}_{1}{{B}}_{1}{{C}}_{1}{{D}}_{1} $ 中, $ {M} $ 、 $ {N} $ 分别是 $ {{{BB}}}_{1} $ 、 $ {{{CC}}}_{1} $ 的中点,则异面直线AM 与 BN 所成的角的大小是\blkda{$arccos\dfrac{1}{5}$}.

    \item 在棱长为 1 的正方体 $ {{ABCD}} - {{A}}_{1}{{B}}_{1}{{C}}_{1}{{D}}_{1} $ 中, $ {O} $ 是正方形 $ {{ABCD}} $ 的中心, $ {E} $ , $ {F} $分别是 $ {{AB}},{{BC}} $ 中点.
    \begin{enumerate}
        \item 异面直线 $ {{A}}_{1}{{D}}_{1} $ 和 $ {{CD}} $ 的距离是\blkda{1};
        
        \item $ {{A}}_{1}{{B}}_{1} $ 与 $ {{{BC}}}_{1} $ 的距离是\blkda{$\dfrac{\sqrt{2}}{2}$}；

        \item 异面直线 $ {{C}}_{1}{O} $ 和 $ {{EF}} $ 的距离是\blkda{$\dfrac{\sqrt{2}}{4}$}.
        
    \end{enumerate}

    \item 如图,在空间四边形 $ {{ABCD}} $ 中, $ {{AC}}\bot {{BD}} $ , $ {{AC}} = {{BD}} $ , 若 $ {E},{F} $ 分别是 $ {{AB}},{{CD}} $ 的中点,则 $ {{EF}} $ 与 $ {{AC}} $ 所成的角的大小为\blkda{$\dfrac{\pi}{4}$}.

    \item 若 $ {P} $ 为异面直线 $ {a} $ 、 $ {b} $ 外一点,则过 $ {P} $ 与 $ {a} $ 、 $ {b} $ 同时：
    \begin{enumerate}
        \item 平行的直线有\blkda{0}条；
        \item 垂直的直线有\blkda{1}条；
        \item 相交的直线有\blkda{0或1}条.
    \end{enumerate}
    
    \item 异面直线 $ {a} $ 、 $ {b} $ 分别在平面 $ {a} $ 、 $ {\beta } $ 内, $ {a} \cap  {\beta } = {L} $ ,则直线 $ {L} $ 与 $ {a} $ 、 $ {b} $ 的位置关系是\dotfill{(\kh{B})}
    \xx{与a、b都相交} {至少与a、b中的一条相交}{与 a、b 都不相交}{至多与 $ {a}\text{、}{b} $ 中的一条相交}

    \item 如图所示,在空间四边形 $ {{ABCD}} $ 中,E、F、G、H 分别是 $ {{AB}} $ 、 $ {{BC}} $ 、 $ {{CD}} $ 、 $ {{DA}} $的中点. 
    \begin{enumerate}
        \item 若 $ {{AC}} $ 与 $ {{BD}} $ 成 $ {60}^{ \circ  } $ 角,且 $ {{AC}} = {{BD}} = a $ ,则 $ {{EFGH}} $ 的面积为\blkda{$\dfrac{\sqrt{3}}{8}a^2$}
        \item 若 $ M $ 、 $ N $ 分别是对角线 $ {AC} $ 、 $ {BD} $ 的中点. 求证： $ {EG} $ 、 $ {FH} $ 、$ {MN} $ 交于一点且在该点互相平分.
    \end{enumerate}
    \begin{center}

        \includegraphics[width=0.7\linewidth]{f;kshdfuhsdf.png}
    \end{center}

        \begin{jieda}
            在平行四边形EHGF中，EG与FH是平行四边形的对角线，因此互相平分；在平行四边形ENGM中，EG与NM为平行四边形的对角线，因此互相平分.
        \end{jieda}

    \item 如图,在空间四边形 $ {{ABCD}} $ 中,点 $ {E} $ 、 $ {H} $ 分别是边 $ {{AB}} $ 、 $ {{AD}} $ 上的点, $ {F} $ 、 $ {G} $分别是边 $ {{BC}}\text{、}{{CD}} $ 上的点. 若 $ {{EFGH}} $ 是平行四边形,且 $ {{BD}} = a,{{AC}} = b,{{AC}} $ 与 $ {{BD}} $所成的角为 $ \theta $ ,四边形 EFGH 面积的最大值为\blkda{$\dfrac{1}{4}absin\theta$}.
    \begin{jieda}
        因为EFGH是平行四边形，因此$AE:EB=$
    \end{jieda}

    \item 如图,空间四边形 $ {ABCD} $ 中, $ {AD} = {BC} = 4,E\text{、}F $ 分别是 $ {AB}\text{、}{CD} $ 的中点, $ {EF} = 2\sqrt{3} $ .则 AD、BC 所成角为\blkda{$\dfrac{\pi}{3}$}.

    \item 如图所示,在正方体 $ {ABCD} - {A}_{1}{B}_{1}{C}_{1}{D}_{1} $ 中, $ O $ 是底面 $ {ABCD} $的中心, $ E $ 是 $ C{C}_{1} $ 的中点. 那么异面直线 $ {EO} $ 和 $ {D}_{1}A $ 所成的角的余弦值等于\blkda{$\dfrac{\sqrt{6}}{3}$}.
    
    \begin{center}

        \includegraphics[width=0.25\linewidth]{imsfuh'afueafage.png}
    \end{center}

    \item 对于已知直线 $ a $ ,假设直线 $ b $ 同时满足下列三个条件：(1)与 $ a $ 成异面直线；(2)与 $ a $ 的夹角为定值 $ \theta $ ；(3)与 $ a $ 的距离为定值 $ d $ .那么,这样的直线 $ b $ 有\blkda{无数}条.

    
    \item 在正方体的 12 条棱、12 条面对角线中,总共可以组成\blkda{126}对异面直线.
    \begin{jieda}
        棱与棱可以形成$\dfrac{12*4}{2}=24$对异面直线，线与线可以形成$\dfrac{12*5}{2}=30$对异面直线，棱与线可以形成$6*12=72$对异面直线，
    \end{jieda}
    
    \item 若 $ a,b $ 为异面直线, $ m \cap  a = A,m \cap  b = B,n \cap  a = M,n \cap  b = N $ ,求证： $ m,n $是异面直线.
    \begin{jieda}
        因为直线$a$与直线$b$异面，所以$A\notin b$，从而A与b可确定平面$\alpha$.因为$B\in b, b\subset\alpha$，所以$B\in \alpha$，又因为$A\in \alpha$，所以$AB \subset \alpha$.因为直线$a$与直线$b$异面，所以$M\notin \alpha$.因为$N\in b\subset \alpha$,且$N\notin m$，由异面直线的判定定理，n与m异面.
        \begin{center}
            \includegraphics[width=0.25\linewidth]{a'huf'auhfdf.png}
        \end{center}
    \end{jieda}
    
    \item 如图,在正方体 $ {{ABCD}} - {{A}}_{1}{{B}}_{1}{{C}}_{1}{{D}}_{1} $ 中, $ {E} $ 是 $ {{B}}_{1}{{C}}_{1} $ 的中点, $ {O} $ 是正方形 $ {{A}}_{1}{{B}}_{1}{{C}}_{1}{{D}}_{1}{R} $的中心,连接 $ {{AO}},{{CE}} $ ,求异面直线 $ {{AO}} $ 与 $ {{CE}} $ 所成的角的大小.

    \begin{center}

        \includegraphics[width=0.25\linewidth]{fhsuhfajsdf.png}
    \end{center}

    \begin{jieda}
        取$CD$的中点M，可证得$\angle{AOM}$即为异面直线所成角，$|AO|=\dfrac{\sqrt{6}}{2},|MO|=\dfrac{\sqrt{5}}{2},|AM|=\dfrac{\sqrt{5}}{2}$，由余弦定理可得$cos\angle{AOM=\dfrac{\sqrt{30}}{10}}$，综上异面直线AO与CE所成角为$arccos \dfrac{\sqrt{30}}{10}$.
        \begin{center}
    
            \includegraphics[width=0.25\linewidth]{asoufha'oiusf.png}
        \end{center}
    \end{jieda}

    
    \item 在正方体 $ {ABCD} - {A}_{1}{B}_{1}{C}_{1}{D}_{1} $ 中,正方形 $ {AB}{B}_{1}{A}_{1} $ 和正方形 $ {BC}{C}_{1}{B}_{1} $ 的中心分别为 $ O $和 $ {O}_{1} $ ,试判断直线 $ O{O}_{1} $ 与直线 $ B{D}_{1} $ 的位置关系,并证明你的结论.
    \begin{jieda}
        异面.记平面$A_1BC_1$为$\alpha$，因为$OO_1\subset\alpha, B\in \alpha 且 B\notin OO_1，D\notin \alpha$，由异面直线判定定理可知，两直线异面.
        \begin{center}
    
            \includegraphics[width=0.25\linewidth]{imauohf'aouhfage.png}
        \end{center}
    
    \end{jieda}

    \item 已知正方体 $ {ABCD} - {A}_{1}{B}_{1}{C}_{1}{D}_{1} $ 的棱长为 $ a,M\text{、}N $ 分别为 $ {A}_{1}{B}_{1}\text{、}{B}_{1}B $ 的中点,求 $ {AM} $ 与 $ {CN} $ 所成的角.

    \begin{center}

        \includegraphics[width=0.25\linewidth]{imfsdhfffage.png}
    \end{center}

    \begin{jieda}
        取$D_1C_1$的中点$M_1$，$AA_1$的中点$N_1$，可证明$\angle{N_1DM_1}$即为异面直线的夹角，利用余弦定理可得$arccos \dfrac{2}{5}$
        \begin{center}
    
            \includegraphics[width=0.25\linewidth]{ahfa'ouhfaaf.png}
        \end{center}
    \end{jieda}
\end{enumerate}



\section*{直线与平面的位置关系}
\subsection*{直线与平面的位置关系（1）}

知识点：线面平行

A组：
\begin{enumerate}
    \item 判断下列命题是否正确:
        \begin{enumerate}
            \item 过平面外一点可作无数条直线与这个平面平行. \blkda{\ding{51}}
            
            \item 若直线 $ l \not \subset \alpha $ ,则 $ l $ 不可能与 $ \alpha $ 内无数条直线相交. \blkda{\ding{55}}
            
            \item 若直线 $ l $ 与平面 $ \alpha $ 不平行,则 $ l $ 与 $ \alpha $ 内任一直线都不平行.\blkda{\ding{55}}
            
            \item 经过两条平行线中一条直线的平面平行于另一条直线.\blkda{\ding{55}}
            
            \item 若平面 $ \alpha $ 内有一条直线和直线 $ l $ 异面,则 $ l \not \subset \alpha$ . \blkda{\ding{51}}
        \end{enumerate}
 
    \item 已知平面 $ \alpha $ 与直线 $ l $ ,则 “ $ l//\alpha $ ” 是 “直线 $ l $ 与平面 $ \alpha $ 无公共点” 的\dotfill{(\kh{C})}
    \xx{充分不必要条件} {必要不充分条件}{充要条件}{既不充分也不必要条件}

    \item 下列命题中,正确的是\dotfill{(\kh{D})}
    \xx{若直线 $ a $ 与平面 $ \alpha $ 平行,则 $ a $ 平行于 $ \alpha $ 内的任何直线}
    {若两直线 $ a,b $ 都与平面 $ {\alpha } $ 平行,则 $ {a}//{b} $}
    {若直线 $ a $ 平行于平面 $ \alpha $ ,直线 $ b $ 在平面 $ \alpha $ 内,则 $ a//b $}
    {若直线 $ l $ 与平面 $ \alpha $ 平行,则平面 $ \alpha $ 内有无数条直线与 $ l $ 平行}
    
    \item 如图,在下列四个正方体中, $ {A},{B} $ 为正方体的两个顶点, $ {M},{N},{Q} $ 为所在棱的中点, 则在这四个正方体中, 直线 $ {{AB}} $ 与平面 $ {{MNQ}} $ 平行的有\blkda{3}个.
    \begin{jieda}
        (1)中令$NM$的中点为T,则$AB$与$QT$相交（因为共面且不平行）;(2)中$AM // MQ$；(3)中$AB // MQ$;（4）中$AB//NQ$
    \end{jieda}

    \begin{center}

        \includegraphics[width=0.75\linewidth]{imaarfharhfafge.png}
    \end{center}

\end{enumerate}

B组：
\begin{enumerate}
    \item 在正方体 $ {{ABCD}} - {{A}}_{1}{{B}}_{1}{{C}}_{1}{{D}}_{1} $ 中,和平面 $ {{A}}_{1}{{DB}} $ 平行的侧面对角线有\blkda{3}条；分别是\blkda[4]{$CD_1,CB_1,D_1B_1$}.
    
    \item 在正方体 $ {{ABCD}} - {{A}}_{1}{{B}}_{1}{{C}}_{1}{{D}}_{1} $ 中, $ {E} $ 为 $ {{{DD}}}_{1} $ 的中点,则 $ {{{BD}}}_{1} $ 与平面 $ {{AEC}} $ 的位置关系为\blkda{平行}.
    
    \item 过直线外一点和这条直线平行的平面有\blkda{无数}个.

    \item 在空间四边形 $ {ABCD} $ 中, $ E $ 、 $ F $ 分别为 $ {AB} $ 、 $ {AD} $ 上的点,且 $ {{AE}} : {{EB}} = {{AF}} : {{FD}} = 1 : 4 $ ,又 $ H\text{、}G $ 分别为 $ {BC}\text{、}{CD} $ 的中点,则\blkda{（2）}(填序号).
    \begin{enumerate}
        \item  $ {{BD}}// $ 平面 $ {EFG} $ ,且四边形 $ {EFGH} $ 是平行四边形

        \item $ {EF}// $ 平面 $ {BCD} $ ,且四边形 $ {EFGH} $ 是梯形

        \item $ {{HG}}// $ 平面 $ {ABD} $ ,且四边形 $ {EFGH} $ 是平行四边形

        \item $ {{EH}}// $ 平面 $ {ADC} $ ,且四边形 $ {EFGH} $ 是梯形
        
    \end{enumerate}

    \begin{center}

        \includegraphics[width=0.25\linewidth]{fsajhsdfa'hsud.png}
    \end{center}

    
    \item 若直线 $ a//b $ , $ a// $ 平面 $ \alpha $ ,则 $ b $ 与平面 $ \alpha $ 的位置关系是\blkda[4]{$b//\alpha \text{或} b\subset \alpha$}.

    \item 已知直线 $ a// $ 平面 $ a $ ,点 $ A \in  a $ ,点 $ A \in $ 直线 $ b $ ,且 $ a//b $ ,则直线 $ b $ 与平面 $ a $ 的位置关系是\blkda[2]{$b\subset \alpha$}.

    
    \item 设 $ m,n $ 是平面 $ a $ 外的两条直线,给出三个论断： $ \left( 1\right) m//n,\left( 2\right) m//a,\left( 3\right) n//a $ ,以其中两个为条件, 余下的一个为结论, 写出你认为正确的一个命题：\blkda[6]{(1)(2)推(3);(1)(3)推(2)}.

    
    \item 已知 $ \bigtriangleup  {ABC} $ 中, $ {AB} = 5,{AC} = 7,A = {60}^{ \circ  } $ , $ G $ 是重心,若过 $ G $ 的平面 $ a $ 与 $ {BC} $ 平行, $ {AB} \cap  a = M,{AC} \cap  a = N $ ,则 $ {MN} = $ \blkda{$\dfrac{2\sqrt{39}}{3}$}.

    
    \item 直线与平面平行的充要条件是 \dotfill{(\kh{C})}
    \xx{直线与平面内的一条直线平行} {直线与平面内的两条直线不相交}
    {直线与平面内的任一直线都不相交} {直线与平行内的无数条直线平行}

    \item 下列命题正确的是 \dotfill{(\kh{C})}
    \xx{一直线与平面平行,则它与平面内任一直线平行}{一直线与平面平行,则平面内有且只有一个直线与已知直线平行}{ 一直线与平面平行,则平面内有无数直线与已知直线平行,它们在平面内彼此平行}{一直线与平面平行,则平面内任意直线都与已知直线异面}


    \item 若直线 $ a $ 在平面 $ \alpha $ 内,直线 $ a $ 、 $ b $ 是异面直线,则直线 $ b $ 和 $ \alpha $ 平面的位置关系是 \dotfill{(\kh{C})}
    \xx{相交} {平行} {相交或平行} {相交且垂直}

    \item 下列各命题中假命题的个数为\dotfill{(\kh{B})}
    \begin{enumerate}
        \item 经过两条平行直线中一条直线的平面必平行于另一条直线；
        \item 若一条直线平行于两相交平面,则这条直线和交线平行；
        \item 空间四边形中三条边的中点所确定平面和这个空间四边形的两条对角线都
平行.
    \end{enumerate}
    \xx{0}{1}{2}{3}

    \item 已知 $ a $ 、 $ b $ 为两异面直线,下列结论正确的是\dotfill{(\kh{D})}
    \xx{过不在 $ a\text{、}b $ 上的任何一点,可作一个平面与 $ a\text{、}b $ 都平行}
    {过不在 $ a\text{、}b $ 上的任一点,可作一直线与 $ a\text{、}b $ 都相交}
    {过不在 $ a\text{、}b $ 上任一点,可作一直线与 $ a\text{、}b $ 都平行}
    {过 $ a $ 可以并且只可以作一个平面与 $ b $ 平行}

    
    \item 如图,已知 $ {{ABCD}} $ 与 $ {{BAFE}} $ 是两个不共面的正方形,且 $ {{AM}} = {{FN}} $ .求证: $ {{MN}}// $ 平面 $ {{BEC}} $ .
    
    \begin{center}
        \includegraphics[width=0.25\linewidth]{asdfhasldf.png}
    \end{center}
    \begin{jieda}

     \begin{center}
        \includegraphics[width=0.25\linewidth]{bodyxb//src/fa;skjdhfa'hfa.png}
    \end{center}   
    \end{jieda}
    
    \begin{jieda}
        过M作$MG//AB$，交$BC$于点G，过N作$NK//AB$，交$BE$于点G，可以证明$GK\subset \alpha,MN\not\subset,GK//MN$,故$MN//平面BEC$.
    \end{jieda}

    
    \item 用平行于空间四边形 $ {{ABCD}} $ 的一组对边 $ {{AB}}\text{、}{{CD}} $ 的平面截此空间四边形。
    \begin{enumerate}
        \item 求证：所得截面 $ {{MNPQ}} $ 是平行四边形；
        \item 如果 $ {{AB}} = {{CD}} = a $ ,求证：四边形 $ {{MNPQ}} $ 的周长为定值.
    \end{enumerate}
  
    \begin{center}

        \includegraphics[width=0.25\linewidth]{cjshf'kashdfws.png}
    \end{center}
    
    \begin{jieda}
    \begin{enumerate}
        \item AB//平面MNPQ，$AB\subset 平面ABC, 平面ABC\cap 平面MNPQ=MN$，由线面平行的判定定理可知，$AB//MN$，同理可得$AB//PQ, CD//MQ, CD//NP$，因此$PQ//MN,NP//MQ$，因此为平行四边形.
        \item 因为$AB//MN,MQ//CD$，因此$\dfrac{BN}{NC}=\dfrac{AM}{MC}=\dfrac{AQ}{QD}=\dfrac{1-t}{t}$，从而周长$=2(ta+(1-t)a)=2a$.
    \end{enumerate}

    \end{jieda}

    
\end{enumerate}


    





\subsection*{直线与平面的位置关系（2）}

知识点：线面垂直

A 组：

\begin{enumerate}

\item 下列关于直线与平面垂直的判断中，正确的是\dotfill{(\kh{C})}

\xx{若直线 ${a}$ 与平面 ${\alpha }$ 内的一条直线垂直，则直线 ${a}$ 与平面 ${\alpha }$ 垂直}
{若直线 ${a}$ 与平面 ${\alpha }$ 内的两条平行直线垂直，则直线 ${a}$ 与平面 ${\alpha }$ 垂直}
{若直线 ${a}$ 与平面 ${\alpha }$ 内的两条相交直线垂直，则直线 ${a}$ 与平面 ${\alpha }$ 垂直}
{若直线 ${a}$ 与平面 ${\alpha }$ 内的无数条直线垂直，则直线 ${a}$ 与平面 ${\alpha }$ 垂直}

\item 给出下列命题：（1）若直线 ${l}$ 与平面 ${\alpha }$ 中的无数条直线垂直，则 ${l}\bot {\alpha }$ ；（2）若直线 ${l}\bot$ 平面 ${\alpha }$ ,且直线 ${m} \bot$ 平面 ${\alpha }$ ,则 ${l}//{m}$ ; (3) 若 ${A} \in  {l},{B} \in  {l}$ 且 ${A} \in  {\alpha },{B} \in  {\alpha }$ ,可得 ${l} \subset  {\alpha }$ . 其中真命题的个数是\dotfill{(\kh{C})}
\xx{0}{1}{2}{3}

\item m，n为空间两条不重合直线，α为空间平面，下列命题正确的是\dotfill{(\kh{B})}

\xx{${m} \bot  {\alpha },{n} \bot  {m}$ ,则 ${n}// {\alpha }$}
{${m},{n}$ 与 ${\alpha }$ 所成角均为 ${90}^{ \circ  }$ ,则 ${m}// {n}$}
{${m}// {\alpha },{n}// {\alpha },{m}// {n}$ ,则直线 ${m},{n}$ 到 ${\alpha }$ 的距离相等}
{${m}// {\alpha },{n}// {\alpha }$ ,则 ${m}// {n}$}

\item 下列四个正方体图形中, $l$ 是正方体的一条对角线,点 ${D}\text{ 、 }{E}\text{ 、 }{F}$ 分别为其所在棱的中点, 能得出 ${I} \bot$ 平面 ${DEF}$ 的是 \dotfill{(\kh{C})}

\begin{center}
    \centering
    \includegraphics[width=0.9\linewidth]{bodyxb//src/aouihfiaqgefiq'ihwerfq.png}
\end{center}

\item 如图，已知正方体 ${ABCD} - {A}_{1}{B}_{1}{C}_{1}{D}_{1}$ 的棱长为 1 .

\begin{center}
    \centering
    \includegraphics[width=0.25\linewidth]{bodyxb//src/akuihfai'bf'ihqwiherfvq.png}
\end{center}
\begin{enumerate}
    \item 点 $A$ 到平面 $B{B}_{1}{C}_{1}C$ 的距离为\blkda{1}；
    \item 直线 ${B}_{1}{D}_{1}$ 和平面 ${ABCD}$ 的距离为\blkda{1}；
    \item 直线 ${A}_{1}{B}_{1}$ 和平面 ${AB}{C}_{1}{D}_{1}$ 的距离为\blkda{$\dfrac{\sqrt{2}}{2}$}.
\end{enumerate}

\end{enumerate}

B 组:
\begin{enumerate}

\item 下列命题中正确的是 \dotfill{(\kh{D})}

\xx{若直线 $a \bot  a$ 内的一条直线,则 $a \bot  a$}
{若直线 $a \bot  a$ 内的无数条直线,则 $a \bot  a$}
{若 $a$ 与 $a$ 不垂直,则 $a$ 内不存在与 $a$ 垂直的直线}
{若 $a$ 与 $a$ 不垂直，则 $a$ 内不存在与 $a$ 垂直的两条相交直线}

\item 下列命题中，假命题是 \dotfill{(\kh{D})}

\xx{如果一条直线与一个平面垂直，则这条直线与这个平面内的任意一条直线垂直}
{如果两条直线同垂直于一个平面，则这两条直线平行}
{平面的垂线与这个平面一定相交}
{如果一条直线上有两点到一个平面的距离相等，则这条直线与这个平面平行}

\item 下列四个命题，其中正确命题的个数是 \dotfill{(\kh{C})}
\begin{enumerate}
    \item 点 $A \notin$ 平面 $\alpha$ ，直线 $a \subset  a$ ，则 $A$ 到 $a$ 的距离等于 $A$ 到 $\alpha$ 的距离；
    
    \item 直线 $a//$ 平面 $\alpha$ ， $b \subset  a$ ， $a// b$ ，则 $a$ 、 $b$ 间的距离等于 $a$ 与 $\alpha$ 间的距离；
    
    \item $A \in$ 直线 $a,a.m$ 直线 $b$ ，则 $a\text{ 、 }b$ 间的距离等于 $A$ 到 $b$ 的距离；
        
    \item 直线 $a//$ 平面 $\alpha ,A \in  a$ ，则 $a\text{ 、 }\alpha$ 间的距离等于 $A$ 到 $\alpha$ 的距离。
\end{enumerate}

\xx{4 个} {3 个} {2 个} {1 个}

\item 从平面 $a$ 外一点 P 作与 $a$ 相交的直线，使 P 与交点的距离为 1，则满足条件的直线不可能是\dotfill{(\kh{C})}

\xx{0 条} {1 条} {2 条} {无数条}

\item 若正 $\bigtriangleup {ABC}$ 的边长为 $\dfrac{4\sqrt{3}}{3}$ ，则到三个顶点的距离都为 1 的平面有 \dotfill{(\kh{C})}

\xx{1 个} {3 个} {5 个} {7 个}

\item “直线 $l$ 垂直于平面 $\alpha$ 内的无数条直线”是“ $l\bot a$ ”的\blkda{必要非充分}条件.

\item 已知 $m,n$ 是空间两条相交直线， ${l}_{1},{l}_{2}$ 是与 $m,n$ 都垂直的两条直线，直线 $l$ 与 ${l}_{1},{l}_{2}$ 都相交，则直线 $l$ 与 ${l}_{1},{l}_{2}$ 所成的角的大小关系为\blkda{相等}.

\item 已知三条相交于一点的线段 ${PA}\text{ 、 }{PB}\text{ 、 }{PC}$ 两两垂直，且 $A\text{ 、 }B\text{ 、 }C$ 在同一平面内， $P$ 在平面 ${ABC}$ 外， ${PH}\bot$ 平面 ${ABC}$ 于 $H$ ，则垂足 $H$ 是 $\bigtriangleup  {ABC}$ 的\blkda{垂}心


\item 已知BC是Rt $\bigtriangleup  {ABC}$ 的斜边， ${AP}\bot$ 平面 ${ABC}$ ，连结 ${PB}\text{ 、 }{PC}$ ，作 ${PD}\bot {BC}$ 于 $D$ ，连结 ${AD}$ ， 则图中共有直角三角形\blkda{8}个， $\bigtriangleup  {PBC}$ 是\blkda{锐角}三角形.

\item 在棱长为 $\alpha$ 的正方体 ${ABCD} - {A}_{1}{B}_{1}{C}_{1}{D}_{1}$ 中， $A$ 到直线 ${B}_{1}C$ 的距离为\blkda{$\dfrac{\sqrt{6}a}{2}$}， $A{A}_{1}$ 到平面 $B{B}_{1}{D}_{1}D$ 的距离为\blkda{$\dfrac{\sqrt{2}a}{2}$}， $A$ 到平面 ${A}_{1}{BD}$ 的距离为\blkda{$\dfrac{\sqrt{3}a}{3}$}.


\item 已知线段 ${AB}$ 的两端点 $A$ 、 $B$ 到平面 $\alpha$ 的距离分别为 $a$ 和 $b$ ， （1）若 $A$ 、 $B$ 在 $\alpha$ 的同侧，则 ${AB}$ 中点到 $\alpha$ 的距离为\blkda{$\dfrac{a+b}{2}$}.

（2）若 $A$ 、 $B$ 在 $\alpha$ 的异侧，则 ${AB}$ 中点到 $\alpha$ 的距离为\blkda{$\dfrac{|a-b|}{2}$}.

\item 若三角形 ${ABC}$ 三个顶点到平面 $\alpha$ 的距离分别为 $a$ 、 $b$ 、 $c$ ， $\bigtriangleup {ABC}$ 的重心为 $G$ ，则点 $G$ 到平面 $\alpha$ 的距离为\blkda[10]{$\dfrac{a+b+c}{3}$,$\dfrac{|-a+b+c|}{3}$,$\dfrac{|a-b+c|}{3}$,$\dfrac{|a+b-c|}{3}$}.

\item 如图 ABCD 是矩形， ${PA} \bot$ 平面 ${ABCD}$ ， $\bigtriangleup  {PAD}$ 是等腰三角形，M，N 分别是 AB，PC 的中点,求证: ${MN} \bot$ 平面 ${PCD}$ .

\begin{center}
    \centering
    \includegraphics[width=0.25\linewidth]{bodyxb//src/aufh'fibqirvbf'qiyefq.png}
\end{center}

\begin{jieda}
    取 $ {PD} $ 的中点 $ E $ ,连接
$ {AE},{NE} $ ,
$ \because N $ 为 $ {PC} $ 的中点,
$ \therefore {NE} $ 为 $ \bigtriangleup  {PCD} $ 的中位线,
$ \therefore {NE}\overset{//}{ = }\dfrac{1}{2}{CD} $ .
在矩形 $ {ABCD} $ 中, $ {AB} ⫫ {CD} $ ,
又 $ M $ 为 $ {AB} $ 的中点, $ \therefore {AM}\overset{//}{ = }\dfrac{1}{2}{CD} $ .
$ \therefore {AM}\overset{//}{ = }{NE} $ ,
$ \therefore $ 四边形 $ {AMNE} $ 为平行四边形, $ \therefore {AE}//{MN} $ .
$ \because {PA} \bot $ 平面 $ {ABCD},{AD} \subset $ 平面 $ {ABCD} $ ,
$ \therefore {PA}\bot {AD} $ .
又 $ \bigtriangleup {PAD} $ 为等腰三角形, $ E $ 为 $ {PD} $ 的中点,
$ \therefore {PA} = {AD},{AE} \bot  {PD},\therefore {MN} \bot  {PD} $ .
$ \because {PA}\bot $ 平面 $ {ABCD},{CD} \subset $ 平面 $ {ABCD} $ ,
$ \therefore {PA}\bot {CD} $ .
在矩形 $ {ABCD} $ 中, $ {CD} \bot  {AD} $ ,
又 $ {PA} \cap  {AD} = A,{PA},{AD} \subset $ 平面 $ {PAD} $ ,
$ \therefore {CD} \bot $ 平面 $ {PAD} $ ,
$ \because {AE} \subset $ 平面 $ {PAD},\therefore {AE} \bot  {CD},\therefore {MN} \bot  {CD} $ ,
又 $ {PD} \cap  {CD} = D,{PD},{CD} \subset $ 平面 $ {PCD} $ ,
$ \therefore {MN}\bot $ 平面 $ {PCD} $ .
\end{jieda}

\item 如图，已知 ${ABCD}$ 是矩形， ${AB} = a,{AD} = b$ ， ${PA} \bot$ 平面 ${ABCD}$ ， ${PA} = {2c}$ ， $Q$ 是 ${PA}$ 的中点. 求
\begin{enumerate}
    \item $Q$ 到 ${BD}$ 的距离；
    \item $P$ 到平面 ${BQD}$ 的距离.
\end{enumerate}

\begin{center}
    \centering
    \includegraphics[width=0.25\linewidth]{bodyxb//src/afuhg'igq'iwyf'yiq.png}
\end{center}

\begin{jieda}
    \begin{enumerate}
        \item 在矩形 $ {ABCD} $ 中,作 $ {AE}\bot {BD},E $ 为垂足,连
结 $ {QE} $ .
$ \because {QA} \bot $ 平面 $ {ABCD} $ ,由三垂线定理得 $ {QE} \bot  {BE} $ ,
$ \therefore {QE} $ 的长为 $ Q $ 到 $ {BD} $ 的距离.
在矩形 $ {ABCD} $ 中, $ {AB} = a,{AD} = b $ ,
$ \therefore {AE} = \dfrac{ab}{\sqrt{{a}^{2} + {b}^{2}}} $ , 图 3-2-101
在 Rt $ \bigtriangleup {QAE} $ 中, $ {QA} = \dfrac{1}{2}{PA} = c $ ,
$ \therefore {QE} = \sqrt{Q{A}^{2} + A{E}^{2}} = \sqrt{{c}^{2} + \dfrac{{a}^{2}{b}^{2}}{{a}^{2} + {b}^{2}}}\therefore Q $ 到 $ {BD} $ 的距离为 $ \sqrt{{c}^{2} + \dfrac{{a}^{2}{b}^{2}}{{a}^{2} + {b}^{2}}} $ .
    \item 解法一 因为平面 $ {BQD} $ 经过线段 $ {PA} $ 的中点,
$ \therefore P $ 到平面 $ {BQD} $ 的距离等于 $ A $ 到平面 $ {BQD} $ 的距离.
在 $ \bigtriangleup {AQE} $ 中, $ {AH}\bot {QE},H $ 为垂足.
$ \because {BD} \bot  {AE},{BD} \bot  {QE},\therefore {BD} \bot $ 平面 $ {AQE}.\therefore {BD} \bot  {AH} $ .
$ \therefore {AH} \bot $ 平面 $ {BQE} $ ,即 $ {AH} $ 为 $ A $ 到平面 $ {BQD} $ 的距离.
在 Rt $ \bigtriangleup {AQE} $ 中, $ \because {AQ} = c,{AE} = \dfrac{ab}{\sqrt{{a}^{2} + {b}^{2}}} $ ,
$ \therefore {AH} = \dfrac{abc}{\sqrt{\left( {{a}^{2} + {b}^{2}}\right) {c}^{2} + {a}^{2}{b}^{2}}} $ .
$ \therefore P $ 到平面 $ {BD} $ 的距离为 $ \dfrac{abc}{\sqrt{\left( {{a}^{2} + {b}^{2}}\right) {c}^{2} + {a}^{2}{b}^{2}}} $ .\\
解法二 本题也可用体积法求解
设点 $ A $ 到平面 $ {QBD} $ 的距离为 $ h $ ,
由 $ {V}_{A - {BQD}} = {V}_{Q - {ABD}} $ 得
$ \dfrac{1}{3}{S}_{\bigtriangleup {BQD}} \cdot  h = \dfrac{1}{3}{S}_{\bigtriangleup {ABD}} \cdot  {AQ} $
$ h = \dfrac{{S}_{\bigtriangleup {ABD}} \cdot  {AQ}}{{S}_{\bigtriangleup {BQD}}} = \cdots  = \dfrac{abc}{\sqrt{\left( {{a}^{2} + {b}^{2}}\right) {c}^{2} + {a}^{2}{b}^{2}}} $
    \end{enumerate}
\end{jieda}



\end{enumerate}

C 组:
\begin{enumerate}

\item 如图，在边长为 1 的正方体 ${{ABCD}} - {{A}}_{1}{{B}}_{1}{{C}}_{1}{{D}}_{1}$ 中， ${E},{F},{M}$ ， ${N}$ 分别为棱 ${{AB}}$ ， ${AD},{D}_{1}{C}_{1},{C}_{1}{B}_{1}$ 的点,满足 $\dfrac{{C}_{1}M}{{C}_{1}{D}_{1}} = \dfrac{{C}_{1}N}{{C}_{1}{B}_{1}} = \dfrac{AE}{AB} = \dfrac{AF}{AD} = \lambda \left( {\lambda  \in  \left( {0,1}\right) }\right)$ ,过 $E,F$ , ${M},{N}$ 四点作该正方体的截面，则下列说法错误的是 \dotfill{(\kh{BD})}

\begin{center}
    \centering
    \includegraphics[width=0.25\linewidth]{bodyxb//src/oafhqouh3rofhow3uf.png}
\end{center}

\xx{$\lambda  = \dfrac{7}{2}$ 时，该截面是正六边形}
{$\lambda  = \dfrac{7}{8}$ 时，四边形 ${E}{F}{M}{N}$ 为正方形}
{${MN}//$ 平面 $A{D}_{1}{B}_{1}$}
{当四边形${E}{F}{M}{N}$ 为正方形时,它的面积为 $\dfrac{8}{9}$}

\item 如图，在平行四边形 ${{ABCD}}$ 中， ${{BC}} = 2{{AB}} = {4}$ ， $\angle {{ABC}} = {{{60}}}^{ \circ  }$ ， ${{PA}} \bot$ 平面 ${{ABCD}},{{PA}} = 2,{E},{F}$ 分别为 ${{BC}},{{PE}}$ 的中点,求证: ${{AF}} \bot$ 平面 ${{PED}}$ .

\begin{center}
    \centering
    \includegraphics[width=0.25\linewidth]{bodyxb//src/gsh'eouthgouergsbje.png}
\end{center}

\begin{jieda}
    证明：如图,连接 $ {AE} $ ,在平行四边形 $ {ABCD} $ 中,
$ {BC} = {2AB} = 4 $ ,
$ \angle {ABC} = {60}^{ \circ  },E $ 是 $ {BC} $ 的中点,
所以 $ {AE} = 2,{ED} = 2\sqrt{3} $ ,
从而有 $ A{E}^{2} + E{D}^{2} = A{D}^{2} $ , 所以 $ {AE}\bot {ED} $ ,
又 $ {PA} \bot $ 平面 $ {ABCD},{ED} \subset $ 平面 $ B $ $ C $
$ {ABCD} $ ,
所以 $ {PA}\bot {ED} $ ,
因为 $ {PA} \cap  {AE} = A,{PA},{AE} \subset $ 平面 $ {PAE} $ ,
所以 $ {ED}\bot $ 平面 $ {PAE} $ ,
又 $ {AF} \subset $ 平面 $ {PAE} $ ,所以 $ {ED} \bot  {AF} $ .
又 $ {PA} = {AE} = 2,F $ 为 $ {PE} $ 的中点,
所以 $ {AF}\bot {PE} $ ,
又 $ {PE} \cap  {ED} = E,{PE},{ED} \subset $ 平面 $ {PED} $ ,
所以 $ {AF} \bot $ 平面 $ {PED} $ .


\end{jieda}

\item 如图， ${{PA}} \bot$ 矩形 ${ABCD}$ ，有下列结论：① ${{PB}} \bot  {{BC}}$ ， ${{QPD}} \bot  {{CD}}$ ，③ ${{PD}} \bot  {BD}$ ， ④ ${{PA}} \bot  {{BD}}$ . 其中正确的是\blkda{(1)(2)(4)}(填序号).
\begin{center}
    \centering
    \includegraphics[width=0.25\linewidth]{bodyxb//src/lajerfoawbfojaer.png}
\end{center}


\end{enumerate}



\clearpage

\subsection*{直线与平面的位置关系（3）}

知识点：线面所成角

A组：
\begin{enumerate}
\item 直线 ${l}$ 是平面 ${\alpha }$ 的一条斜线， ${l}$ 与平面 ${\alpha }$ 内的直线所成角的取值范围是\blkda{$[<l,\alpha>,\dfrac{\pi}{2}]$}.

\item 已知空间中两条直线 ${a},{b}$ 无公共点，则 “直线 ${a},{b}$ 与平面 ${\alpha }$ 所成的角相等” 是 “ ${a}//{b}$ ” 的\dotfill{(\kh{B})}
\xx{充分不必要条件}{必要不充分条件}{充要条件}{既不充分也不必要条件}

\item 若直线 ${l}$ 与平面 ${\alpha }$ 所成的角为 ${\theta }$ ，直线 ${l}$ 与平面 ${\alpha }$ 上的直线 ${m}$ 所成的角为 ${\varphi }$ ，则总有 \dotfill{(\kh{B})}
\xx{${\theta } < {\varphi }$ }{${\theta } \leq  {\varphi }$}{C. ${\theta } > {\varphi }$ }{D. ${\theta } \geq  {\varphi }$}

\item 直线 $l$ 与平面 ${\alpha }$ 相交， $A,B$ 在 $l$ 上， ${{AB}} = {4},{{AB}}$ 在 ${\alpha }$ 内的射影长为 $2\sqrt{2}$ ， ${l}$ 与 ${\alpha }$ 夹角为\blkda{$\dfrac{\pi}{4}$}.

\item 在正方体 ${{ABCD}} - {{A}}_{1}{{B}}_{1}{{C}}_{1}{{D}}_{1}$ 中， $E$ 是 ${A}{{A}}_{1}$ 的中点，则直线 ${{B}}_{1}{E}$ 与平面 BCE所成角的正弦值为\blkda{$\dfrac{4}{5}$}.
    
\end{enumerate}


B 组:
\begin{enumerate}


\item 下列命题：\ding{172}一条直线在平面内的射影是一条直线；\ding{173} ${Rt} \bigtriangleup  {ABC}$ 的斜边在平面内，直角顶点 $C$ 在平面内的射影为 ${C}^{\prime }$ ，则 $\bigtriangleup  {AB}{C}^{\prime }$ 是钝角三角形；\ding{174}在同一平面内的射影长相等，则斜线长相等；\ding{175}两斜线与平面所成的角相等，则这两斜线互相平行.其中真命题的个数是\blkda{0}.

\item 下列命题中正确的是\dotfill{(\kh{A})}

\xx{若平面 $M$ 外的两条直线在平面 $M$ 内的射影为一条直线及此直线外的一个点, 则它们互为异面直线}
{若平面 $M$ 外的两条直线在平面 $M$ 内的射影为两条相交直线，则这两条直线相交}
{若平面 $M$ 外的两条直线在平面 $M$ 内的射影为两条平行直线，则这两条直线平行}
{若平面 $M$ 外的两条直线在平面 $M$ 内的射影为两条互相垂直的直线，则这两条直线垂直}

\item 已知直线 $a$ 与平面 $\alpha$ 所成的角为 ${30}^{ \circ  },b \subset  a$ 且与 $a$ 异面，若直线 $a$ 与 $b$ 所成的角为 $\varphi$ ,则\dotfill{(\kh{C})}

\xx{${0}^{ \circ  } < \varphi  \leq  {30}^{ \circ  }$ }{ ${0}^{ \circ  } < \varphi  \leq  {90}^{ \circ  }$}{  ${30}^{ \circ  } \leq  \varphi  \leq  {90}^{ \circ  }$}{ ${30}^{ \circ  } \leq  \varphi  \leq  {180}^{ \circ  }$}


\item 已知 ${AB}\bot$ 平面 $a$ 于 $B$ ， ${BC}$ 为 ${AC}$ 在 $a$ 内的射影， ${CD} \subset  a$ ，若 $\angle {ACD} = {60}^{ \circ  }$ ， $\angle {BCD} = {45}^{ \circ  }$ ，则 ${AC}$ 与平面 $a$ 所成的角为\blkda{$\dfrac{\pi}{4}$}.

\item 在正方体 ${ABCD} - {A}_{1}{B}_{1}{C}_{1}{D}_{1}$ 中，直线 $A{B}_{1}$ 与平面 ${AB}{C}_{1}{D}_{1}$ 所成的角的大小为\blkda{$\dfrac{\pi}{6}$}.

\item 已知 $P$ 是 $\bigtriangleup  {ABC}$ 所在平面外一点， $O$ 是 $P$ 点在平面 $\alpha$ 上的射影。若 $P$ 到 $\bigtriangleup  {ABC}$ 三边的距离相等，则 $O$ 是 $\bigtriangleup  {ABC}$ 的\blkda{内心}，若 $P$ 到 $\bigtriangleup  {ABC}$ 三个顶点的距离相等，则 $O$ 是 $\bigtriangleup  {ABC}$ 的\blkda{外心}，若 ${PA}\text{ 、 }{PB}\text{ 、 }{PC}$ 两两互相垂直，则 $O$ 是 $\bigtriangleup  {ABC}$ 的\blkda{垂心}，

\item 已知 $\bigtriangleup  {ABC}$ 中， $A \in  a$ ， ${BC}//a$ ， ${BC} = 6$ ， $\angle {BAC} = {90}^{ \circ  }$ ， ${AB}$ 、 ${AC}$ 与平面 $a$ 分别成 ${30}^{ \circ  }$ 、 ${45}^{ \circ  }$ 的角。则 ${BC}$ 到平面 $\alpha$ 的距离为\blkda{$\sqrt{6}$}.

\item 已知直角 $\bigtriangleup  {ABC}$ 的斜边 ${AB}$ 在平面 $M$ 内，而直角边 ${AC}$ ， ${BC}$ 分别与 $M$ 成 ${30}^{ \circ  }$ ， ${45}^{ \circ  }$ 角，则斜边上的高 ${CD}$ 与平面 $M$ 所成的角为\blkda{$\dfrac{\pi}{3}$}.

\item 已知 ${PA}\text{ 、 }{PB}\text{ 、 }{PC}$ 是从点 $P$ 发出的不共面的三条射线，每两条射线的夹角都是 ${60}^{ \circ  }$ ， 则直线 PC 与平面 ${PAB}$ 所成的角的余弦值为\blkda{$\dfrac{\sqrt{3}}{3}$}.

\item 到空间不共面的四点距离相等的平面有\blkda{7}个.

\item 已知长为 3 的线段 ${AB}$ 的两个端点到平面 ${\alpha }$ 的距离分别为 1 和 2，则直线 ${AB}$ 与平面 ${\alpha }$ 的所成角大小为\blkda{$arcsin\dfrac{1}{3}or\dfrac{\pi}{2}$}.

\item 已知 ${\theta }$ 是锐角，则 “直线 ${l}$ 与平面 ${\alpha }$ 所成角的大小为 ${\theta }$ ” 是 “直线 ${l}$ 与平面 ${\alpha }$ 内无数条直线所成角的大小为 ${\theta }$ ” 的\dotfill{(\kh{B})}条件.

\xx{必要不充分 }{充分不必要}
{充分必要} {既不充分也不必要}


\item 如图，已知正方体 ${ABCD} - {EFGH}$ 的棱长为 $a$ ，点 $P$ 在 ${AC}$ 上， $Q$ 在 ${BG}$ 上， ${AP} = {BQ} = a$ . 求直线 ${PQ}$ 与平面 ${ABCD}$ 所成角的大小.

\begin{center}
    \centering
    \includegraphics[width=0.25\linewidth]{bodyxb//src/hfa'suodhf'os.png}
\end{center}

\begin{jieda}
    过点$Q$作$BC$的垂线，交BC于点$E$,可证明$E$即为$Q$在平面$ABCD$的投影，因此 $\angle{QPE}$即为所求角，$|QE|=\dfrac{\sqrt{2}}{2}a,|EP|=(1-\dfrac{\sqrt{2}}{2})a$,$tan\angle{QPE}=\sqrt{2}+1$，综上所求角为$arctan(\sqrt{2}+1)$

    
    \begin{center}

        \includegraphics[width=0.25\linewidth]{bodyxb//src/ashfaohfaoijf.png}
    \end{center}
    
\end{jieda}

\item 如图，在 ${Rt\Delta BMC}$ 中，斜边 ${BM} = 5$ ，它在平面 ${ABC}$ 上的射影 $AB$ 长为 4， $\angle {MBC} = {60}^{ \circ  }$ ,求 MC 与平面 MAB 所成角的大小.

\begin{center}
    \centering
    \includegraphics[width=0.25\linewidth]{bodyxb//src/ahfa'hofuafa.png}
\end{center}

\begin{jieda}
    过C作AB的垂线，垂足为N，因为$CN// AB$,$CN// AM$，因此$CN// 面ABM$,所以$\angle{NMC}$为所成角. 因为$BC// CM$，$BC// AM$，因此$BC// 面ACM$,所以$BC// AC$，由等面积法$S_{\triangle ABC}=\dfrac{1}{2}|BC||AC|=\dfrac{1}{2}|AB||CN|$，求得$|CN|=\dfrac{5\sqrt{39}}{16}$. 又因为$|CM|=\dfrac{5\sqrt{3}}{2}$,所以$sin\angle{NMC}=\dfrac{CN}{CM}=\dfrac{\dfrac{5\sqrt{39}}{16}}{\dfrac{5\sqrt{3}}{2}}=\dfrac{\sqrt{13}}{8}$。综上，所求角为$arcsin \dfrac{\sqrt{13}}{8}$.
    \begin{center}
        \centering
        \includegraphics[width=0.25\linewidth]{bodyxb//src/fasuhfaohfa.png}
    \end{center}
\end{jieda}



\item 如图：在平面 $\beta$ 内有 $\bigtriangleup  {ABC}$ ，在平面 $\beta$ 外有点 $S$ ，斜线 ${SA}\bot {AC}$ ， ${SB}\bot {BC}$ ， 且斜线 ${SA}\text{ 、 }{SB}$ 分别与平面 $\beta$ 所成的角相等.
\begin{enumerate}
    \item 求证： ${AC} = {BC}$ ；
    \item 又设点 $S$ 与平面 $\beta$ 的距离是 ${4cm}$ ， ${AC}\bot {BC}$ ，且 ${AB} = {6cm}$ ， 求 ${SC}$ 与平面 $\beta$ 所成的角.
\end{enumerate}

\begin{center}
    \centering
    \includegraphics[width=0.5\linewidth]{bodyxb//src/faushf'ahfoaihjfa.png}
\end{center}

\begin{jieda}
    \begin{enumerate}
        \item 设S在$\beta$的投影为D，因为$\angle{SAD}=\angle{SBD}$，所以$SA=SB$,又因为$\angle{SAC}=\angle{SBC}$,所以$AC=BC$.
        \item $AC//面SAC$,所以$AC// AD$,同理$CB// BD$，又因为AC=BC,所以四边形ACBD为正方形. $CD=AB=6$,$SD=4$,所以$tan\angle{SCD}=\dfrac{4}{6}=\dfrac{2}{3}$.综上，所成角为$arctan\dfrac{2}{3}$
    \end{enumerate}
\end{jieda}

\item 在空间四边形 ${ABCD}$ 中， ${AB} = {CD} = 8$ ， $E$ 、 $F$ 分别为 ${AC}$ 、 ${BD}$ 的中点.

\begin{enumerate}
    \item 若 ${EF} = {4\sqrt{3}}$ ，求 ${AB}$ 与 ${CD}$ 所成的角；
    \item 若 ${AB}$ 与 ${CD}$ 所成的角为 ${60}^{ \circ  }$ ， 求 EF 的长.
\end{enumerate}

\begin{jieda}
    \begin{enumerate}
        \item 取BC的中点G，根据余弦定理，设所成角为$\theta$，有$cos\theta=|cos\angle{EGF}|=|\dfrac{4^2+4^2-(4\sqrt{3})^2}{2*4*4}|=\dfrac{1}{2}$，所以所成角为$\dfrac{\pi}{3}$.
        \item 当$\angle{EGF=\dfrac{\pi}{3}}$时，$EF=4$;当$\angle{EGF=\dfrac{2\pi}{3}}$时，$EF=4\sqrt{3}$.
    \end{enumerate}
\end{jieda}

\item 在空间四边形 ${ABCD}$ 中,已知 ${AB} = {BC} = {CD} = {DA} = {AC} = {BD} = a,E,F$ 分别是 ${AB},{CD}$ 的中点. 
\begin{enumerate}
    \item 求证: ${EF}$ 和 ${AD}$ 为异面直线;
    \item 求证: ${EF}$ 是异面直线 ${AB}$ 、 ${CD}$ 的公垂线；
    \item 求异面直线 ${AB}$ 、 ${CD}$ 间的距离.
\end{enumerate}

\begin{jieda}
    \begin{enumerate}
        \item 反证法，假设E、F、A、D共面，则$B\in AE\subset \alpha$,$C\in DF\subset \alpha, A\in \alpha, D\in \alpha$,从而$AD,BC\subset \alpha$,矛盾.
        \item 在等腰三角形AFB中，E为AB的中点，所以$EF// AB$,同理$EF// CD$,且EF与AB、CD均相交，因此是公垂线.

        \item $|EF|=\dfrac{a}{2}$,$|FG|=\dfrac{a}{2}$,$\angle{EGF}=\dfrac{\pi}{2}$ ,因此$d=|EF|=\dfrac{\sqrt{2}}{2}a$
    \end{enumerate}
\end{jieda}



\item 
\begin{enumerate}
    \item 如图，已知空间四边形 ${ABCD}$ 中， ${AB} = {CD} = 3$ ， $E\text{ 、 }F$ 分别是 ${BC}\text{ 、 }{AD}$ 上的点，且 ${BE} : {EC} = {AF} : {FD} = 1 : 2$ , ${EF} = 2\sqrt{2}$ ，求 ${AB}$ 和 ${CD}$ 所成角的大小.
    
\begin{center}
    \centering
    \includegraphics[width=0.25\linewidth]{bodyxb//src/fasha'sjdhf'asfh.png}
\end{center}
    
    \item 如图，已知空间四边形 ${ABCD}$ 中各边长和对角线长 $a,E\text{ 、 }F$ 分别为边 ${AD}\text{ 、 }{BC}$ 的中点，求 ${AF}$ 与 ${CE}$ 所成的角的大小.
    
\begin{center}
    \centering
    \includegraphics[width=0.25\linewidth]{afha'ouhfaof.png}
\end{center}

    \item 在空间四边形 ${ABCD}$ 中， ${AD} = 1$ ， ${BC} = \sqrt{3}$ ， ${AD}\bot {BC}$ ，对角线 ${BD} = \dfrac{\sqrt{13}}{2}$ ， ${AC} = \dfrac{\sqrt{3}}{2}$ ，求 ${AC}$ 与 ${BD}$ 所成角的大小
    
\begin{center}
    \centering
    \includegraphics[width=0.25\linewidth]{bodyxb//src/aisufghaiuhfa.png}
\end{center}
\end{enumerate}


\begin{jieda}
    \begin{enumerate}
        \item 取BD的靠近B的三等分点，$\angle{EMF}$即为所成角或其补角，$cos\angle{EMF}=\dfrac{-3}{4}$，综上所成角为$arccos\dfrac{3}{4}$

        \item 连接FD，取FD中点为O，$\angle{OEC}$即为所求角或其补角，\newline $cos\angle{OEC}=\dfrac{{\dfrac{\sqrt{3}a}{4}}^2+{\dfrac{\sqrt{3}a}{2}}^2-{\dfrac{\sqrt{7}a}{4}}^2}{2*\dfrac{\sqrt{3}a}{4}*\dfrac{\sqrt{3}a}{2}}=\dfrac{2}{3}$，综上所求即为$arccos\dfrac{2}{3}$.

        \begin{center}
    
            \includegraphics[width=0.25\linewidth]{bodyxb//src/aofuahsoduhfaouhfa.png}
        \end{center}

        \item 取AB、AC、AD、CD的中点分别为E、F、G、H，则$\angle{EGH}$为所成角，$|EG|=\dfrac{\sqrt{13}}{4}$,$|GH|=\dfrac{\sqrt{3}}{4}$,$|EH|=\sqrt{EF^2+FH^2}=1$,所以$cos\angle{EGH}=0$，故$\dfrac{\pi}{2}$

        \begin{center}
    
            \includegraphics[width=0.25\linewidth]{asuhfouhfaoh.png}
        \end{center}
    \end{enumerate}
\end{jieda}
    
\end{enumerate}


\clearpage



\subsection*{直线与平面的位置关系（4）}

知识点：三垂线定理

三垂线定理: 平面上的一条直线和这个平面的一条斜线垂直的充要条件是它和这条斜线在\blkda[3]{平面内的投影}垂直.

A 组：

\begin{enumerate}

\item 平面 $ a $ 外的直线 $ a $ 与平面 $ a $ 所成的角是 $ \theta $ ，则 $ \theta $ 的取值范围是\blkda{$[0,\dfrac{\pi}{2}]$}.

\item 在正四面体 $ {ABCD} $ 中，直线 $ {BC} $ 与 $ {AD} $ 所成角的大小为\blkda{$\dfrac{\pi}{2}$}.

3、定点 $ A $ 和 $ B $ 都在平面 $ \alpha $ 内,定点 $ P \notin  \alpha ,{PB} \bot  \alpha ,C $ 是平面 $ \alpha $ 内异于 $ A $ 和 $ B $ 的动点，且 $ {PC}\bot {AC} $ ，则 $ \bigtriangleup  {ABC} $ 为
\dotfill{(\kh{B})}

\xx{锐角三角形}{直角三角形} {钝角三角形} {无法确定}



\item 如图，如果 $ {MC}\bot $ 菱形 $ {ABCD} $ 所在平面，那么 $ {MA} $ 与 $ {BD} $ 的位置关系是\dotfill{(\kh{C})}
\xx{平行} {垂直相交} {垂直但不相交} {相交但不垂直}
\begin{center}
    \includegraphics[width=0.3\textwidth]{bodyxb//src/afh'auhf'aouhfa.png}
\end{center}


\item 在空间四边形 $ {ABCD} $ 中， $ E $ ， $ F $ ， $ G $ ， $ H $ 分别是边 $ {AB} $ ， $ {BC} $ ， $ {CD} $ ， $ {DA} $ 的中点， 若四边形对角线 $ {AC} = {BD} = 2 $ ，对角线 $ {AC} $ 与 $ {BD} $ 所成的角为 $ \dfrac{\pi }{3} $ ，则 $ {FH} = $ \blkda{$1or\sqrt{3}$}.

\end{enumerate}


B 组:
\begin{enumerate}

\item 直线 $ m $ 和平面 $ a $ 所成角为 $ \dfrac{\pi }{6} $ ，则直线 $ m $ 和平面 $ a $ 内任意直线所成角的取值范围是\blkda{$[\dfrac{\pi}{6},\dfrac{\pi}{2}]$}.

\item 在《九章算术》中，将底面为矩形且有一条侧棱与底面垂直的四棱锥称为 “阳马”. 如图，四棱锥 $ P - {ABCD} $ 为阳马，侧棱 $ {PA}\bot $ 底面 $ {ABCD} $ ， $ {PA} = {AB} = {AD} $ ， $ E $ 为棱 $ {PA} $ 的中点，则直线 $ {CE} $ 与平面 $ {PAD} $ 所成角的余弦值为\blkda{$\dfrac{\sqrt{5}}{3}$}.
\begin{center}
    \includegraphics[width=0.3\textwidth]{bodyxb//src/a'ofuh'aushdfa.png}
\end{center}



\item 如图，在坡面 $ a $ 与水平面 $ \beta $ 所成二面角为 $ {60}^{ \circ  } $ 的山坡上，有段直线型道路 $ {AB} $ 与坡脚 $ l $ 成 $ {30}^{ \circ  } $ 的角,这段路直通山顶 $ \mathrm{A} $ ，已知此山高 $ {135}\sqrt{3} $ 米，若小李从 $ B $ 沿着这条路上山，并且行进速度为每分钟 30 米，那么小李到达山顶 A 需要的时间是\blkda{18}分钟.
\begin{center}
    \includegraphics[width=0.3\textwidth]{askfgagfafasafa.png}
\end{center}


\item 已知点 $ P $ 是 $ \bigtriangleup  {ABC} $ 所在平面外一点，且 $ {PA} = {PB} = {PC} $ ，则点 $ P $ 在平面 $ {ABC} $ 上的射影一定是 $ \bigtriangleup  {ABC} $ 的\blkda{外}心.


\item 下列说法正确的是\dotfill{(\kh{B})}

\xx{过平面外一点作这个平面的垂直平面是唯一的}
{过平面的一条斜线作这个平面的垂直平面是唯一的}
{过直线外一点作这直线的平行平面是唯一的}
{过直线外一点作这直线的垂线是唯一的}



\item 如图， $ {AB} $ 为 $ \odot  O $ 的直径， $ {PA} $ 垂直于 $ \odot  O $ 所在的平面， $ M $ 为圆周上任意一点， $ {AN}\bot {PM} $ ， $ N $ 为垂足.

\begin{enumerate}
    \item 求证： $ {AN}\bot $ 平面 $ {PBM} $ ；
    \item 若 $ {AQ}\bot {PB} $ ，垂足为 $ Q $ ，求证： $ {NQ}\bot {PB} $ .
\end{enumerate}

\begin{center}
    \includegraphics[width=0.3\textwidth]{asfashf'oauhfasd.png}
\end{center}


\begin{jieda}
    \begin{enumerate}
        \item $ \because \mathrm{{AB}} $ 为 $ \odot  \mathrm{O} $ 的直径, $ \therefore \mathrm{{AM}}\bot \mathrm{{BM}} $ .
又 $ \mathrm{{PA}} \bot $ 平面 $ \mathrm{{ABM}} $ , $ \therefore \mathrm{{PA}} \bot  \mathrm{{BM}} $ ,
又 $ \because \mathrm{{PA}} \cap  \mathrm{{AM}} = \mathrm{A},\therefore \mathrm{{BM}}\bot $ 平面 $ \mathrm{{PAM}} $ .
又 $ \mathrm{{AN}} \subset $ 平面 $ \mathrm{{PAM}},\therefore \mathrm{{BM}}\bot \mathrm{{AN}} $ .
又 $ \mathrm{{AN}} \bot  \mathrm{{PM}} $ ,且 $ \mathrm{{BM}} \cap  \mathrm{{PM}} = \mathrm{M} $ ,
$ \therefore \mathrm{{AN}} \bot $ 平面 $ \mathrm{{PBM}} $ .
    \item 由(1)知 $ \mathrm{{AN}}\bot $ 平面 $ \mathrm{{PBM}},\mathrm{{PB}} \subset $ 平面 $ \mathrm{{PBM}},\therefore \mathrm{{AN}}\bot \mathrm{{PB}} $ .
又 $ \because \mathrm{{AQ}}\bot \mathrm{{PB}},\mathrm{{AN}} \cap  \mathrm{{AQ}} = \mathrm{A} $ ,
$ \therefore \mathrm{{PB}}\bot $ 平面 $ \mathrm{{ANQ}} $ .又 $ \mathrm{{NQ}} \subset $ 平面 $ \mathrm{{ANQ}} $ .
$ \therefore \mathrm{{PB}}\bot \mathrm{{NQ}} $ .
    \end{enumerate}
\end{jieda}



\item 已知长方体 $ {ABCD} - {A}_{1}{B}_{1}{C}_{1}{D}_{1} $ 中， $ {AD} = {AB} = 4 $ ， $ A{A}_{1} = 2,E,F $ 分别是 $ {AB},{A}_{1}{D}_{1} $ 的中点.

\begin{enumerate}
    \item 求证：直线 $ {EF}// $ 平面 $ B{B}_{1}{D}_{1}D $ ；
    \item 求直线 $ {EF} $ 与平面 $ {BC}{C}_{1}{B}_{1} $ 所成角的正弦值.
\end{enumerate}

\begin{center}
    \includegraphics[width=0.3\textwidth]{bodyxb//src/aufhga'uhsf'afa.png}
\end{center}

\begin{jieda}
    \begin{enumerate}
        \item 取$BD_1$的中点G，连接FG和BG，可以证明$FE//BG$,从而证明直线 $ {EF}// $ 平面 $ B{B}_{1}{D}_{1}D $ .

        \item 取$B_1C_1$的中点M，可证明GB在平面$ {BC}{C}_{1}{B}_{1} $的投影为$BM$，因此$\angle{GBM}$即为所求线面角，$|GM|=2$,$|MB|=2\sqrt{2}$，因此$sin\angle{GBM}=\dfrac{\sqrt{3}}{3}$。
    \end{enumerate}
    \begin{center}
    \includegraphics[width=0.3\textwidth]{ao'uhfaoushdfa'df.png}
\end{center}
\end{jieda}


\item 如图，在四棱锥 $ P - {ABCD} $ 中，底面 $ {ABCD} $ 是矩形.已知 $ {AB} = 3,{AD} = 2,{PA} = 2,{PD} = 2\sqrt{2},\angle {PAB} = {60}^{ \circ  } $ . 
\begin{enumerate}
    \item 证明： $ {AD}\bot $ 平面 $ {PAB} $ ；
    \item 求异面直线 $ {PC} $ 与 $ {AD} $ 所成角的大小.
\end{enumerate}

\begin{center}
    \includegraphics[width=0.3\textwidth]{'asofuhaouihf'ashf.png}
\end{center}


\begin{jieda}
    \begin{enumerate}
        \item 由于四边形 $ {ABCD} $ 是矩形,所以 $ {AD}\bot {AB} $ ,
由于 $ P{A}^{2} + A{D}^{2} = P{D}^{2} $ ,所以 $ {AD}\bot {PA} $ ,
由于 $ {AB} \cap  {PA} = A,{AB},{PA} \subset $ 平面 $ {PAB} $ ,所以 $ {AD}\bot $ 平面 $ {PAB} $ .
        \item 由于 $ {AD}//{BC} $ ,所以 $ {BC}\bot $ 平面 $ {PAB} $ ,
由于 $ {PB} \subset $ 平面 $ {PAB} $ ,所以 $ {BC} \bot  {PB} $ ,
所以 $ \angle {PCB} $ 是异面直线 $ {PC} $ 与 $ {AD} $ 所成角,
$ {PB} = \sqrt{{2}^{2} + {3}^{2} - 2 \times  2 \times  3 \times  \cos {60}^{ \circ  }} = \sqrt{7}, $
所以 $ \tan \angle {PCB} = \dfrac{\sqrt{7}}{2} $ ,由于 $ \angle {PCB} $ 是锐角,所以 $ \angle {PCB} = \arctan \dfrac{\sqrt{7}}{2} $ .
    \end{enumerate}
\end{jieda}

\item 如图， $ \bigtriangleup  {{ACD}} $ 和 $ \bigtriangleup  {{BCD}} $ 都是边长为 2 的等边三角形，

$ {{AB}} = \sqrt{{6}},{{EB}} \bot $ 平面 $ {BCD} $ .

\begin{enumerate}
    \item 证明： $ {EB}// $ 平面 $ {ACD} $ ；
    \item 求直线 $ {EB} $ 到平面 $ {ACD} $ 的距离。
\end{enumerate}

\begin{figure}
    \centering
    \includegraphics[width=0.25\linewidth]{asfudha'uhfawdfa.png}
\end{figure}

\begin{jieda}
    \begin{enumerate}
        \item 由于四边形 $ {ABCD} $ 是矩形,所以 $ {AD}\bot {AB} $ ,
由于 $ P{A}^{2} + A{D}^{2} = P{D}^{2} $ ,所以 $ {AD}\bot {PA} $ ,
由于 $ {AB} \cap  {PA} = A,{AB},{PA} \subset $ 平面 $ {PAB} $ ,所以 $ {AD}\bot $ 平面 $ {PAB} $ .
        \item  由第一问的证明，线到面的距离即点B到平面ACD的距离.取CD的中点O，因为BO垂直CD和AO，因此$BO// ACD$,所以$d=|BO|=\sqrt{3}$
    \end{enumerate}

    \begin{center}
    \includegraphics[width=0.3\textwidth]{‘faoushf'aouhfasdf.png}
\end{center}
\end{jieda}





\item 如图，已知 $ A{A}_{1} \bot $ 平面 $ {ABC},B{B}_{1}//A{A}_{1},{AB} = {AC} = 3 $ ， $ {BC} = {2\sqrt{5}},{A{A}_{1}} = \sqrt{7},{B{B}_{1}} = {2\sqrt{7}} $ ，点 $ E $ 和 $ F $ 分别为 $ {BC} $ 和 $ {A}_{1}C $ 的中点.

\begin{enumerate}
    \item 求证： $ {AE}\bot $ 平面 $ {BC}{B}_{1} $ ；
    \item 求直线 $ {A}_{1}{B}_{1} $ 与平面 $ {BC}{B}_{1} $ 所成角的大小.
\end{enumerate}


\begin{center}
    \includegraphics[width=0.3\textwidth]{asufha'suodhfafa.png}
\end{center}


\begin{jieda}
    \begin{enumerate}
        \item 因为$AE// BB_1$,$AE// BC$,因此$AE// BCB_1$
        
        \item 如图,取 $ B{B}_{1} $ 的中点 $ M $ 和 $ {B}_{1}C $ 的中点 $ N $ ,连
接 $ {A}_{1}M,{A}_{1}N,{NE} $ .
因为 $ N $ 和 $ E $ 分别为 $ {B}_{1}C $ 和 $ {BC} $ 的中点,所以 $ {NE}// $  $ {B}_{1}B,{NE} = \dfrac{1}{2}{B}_{1}B $ ,故 $ {NE}//{A}_{1}A $ 且 $ {NE} = {A}_{1}A $ ,
所以四边形 $ {NEA}{A}_{1} $ 为平行四边形,所以 $ {A}_{1}N//{AE} $ 且 $ {A}_{1}N = {AE} $ .
因为 $ {AB} = {AC},E $ 为 $ {BC} $ 的中点,所以 $ {AE}\bot {BC},{A}_{1}N\bot {BC} $ .
因为 $ A{A}_{1} \bot $ 平面 $ {ABC},B{B}_{1}//A{A}_{1} $ ,则易证得 $ B{B}_{1} \bot $ 平面 $ {ABC} $ .
因为 $ {AE} \subset $ 平面 $ {ABC} $ ,所以 $ B{B}_{1} \bot  {AE},{A}_{1}N \bot  B{B}_{1} $ ,
又 $ B{B}_{1} \cap  {BC} = B $ ,所以 $ {A}_{1}N \bot $ 平面 $ {BC}{B}_{1} $ ,
从而 $ \angle {A}_{1}{B}_{1}N $ 为直线 $ {A}_{1}{B}_{1} $ 与平面 $ {BC}{B}_{1} $ 的夹角.
在 $ \bigtriangleup {ABC} $ 中,可得 $ {AE} = 2 $ ,所以 $ {A}_{1}N = {AE} = 2 $ .
因为 $ {BM}//{AA}_{1},{BM} = {AA}_{1} $ ,所以 $ {A}_{1}M//{AB},{A}_{1}M = {AB} $ ,又由
$ {AB} \bot  B{B}_{1} $ ,得 $ {A}_{1}M \bot  B{B}_{1} $ .
在 $ \mathrm{{Rt}}\bigtriangleup {A}_{1}M{B}_{1} $ 中,可得 $ {A}_{1}{B}_{1} = \sqrt{{B}_{1}{M}^{2} + {A}_{1}{M}^{2}} = 4 $ .
在 $ \mathrm{{Rt}}\bigtriangleup {A}_{1}N{B}_{1} $ 中, $ \sin \angle {A}_{1}{B}_{1}N = \dfrac{{A}_{1}N}{{A}_{1}{B}_{1}} = \dfrac{1}{2} $ ,因此 $ \angle {A}_{1}{B}_{1}N = {30}^{ \circ  } $ .
所以直线 $ {A}_{1}{B}_{1} $ 与平面 $ {BC}{B}_{1} $ 的夹角为 $ {30}^{ \circ  } $ .
    \end{enumerate}
    \begin{center}
    \includegraphics[width=0.3\textwidth]{aoufuh'aouhf'aodfa.png}
\end{center}
\end{jieda}

\item 如图，已知 $ {P} $ 是平面 $ {A}{B}{C} $ 外一点， $ {P}{A} \bot $ 平面 $ {ABC},{AC} \bot  {BC} $ .

\begin{enumerate}
    \item 证明： $ {{BC}} \bot $ 平面 $ {{PAC}} $ ；
    \item 过 $ {A} $ 点作 $ {{AD}} $ 垂直 $ {{PC}} $ 于 $ {D} $ ，证明： $ {{AD}} \bot  {{PB}} $ ； 
    \item 若 $ {{BC}} = {1},{{AC}} = \sqrt{{3}} $ ，求点 $ {C} $ 到平面 $ {{PAB}} $的距离.
\end{enumerate}

\begin{center}
    \includegraphics[width=0.3\textwidth]{afouh'aoushfadfa.png}
\end{center}

\begin{jieda}
    \begin{enumerate}
        \item 因为 $ {PA}\bot $ 平面 $ {ABC} $ , $ {BC} \subset $ 平面 $ {ABC} $ ,所以 $ {PA}\bot {BC} $ ,
又 $ {AC}\bot {BC},{PA} \cap  {AC} = A,{PA},{AC} \subset $ 平面 $ {PAC} $ ,
所以 $ {BC} \bot $ 平面 $ {PAC} $ ;
        \item 因为 $ {BC}\bot $ 平面 $ {PAC},{AD} \subset $ 平面 $ {PAC} $ ,
所以 $ {BC}\bot {AD} $ ,又 $ {AD}\bot {PC},{PC} \cap  {BC} = C,{PC},{BC} \subset $ 平面 $ {PBC} $ ,
所以 $ {AD}\bot $ 平面 $ {PBC} $ ,又 $ {PB} \subset $ 平面 $ {PBC} $ ,所以 $ {AD}\bot {PB} $ ；
    \item 在平面 $ {ABC} $ 内过点 $ C $ 作 $ {CE}\bot {AB} $ 交 $ {AB} $ 于点 $ E $ ,
因为 $ {BC} = 1,{AC} = \sqrt{3},{AC}\bot {BC} $ ,所以 $ {AB} = \sqrt{{1}^{2} + {\left( \sqrt{3}\right) }^{2}} = 2 $ ’
所以 $ {CE} = \dfrac{{AC} \cdot  {BC}}{AB} = \dfrac{\sqrt{3}}{2} $ ,
又 $ {PA}\bot $ 平面 $ {ABC},{CE} \subset $ 平面 $ {ABC} $ ,所以 $ {PA}\bot {CE} $ ,
又 $ {PA} \cap  {AB} = A,{PA},{AB} \subset $ 平面 $ {PAB} $ ,所以 $ {CE}\bot $ 平面 $ {PAB} $ ,
所以点 $ C $ 到平面 $ {PAB} $ 的距离为 $ {CE} = \dfrac{\sqrt{3}}{2} $ .

    \begin{center}
    \includegraphics[width=0.3\textwidth]{asdfhauhf'ads.png}
\end{center}
    \end{enumerate}
\end{jieda}



\end{enumerate}





\section*{平面与平面的位置关系}


\subsection*{平面与平面的位置关系（1）}

知识点：面面平行

A 组：
\begin{enumerate}

\item 判断下列命题是否正确：

    \begin{enumerate}
        \item 若一个平面内有两条直线平行于另一个平面，则这两个平面平行；\blkda{\ding{55}}

        \item 若一个平面内有无数条直线平行于另一个平面，则这两个平面平行； \blkda{\ding{55}}

        \item 若一个平面内任意一条直线平行于另一个平面，则这两个平面平行； \blkda{\ding{51}}

        \item 若一个平面内有两条相交直线，分别平行于另一个平面内的两条直线，则这两个平面平行; \blkda{\ding{51}}

        \item 若一个平面内一个角（锐角或钝角）的两边和另一个平面内的一个角的两边分别平行， 则这两个平面平行. \blkda{\ding{51}}

        
    \end{enumerate}

    \item 已知 $ a \subset  a $ ，则 $ a//\beta $ 是 $ a//\beta $ 的 \blkda[3]{充分不必要}条件

    \item 已知 $ m,n $ 为两条不同直线， $ a,\beta ,\gamma $ 为三个不同平面，则下列说法正确的是\dotfill{(\kh{D})}
    \xx{若 $ m//a,n \subset  a $ ，则 $ m//n $ } {若 $ m//a,a//\beta $ ，则 $ m//\beta $}
    {若 $ m//a,m//\beta $ ，则 $ a//\beta $} {若 $ a//\beta ,\beta //\gamma $ ，则 $ a//\gamma $}

    \item 过平面 $ a $ 外一条直线 $ a $ 与 $ a $ 平行的平面的个数是 \dotfill{(\kh{C})}

    \xx{只有一个} {至少有一个} {至多有一个} { 一个也没有}

    \item 经过平面外一点可以作\blkda{1}个平面平行于这个平面；可以作\blkda{无数}条直线平行于这个平面.
    
    \item 分别在两个平面内的三角形都有两条边与同一条直线垂直，则这两个三角所在面的位置关系是\blkda{平行}.

\end{enumerate}


B组：
\begin{enumerate}

    \item 下列条件中，可以判断平面 $ \alpha $ 与 $ \beta $ 平行的是 \dotfill{(\kh{D})}

    \xx{直线 $ a $ 与 $ a,\beta $ 相交且所成的角相等}
    {$ \alpha $ 内存在不共线的三点到 $ \beta $ 的距离相等}
    {l, $ {m} $ 是 $ \alpha $ 内两条直线,且 $ l//\beta ,{m}//\beta $}
    {$ l,m $ 是两条异面直线，且 $ l//\alpha ,m//\alpha ,l//\beta ,m//\beta $}

    \item 已知平面 $ \alpha // $ 平面 $ \beta $ ，直线 $ a\text{ 、 }b $ 分别与 $ \alpha \text{ 、 }\beta $ 所成的角相等，则直线 $ a\text{ 、 }b $ 的位置关系是\blkda[3]{相交、平行或异面}.

    \item 如图，平面 $ a// $ 平面 $ \beta $ ，自点 $ {O} $ 引三条直线分别交 $ \alpha $ 、 $ \beta $ 于点 $ {A} $ 、B、C 和点 $ {A}_{1} $ 、 $ {B}_{1} $ 、 $ {C}_{1} $ ，则 $ \bigtriangleup  {ABC} $ 与 $ \bigtriangleup  {A}_{1}{B}_{1}{C}_{1} $ 的关系是\blkda{相似}.

    \begin{center}
    
        \includegraphics[width=0.25\linewidth]{bodyxb//src/a'ouhsf'ouahs'fouhasodf.png}
    \end{center}

    
    \item 已知 $ {a}\text{ 、 }{b}\text{ 、 }{c} $ 为三条不重合的直线， $ a,\beta ,\gamma $ 为三个不重合的平面， 下列六个命题中正确的命题是\blkda{(1)(4)}.

    \begin{enumerate}
        \item $ a//c,b//c \Rightarrow  a//b; $ 
        \item $ a//\gamma ,b//\gamma  \Rightarrow  a//b $ ;
        \item  $ {c}//a,{c}//\beta  \Rightarrow  a//\beta ; $
        \item $ a//\gamma ,\beta //\gamma  \Rightarrow  a//\beta $ ;
        \item $ {a}//{c},{c}//a \Rightarrow  {a}//a; $ 
        \item $ {a}//\gamma ,a//\gamma  \Rightarrow  {a}//a $ .
    \end{enumerate}


    \item 已知平面 $ a//\beta $ ，点 $ A,C \in  a,B,D \in  \beta $ ，直线 $ {{AB}} $ 与 $ {{CD}} $ 交于点 $ {S} $ ，且 $ {{AS}} = 8,{{BS}} = 9 $ ， CD=34，则 SC 的长度是\blkda[3]{272或16}.

    \item 已知平面 $ a//\beta $ ，且 $ a,\beta $ 的距离为 $ {d} $ ， $ a \subset  a $ ，则在平面 $ \beta $ 中 \dotfill{(\kh{C})}
    \xx{有且只有一条直线与 $ {a} $ 的距离为 $ {d} $}
    {所有直线与 $ {a} $ 的距离都等于 $ {d} $}
    {有无数条直线与 $ a $ 的距离等于 $ d $}
    {所有直线与 $ a $ 距离都不等于 $ d $}

    \item 已知 AB、CD 是夹在两平行平面 $ a,\beta $ 之间的两条线段， $ {AB}\bot {CD},{AB} = 2,{AB} $ 与平面 $ a $ 成 $ {30}^{ \circ  } $ 角，则线段 CD 的取值范围是 \dotfill{(\kh{B})}

    \xx{$ \left( {\dfrac{2\sqrt{3}}{3},2\sqrt{3}}\right) $} {$ \left\lbrack  {\dfrac{2\sqrt{3}}{3}, + \infty }\right) $} {$ \left( {1,\dfrac{2\sqrt{3}}{3}}\right) $ }{$ \lbrack 1, + \infty ) $}
    
    \item 已知 $ {P} $ 是 $ \bigtriangleup  {ABC} $ 所在平面外一点， $ {A}_{1} $ 、 $ {B}_{1} $ 、 $ {C}_{1} $ 分别是 $ \bigtriangleup  {PBC} $ 、 $ \bigtriangleup  {PCA} $ 、 $ \bigtriangleup  {PAB} $ 的重心. 
    \begin{enumerate}
        \item 求证: 平面 $ {A}_{1}{B}_{1}{C}_{1}// $ 平面 $ {{ABC}} $ ;
        \item 求 $ {S}_{\bigtriangleup {A}_{1}{B}_{1}{C}_{1}} : {S}_{\bigtriangleup {ABC}}$.
    \end{enumerate}

    \begin{jieda}
    \begin{enumerate}

        \item 连接 $ P{A}^{\prime } $ 、 $ P{C}^{\prime \prime } $ ,并延长分别交BC、AB于点M,N,连接MN. $ \because {A}^{\prime }\text{、}{C}^{\prime \prime } $ 分别是 $ \bigtriangleup  {{PBC}} $ 、 $ \bigtriangleup  {{PAB}} $ 的重心,
$ \therefore P{A}^{\prime } = \dfrac{2}{3}{PM},P{C}^{\prime \prime } = \dfrac{2}{3}{PN} $ ,
$ \therefore {A}^{\prime }{C}^{\prime \prime }// {MN} $ .
$ \because {A}^{\prime }{C}^{\prime \prime } \not \subset$ 平面 $ {{ABC}},{MN} \subset $ 平面 $ {{ABC}} $ ,
$ \therefore {A}^{\prime }{C}^{\prime \prime }// $ 平面 $ {{ABC}} $ .
同理, $ {A}^{\prime }{B}^{\prime \prime }// $ 平面 $ {{ABC}} $ .
$ \because {A}^{\prime }{C}^{\prime \prime } \cap  {A}^{\prime }{B}^{\prime } = {A}^{\prime } $ ,且 $ {A}^{\prime }{C}^{\prime \prime }\text{、}{A}^{\prime }{B}^{\prime } \subset $ 平面 $ {A}^{\prime }{B}^{\prime }{C}^{\prime \prime } $ ,
$ \therefore $ 平面 $ {A}^{\prime }{B}^{\prime }{C}^{\prime \prime }// $ 平面 $ {ABC} $ .
        \begin{center}
            \includegraphics[width=0.25\linewidth]{bodyxb//src/auhf'aoshfasd.png}
        \end{center}

\item 由(1)知 $ {A}^{\prime }{C}^{\prime \prime } = \dfrac{2}{3}{MN},{A}^{\prime }{C}^{\prime \prime }// {MN} $ .
$ \because {MN} = \dfrac{1}{2}{AC},{MN}// {AC} $ ,
$ \therefore {A}^{\prime }{C}^{\prime \prime } = \dfrac{1}{3}{AC},{A}^{\prime }{C}^{\prime \prime }// {AC} $ .
同理可得： $ {A}^{{}^{\prime }}{B}^{{}^{\prime }} = \dfrac{1}{3}{AB} $ , $ {A}^{{}^{\prime }}{B}^{{}^{\prime }} //  {AB} $ , $ {B}^{{}^{\prime }}{C}^{{}^{\prime \prime }} = \dfrac{1}{3}{BC} $ , $ {B}^{{}^{\prime }}{C}^{{}^{\prime \prime }} //  {BC} $ ,
则 $ \Delta {A}^{\prime }{B}^{\prime }{C}^{\prime \prime } \sim  {\Delta ABC} $ .
故 $ \Delta {A}^{\prime }{B}^{\prime }{C}^{\prime \prime } $ 与 $ \bigtriangleup  {{ABC}} $ 的面积之比为 $ \dfrac{{S}_{\Delta {A}^{\prime }{B}^{\prime }{C}^{\prime \prime }}}{{S}_{\Delta ABC}} = \dfrac{1}{9} $ .
    \end{enumerate}
    \end{jieda}
    \vspace{3cm}



    \item 如图，在四棱锥 $ C - {ABED} $ 中，四边形 $ {ABED} $ 是正方形，点 $ G $ ， $ F $ 分别是线段 $ {EC} $ ， $ {BD} $ 的中点。

    \begin{enumerate}
        \item 求证： $ {GF}// $ 平面 $ {ABC} $ ；
        \item 若 $ H $ 是线段 $ {BC} $ 的中点，证明：平面 $ {GFH}// $ 平面 $ {ACD} $ .
    \end{enumerate}
    
\begin{center}

    \includegraphics[width=0.25\linewidth]{auufh'aoushf'aouhfa.png}
\end{center}

\begin{jieda}
    \begin{enumerate}
        \item 证明：由四边形 $ {ABED} $ 为正方形可知,
        连接 $ {AE} $ 必与 $ {BD} $ 相交于中点 $ F $ ,故 $ {GF} //  {AC} $ ,
        $ \because {GF} \not \subset$ 面 $ {ABC},{AC} \subset $ 平面 $ {ABC} $ ,
        $ \therefore {GF} // $ 面 $ {ABC} $ .

                \begin{center}
                
                    \includegraphics[width=0.25\linewidth]{bodyxb//src/au;ghf;aighf'aisfa.png}
                \end{center}

                \item 由点 $ G,H $ 分别为 $ {CE},{CB} $ 中点可得： $ {GH} //  {EB} //  {AD} $ ,
        $ \because {GH} \not \subset$ 面 $ {ACD},{AD} \subset $ 平面 $ {ACD},\therefore {GH} // $ 平面 $ {ACD} $ ,
        又由( 1 )可知, $ {GF} // $ 平面 $ {ACD} $ ,
        且 $ {GH} \cap  {GF} = G,{GH},{GF} \subset $ 平面 $ {GFH} $ ,
        故平面 $ {GFH}// $ 平面 $ {ACD} $ .
    \end{enumerate}
\end{jieda}


\item 如图，在四棱锥 $ P - {ABCD} $ 中， $ {PA} = 3 $ ， $ {AB} = 2 $ ，四边形 $ {ABCD} $ 为菱形， $ \angle {ABC} = \dfrac{\pi }{3} $ ， $ {PA}\bot $ 平面 $ {ABCD} $ ,E,F,Q 分别是 BC,PC,PD 的中点. 证明: 平面 $ {EFQ}// $ 平面 $ {PAB} $ ;

\begin{jieda}
    因为四边形 $ {ABCD} $ 为菱形,所以 $ {AB}//{CD} $ ,
又 $ E,F,Q $ 分别是 $ {BC},{PC},{PD} $ 的中点,
所以 $ {FQ}//{CD},{EF}//{PB} $ ,故 $ {FQ}//{AB} $ ,
因为 $ {EF} \not \subset$ 平面 $ {PAB},{PB} \subset $ 平面 $ {PAB} $ ,
所以 $ {EF}// $ 平面 $ {PAB} $ ,同理可得 $ {FQ}// $ 平面 $ {PAB} $ .
因为 $ {EF} \cap  {FQ} = F,{EF},{FQ} \subset $ 平面 $ {EFQ} $ ,
所以平面 $ {EFQ}// $ 平面 $ {PAB} $ .
\end{jieda}


\item 如图，已知 $ P $ 是平行四边形 $ {ABCD} $ 所在平面外一点， $ M $ 、 $ N $ 分别是 $ {AB} $ 、 $ {PC} $ 的三等分点 $ \left( {M\text{ 靠近 }B,N\text{ 靠近 }C}\right) $ ；

\begin{enumerate}
    \item 求证： $ {MN}// $ 平面 $ {PAD} $ .
    \item 在 $ {PB} $ 上确定一点 $ Q $ ，使平面 $ {MNQ}// $ 平面 $ {PAD} $ ，并证明.
\end{enumerate}

\begin{center}

    \includegraphics[width=0.25\linewidth]{bodyxb//src/ahkfa'khfg'afsd.png}
\end{center}

\begin{jieda}
    \begin{enumerate}
        \item 证明：过点 $ N $ 作 $ {NE}//{CD} $ ,交 $ {PD} $ 于点 $ E $ ,连接 $ {AE} $ ,
因为 $ N $ 为 $ {PC} $ 的三等分点,可得 $ {NE} = \dfrac{2}{3}{CD} $ ,
又因为 $ M $ 为 $ {AB} $ 的三等分点,可得 $ {AM} = \dfrac{2}{3}{AB} $ ,
因为 $ {AB}//{CD} $ 且 $ {AB} = {CD} $ ,所以 $ {AM}//{NE} $ 且 $ {AM} = {NE} $ ,
所以四边形 $ {AMNE} $ 为平行四边形,所以 $ {MN}//{AE} $ ,
又由 $ {MN} \not \subset$ 平面 $ {PAD},{AE} \subset $ 平面 $ {PAD} $ ,所以 $ {MN}// $ 平面 $ {PAD} $ .
        \item 证明：取 $ {PB} $ 取一点 $ Q $ ,使得 $ {BQ} = \dfrac{1}{3}{BP} $ ,即点 $ Q $ 为 $ {PB} $ 上靠近点 $ B $ 的三等点,
在 $ {\Delta PAB} $ 中,因为 $ M,Q $ 分别为 $ {AB},{PB} $ 的三等分点,可得 $ \dfrac{MB}{AB} = \dfrac{BQ}{BP} $ ,所以 $ {MQ}//{PA} $ ,
因为 $ {MQ} \not \subset$ 平面 $ {PAD},{PA} \subset $ 平面 $ {PAD} $ ,所以 $ {MQ}// $ 平面 $ {PAD} $ ;
又由 (1) 知 $ {MN}// $ 平面 $ {PAD} $ ,且 $ {MN} \cap  {MQ} = M,{MN},{MQ} \subset $ 平面 $ {MNQ} $ ,
所以平面 $ {MNQ}// $ 平面 $ {PAD} $ ,
即当点 $ Q $ 为 $ {PB} $ 上靠近点 $ B $ 的三等点时,能使得平面 $ {MNQ}// $ 平面 $ {PAD} $ .

\begin{center}

    \includegraphics[width=0.25\linewidth]{bodyxb//src/aufhaosufha'ouhdf.png}
\end{center}
    \end{enumerate}
\end{jieda}



\item 如图，在四棱锥 $ P - {ABCD} $ 中，四边形 $ {ABCD} $ 是正方形， $ {PA} = {PC} $ ，E 为侧棱 PD 上的点，且 $ {PE} = {3ED} $ .
\begin{enumerate}
    \item 证明： $ {PD} \bot  {AC} $ ；
    \item 在侧棱PC上是否存在一点 $ {F} $ ，使得 $ {BF}// $ 平面 $ {ACE} $ ？若存在,求 $ \dfrac{PC}{PF} $ 的值; 若不存在,请说明理由.

\end{enumerate}

\begin{center}

    \includegraphics[width=0.25\linewidth]{bodyxb//src/a;kjsfh'akhf'ashdfds.png}
\end{center}

\begin{jieda}
    \begin{enumerate}
        \item 设 $ {BD} $ 交 $ {AC} $ 于点 $ O $ ,连接 $ {OP} $ ,正方形 $ {ABCD} $ 中,则 $ {AO} = {OC} $ , $ {BD}\bot {AC} $ ,
又 $ {PA} = {PC} $ ,则 $ {PO}\bot {AC} $ ,又 $ {PO} \cap  {BD} = O,{PO},{BD} \subset $ 平面 $ {PBD} $ ,
因此 $ {AC} \bot $ 平面 $ {PBD} $ ,而 $ {PD} \subset $ 平面 $ {PBD} $ ,
所以 $ {PD} \bot  {AC} $ .
        
        \begin{center}
        
            \includegraphics[width=0.25\linewidth]{bodyxb//src/;yafga;iygf;afawsfa.png}
        \end{center}

        \item 侧棱 $ {PC} $ 上存在一点 $ F $ ,满足条件,
证明如下：如图,正方形 $ {ABCD} $ 中, $ {BO} = {OD} $ ,
在线段 $ {PE} $ 取一点 $ G $ ,使得 $ {GE} = {ED} $ ,由 $ {PE} = {3ED} $ ,得 $ \dfrac{PG}{PE} = \dfrac{2}{3} $ ,
连接 $ {BG},{FG} $ ,则 $ {BG}//{OE} $ ,而 $ {OE} \subset $ 平面 $ {ACE},{BG} \not \subset$ 平面 $ {ACE} $ ,
$ {则BG}// $ 平面 $ {ACE} $ ,由 $ {BF}// $ 平面 $ {ACE},{BG} \cap  {BF} = B,{BG},{BF} \subset $ 平面 $ {BGF} $ ,
得平面 $ {BGF}// $ 平面 $ {ACE} $ ,而平面 $ {PCD} \cap $ 平面 $ {ACE} = {EC} $ ,平面 $ {PCD} \cap $ 平面 $ {BGF} = {GF} $ ,
于是 $ {GF}//{EC},\dfrac{PF}{PC} = \dfrac{PG}{PE} = \dfrac{2}{3} $ ,
所以 $ \dfrac{PC}{PF} = \dfrac{3}{2} $ .


        \begin{center}
        
            \includegraphics[width=0.25\linewidth]{bodyxb//src/asufha'ufh'auwsdfa.png}
        \end{center}

    \end{enumerate}
\end{jieda}


\item 已知 $ \alpha//\beta //\gamma $ ，A、D $ \in  a $ ，C、F $ \in  \gamma $ ，AC、DF 分别交平面 $ \beta $ 于 $ {B} $ 、E. 
\begin{enumerate}
    \item 求证: $ \dfrac{AB}{BC} = \dfrac{DE}{EF} $ ;
    \item 设 $ {{AF}} $ 交 $ \beta $ 于 $ {M} $ ,若 $ \dfrac{AB}{BC} = 1 $ ，异面直线 $ {{AD}} $ 与 $ {{CF}} $ 所成角为 $ {60}^{ \circ  },{{AD}} = 2,{{CF}} = 4 $ ，求 $ {S}_{\bigtriangleup {BEM}} $ .
\end{enumerate}

\begin{center}

    \includegraphics[width=0.25\linewidth]{bodyxb//src/fa'ufha'ouhfa'odfs.png}
\end{center}

\begin{enumerate}
    \item 连接AF交平面β于G,连接AD,BE,CF,BG,EG.
由 $ \beta //\gamma $ ,平面 $ {{ACF}} \cap  \beta  = {{BG}} $ ,平面 $ {{ACF}} \cap  \gamma  = {{CF}} $ ,所以 $ {{BG}}//{{CF}} $ ,则 $ \dfrac{{{AB}}}{{{BC}}} = \dfrac{{{AG}}}{{{GF}}} $ ,
同理,由 $ \alpha //\beta $ ,可得 $ {{GE}}//{{AD}} $ ,则 $ \dfrac{{{DE}}}{{{EF}}} = \dfrac{{{AG}}}{{{GF}}} $ .
所以 $ \dfrac{{{AB}}}{{{BC}}} = \dfrac{{{DE}}}{{{EF}}} $ ；
    \item 因为BG//CF,GE//AD,所以 $ \angle $ BGE(或其补角)就是直线AD与CF所成的角.
因为 $ \dfrac{{{AB}}}{{{AC}}} = \dfrac{{{BG}}}{{{CF}}} = \dfrac{1}{2},\dfrac{{{GF}}}{{{AF}}} = \dfrac{{{GE}}}{{{AD}}} = \dfrac{1}{2} $ ,所以 $ {{BG}} = 2,{{GE}} = 1 $ ,
又 $ {{BE}} = \sqrt{7},{{CF}} = 4 $ ,
由余弦定理可得 $ \cos \angle {{BGE}} = \dfrac{1 + 4 - 7}{4} =  - \dfrac{1}{2} $ ,得 $ \angle {{BGE}} = {120}^{ \circ  } $ .
即直线 $ {{AD}} $ 与 $ {{CF}} $ 所成的角为 $ {60}^{ \circ  } $ .

\begin{center}

    \includegraphics[width=0.25\linewidth]{bodyxb//src/aksufha'siuhfasd.png}
\end{center}

\end{enumerate}



\item 如图所示，在四棱锥 $ P - {ABCD} $ 中，底面 $ {ABCD} $ 为梯形， $ {BC}//{AD} $ ， $ {BC} = \dfrac{1}{2}{AD} $ ，面 $ {PBC} \cap $ 面 $ {PAD} = l $ ， $ E $ 是 $ {PD} $ 的中点.
\begin{enumerate}
    \item 求证： $ {CE}// $ 平面 $ {PAB} $ ；
    \item 求证： $ l//{AD} $ ；
    \item 若 $ M $ 是线段 $ {CE} $ 上一动点，则线段 $ {AD} $ 上是否存在点 $ N $ ， 使 $ {MN}// $ 平面 $ {PAB} $ ? 说明理由.   
\end{enumerate}

\begin{center}

    \includegraphics[width=0.25\linewidth]{bodyxb//src/a'osuhf'aiuhf'asf.png}
\end{center}

\begin{jieda}
    \begin{enumerate}
        \item 在四棱锥 $ P - {ABCD} $ 中,取 $ {PA} $ 中点 $ F $ ,连 $ {BF} $ 、 $ {EF} $
$ \because {EF} = \dfrac{1}{2}{AD},{EF} ⫫ {AD} $
又 $ \because {BC} = \dfrac{1}{2}{AD},{BC} ⫫ {AD} $
$ \therefore {EF} = {BC},{EF}// {BC} $
$ \therefore $ 四边形 $ {BCEF} $ 为平行四边形,
$ \therefore {CE} ⫫ {BF} $ ,
又 $ \because {CE} \not \subset$ 平面 $ {PAB},{BF} \subset $ 平面 $ {PAB},\therefore {CE} \cup $ 平面 $ {PAB} $ ;

\begin{center}

    \includegraphics[width=0.25\linewidth]{bodyxb//src/asufgh'aiuhf'aiuhsfd.png}
\end{center}

        \item $ \because $ 在梯形 $ {ABCD} $ 中, $ {BC}\square {AD} $ ,
又 $ {BC} \not \subset$ 面 $ {PAD},{AD} \subset $ 面 $ {PAD} $ ,
$ \therefore {BC} // $ 面 $ {PAD} $ ,
$ \because {BC} \subset $ 面 $ {PBC} $ ,面 $ {PBC} \cap $ 面 $ {PAD} = l $
$ \therefore {BC}// l $
$ \because {BC}// {AD}\therefore {l}^{\prime \prime }{AD} $
        
        \item 取 $ {AD} $ 中点 $ N $ ,连接 $ {CN} $ , $ {EN} $ ,
$ \because E,N $ 分别为 $ {PD},{AD} $ 的中点,
$ \therefore {EN}// {PA} $ ,
$ \because {EN} \not \subset$ 平面 $ {PAB},{PA} \subset $ 平面 $ {PAB} $ ,
$ \therefore {EN} // $ 平面 $ {PAB} $ ,
又由 (1) 可得 $ {CE}// $ 平面 $ {PAB},{CE} \cap  {EN} = E,{CE}\text{、}{EN} \subset $ 平面 $ {CEN} $
$ \therefore $ 平面 $ {CEN} ⫫ $ 平面 $ {PAB} $ ,
$ \because M $ 是 $ {CE} $ 上的动点, $ {MN} \subset $ 平面 $ {CEN} $ ,
$ \therefore {MN} ⫫ $ 平面 $ {PAB} $ ,
$ \therefore $ 当 $ N $ 为中点时, $ {MN} // $ 平面 $ {PAB} $ .
    \end{enumerate}
\end{jieda}

\item 如图1，等腰梯形 $ {ABCD} $ 中， $ {AD}//{BC} $ ， $ E $ 为 $ {AD} $ 边上一点，且 $ {AB} = {AE} = {BE} = {CD} = 2 $ ， $ {BC} = {ED} = 4,O $ 为 $ {BE} $ 中点， $ F $ 为 $ {BC} $ 中点. 将 $ \bigtriangleup  {ABE} $ 沿 $ {BE} $ 折起到 $ \bigtriangleup  {{A}^{\prime }{BE}} $ 的位置，如图 2 .

\begin{center}

    \includegraphics[width=0.25\linewidth]{bodyxb//src/aiushf'aih'fudafsd.png}
\end{center}

\begin{enumerate}
    \item 证明： $ {CD} \bot $ 平面 $ {A}^{\prime }{OF} $ ；
    \item 若 $ {A}^{\prime }O \bot  {OF} $ ，求点 $ F $ 到平面 $ {A}^{\prime }{EC} $ 的距离；
    \item 试在线段 $ {A}^{\prime }D $ 上确定一点 $ G $ ,使得平面 $ {GCE}// $ 平面 $ {A}^{\prime }{OF} $ ,并给出证明.
\end{enumerate}
    \begin{jieda}
        \begin{enumerate}
            \item 在等腰梯形 $ {ABCD} $ 中, $ {AD}//{BC} $ , $ {BC} = {ED} = 4 $ ,
则四边形 $ {BCDE} $ 是平行四边形,则 $ {BE}//{CD} $ ,
因为 $ {AB} = {AE} = {BE} = {CD} = 2 $ ,所以 $ \Delta {ABE} $ 为等边三角形,则 $ \angle {BAE} = {60}^{ \circ  } $
因为 $ O $ 为 $ {BE} $ 中点,所以 $ {BE}\bot {AO} $ ,
在等腰梯形 $ {ABCD} $ 中,可得 $ \angle {BAE} = \angle {CDE} = \angle {CBE} = {60}^{ \circ  } $
连接 $ {EC} $ ,在 $ {\Delta BEC} $ 中,由余弦定理可 $ E{C}^{2} = B{C}^{2} + B{E}^{2} - {2BC} \cdot  {BE} \cdot  \cos \angle {EBC} = {12} $ ,
则 $ {EC} = 2\sqrt{3} $ ,所以 $ B{E}^{2} + E{C}^{2} = {16} = B{C}^{2} $ ,则 $ {BE}\bot {EC} $ .
因为 $ O $ 、 $ F $ 分别是 $ {BE} $ 、 $ {BC} $ 中点,所以 $ {OF}//{EC} $ ,所以 $ {OF}\bot {BE} $ ,
从而可得 $ {OF} \bot  {CD},{CD} \bot  {AO} $ ,
因为 $ \dot{AO} \cap  {OF} = O,\dot{AO}\text{、}{OF} \subset $ 平面 $ \dot{AOF} $ ,所以 $ {CD} \bot $ 平面 $ \dot{AOF} $ ；
    \item 由(1)可知 $ {OF}//{EC} $ ,因 $ {OF} \not \subset$ 平面 $ {AEC} $ , $ {EC} \subset $ 平面 $ {AEC} $ ,则 $ {OF}// $ 平面 $ {AEC} $ ,
所以点 $ F $ 到平面 $ A\dot{}{EC} $ 的距离即为点 $ O $ 到平面 $ A\dot{}{EC} $ 的距离,
因为 $ O $ 是 $ {BE} $ 中点,所以点 $ F $ 到平面 $ {AEC} $ 的距离等于点 $ B $ 到平面 $ {AEC} $ 的距离的一半.
取 $ {A}^{\prime }E $ 的中点为 $ H $ ,连接 $ {BH} $ ,
因为 $ \Delta {A}^{\prime }{BE} $ 为等边三角形,所以 $ {BH} \bot  {A}^{\prime }E $ ,

\begin{center}

    \includegraphics[width=0.25\linewidth]{bodyxb//src/asfuha'iuhf'aushdfa.png}
\end{center}


由(1)知 $ {EC}\bot {BE} $ ,因为 $ {A}^{\prime }O\bot {OF} $ ,而 $ {OF}//{EC} $ ,则 $ {EC}\bot {A}^{\prime }O $ ,
又因为 $ {BE} \cap  {AD} = O,{BE}\text{、}{AD} \subset $ 平面 $ {ABE} $ ,所以 $ {EC}\bot $ 平面 $ {ABE} $ ,
因为 $ {BH} \subset $ 平面 $ {A}^{\prime }{BE} $ ,所以 $ {EC} \bot  {BH} $ ,
又因为 $ {EC} \cap  {AE} = E,{EC}\text{、}{AE} \subset $ 平面 $ {AEC} $ ,所以 $ {BH}\bot $ 平面 $ {AEC} $ ,
因为 $ \Delta {A}^{\prime }{BE} $ 是等边三角形,边长为 2,故 $ {BH} = \sqrt{{2}^{2} - {1}^{2}} = \sqrt{3} $ ,
所以点 $ B $ 到平面 $ {A}^{\prime }{EC} $ 的距离为 $ \sqrt{3} $ ,
故点 $ F $ 到平面 $ {AEC} $ 的距离为 $ \dfrac{\sqrt{3}}{2} $ ;
    \item 设 $ {DF} \cap  {CE} = M $ ,则 $ \dfrac{FM}{MD} = \dfrac{CF}{ED} = \dfrac{1}{2} $ ,
在 $ {AD} $ 上取点 $ G $ ,使 $ \dfrac{{A}^{\prime }G}{GD} = \dfrac{1}{2} $ ,从而 $ {GM}//{A}^{\prime }F $ ,
连接 $ {EG} $ ,因为 $ \dot{AF} \not \subset$ 平面 $ {GEC},{GM} \subset $ 平面 $ {GEC} $ ,所以 $ \dot{AF}// $ 平面 $ {GEC} $ ,

\begin{center}

    \includegraphics[width=0.25\linewidth]{bodyxb//src/ufha'shf'aosuhfa'oud.png}
\end{center}
   又由 (1) 可知, $ {OF}//{EC},{OF} \not \subset$ 平面 $ {GEC},{GM} \subset $ 平面 $ {GEC} $ ,
则 $ {OF}// $ 平面 $ {GEC} $ ,
又因为 $ {AF} \cap  {OF} = F,{AF},{OF} \subset $ 平面 $ {AOF} $ ,所以平面 $ {GCE}// $ 平面 $ {AOF} $ ,
故 $ G $ 为线段 $ {AD} $ 上靠近 $ A $ 的三等分点时,平面 $ {GCE}// $ 平面 $ {AOF} $ .

        \end{enumerate}
    \end{jieda}



    
\end{enumerate}



\subsection*{平面与平面的位置关系（2）}


A 组：

\begin{enumerate}


\item 二面角的取值范围是\blkda{$[0,\pi]$}.

\item 二面角是指\dotfill{(\kh{D})}

\xx{两个平面相交所组成的图形}
{一个平面绕这个平面内一条直线旋转所组成的图形}
{从一个平面内的一条直线出发的一个半平面与这个平面所组成的图形}
{从一条直线出发的两个半平面所组成的图形}

\item 如图，已知 $ {PA} \bot  a $ ， $ {PB} \bot  \beta $ ，垂足为 $ {A} $ 、 $ B $ ，若 $ \angle {APB} = {60}^{ \circ  } $ ，则二面角 $ a - l - \beta $ 的大小是\blkda{$\dfrac{2\pi}{3}$}.

\begin{center}

    \includegraphics[width=0.25\linewidth]{bodyxb//src/a'oushfaouhf'shdfas.png}
\end{center}

    
\end{enumerate}

B 组:

\begin{enumerate}


\item 若分别向小于 90° 的二面角的两个半平面作垂线，且它们是异面直线，则这两条直线所成的角与二面角的平面角的关系是 \dotfill{(\kh{A})}

\xx{相等} {互补} {互余} {既不相等又不互补也不互余}



\item 已知三棱锥 $ D - {ABC} $ 的三个侧面与底面全等，且 $ {AB} = {AC} = \sqrt{3},{BC} = 2 $ ，则二面角 $ D - {BC} - A $ 的余弦值为\blkda{0}.


\item 在正三角形 $ {ABC} $ 中， $ {AD}\bot {BC} $ 于 $ D $ ，沿 $ {AD} $ 折成二面角 $ B $ — $ {AD} - C $ 后， $ {BC} = \dfrac{1}{2}{AB} $ ， 这时二面角 B-AD-C 的大小为\blkda{$\dfrac{\pi}{3}$}.


\item 已知 $ P $ 为锐二面角内一点，且 $ P $ 到两个半平面及棱的距离之比为 $ 1 : \sqrt{2} : 2 $ ，则此二面角的度数为\blkda{$75^{\circ}$}.


\item 在 ABCD 为正方形，E 为 AB 中点，将 $ \bigtriangleup  {DAE} $ 和 $ \bigtriangleup  {CBE} $ 分别沿 DE、CE 折起， 使 A 与 B 重合, 重合后的点记为 P, 则平面 PCD 与平面 ECD 所成二面角的大小为\blkda{$\dfrac{\pi}{6}$}.



\item 已知菱形的一个内角是 $ {60}^{ \circ  } $ ，边长为 $ a $ ，沿菱形较短的对角线折成大小为 $ {60}^{ \circ  } $ 的二面角，则菱形中含 $ {60}^{ \circ  } $ 角的两个顶点间的距离为
\blkda{$\dfrac{\sqrt{3}a}{2}$}.


\item 在正方体 $ {{ABCD}} - {{A}}_{1}{{B}}_{1}{{C}}_{1}{{D}}_{1} $ 中，平面 $ {{A}}_{1}{{{DC}}}_{1} $ 和平面 $ {{A}}_{1}{{{BC}}}_{1} $ 所成的二面角 (锐角) 大小为\blkda{$arccos\left(\dfrac{1}{3}\right)$}.




\item 已知二面角 $ a - {AB} - \beta $ 的平面角是锐角 $ \theta $ ， $ a $ 内一点C到 $ \beta $ 的距离为3，点C 到棱AB的距离为4，则 $ \tan \theta  = $ \blkda{$\dfrac{3\sqrt{7}}{7}$}.



\item 在 $30^\circ$ 的二面角的一个面内有一个点，若它到另一个面的距离是 10，则它到棱的距离为\blkda{20}.



\item 若二面角内一点到二面角的两个面的距离分别为 $ a $ 和 $ \sqrt{2}a $ ，到棱的距离为 $ {2a} $ ，则此二面角的度数是\blkda{$75^{\circ}$或$165^{\circ}$}.



\item 在棱长为 $a$ 的正方体 $ {{ABCD}} - {{A}}_{1}{{B}}_{1}{{C}}_{1}{{D}}_{1} $ 中，求二面角 $ {{C}}_{1} - {{BD}} - {C} $ 的大小.

\begin{jieda}
    按照定义法，$\angle{C_1DC}$即为二面角的平面角，因此$arctan\sqrt{2}$.
\end{jieda}

\item 过正方形 $ {ABCD} $ 的顶点 $ A $ 作 $ {PA} \bot $ 平面 $ {ABCD} $ ，设 $ {PA} = {AB} = a $ ，求二面角 $ B - PC-D$ 的大小.

\begin{jieda}
    连接 $ {AC},{BD} $ ,因为 $ {PA} \bot $ 平面 $ {ABD},{BD} \subset $ 平面
$ {ABD} $ ,所以 $ {PA} \bot  {BD} $ ,又 $ {BD} \bot  {AC},{AC},{PA} \subset $ 平面 $ {PAC} $ ,
$ {PA} \cap  {AC} = A $ ,所以 $ {BD}\bot $ 平面 $ {PAC} $ ,又 $ {PC} \subset $ 平面 $ {PAC} $ ,
所以 $ {BD} \bot  {PC} $ ,过 $ {BD} $ 作平面 $ {BDH} \bot  {PC} $ 于 $ H $ ,则 $ {PC} \bot $
$ {DH},{PC} \bot  {BH} $ ,所以 $ \angle {BHD} $ 为二面角 $ B - {PC} - D $ 的平
面角. 由题知 $ {PB} = \sqrt{2}a,{BC} = a $ ,所以 $ {PC} = \sqrt{3}a,\dfrac{1}{2}{PB} $
$ \cdot  {BC} = {S}_{\bigtriangleup {PBC}} = \dfrac{1}{2}{PC} \cdot  {BH} $ ,所以 $ {BH} = \dfrac{\sqrt{6}}{3}a = {DH} $ ,又
$ {BD} = \sqrt{2}a $ ,在 $ \bigtriangleup {BHD} $ 中,由余弦定理得 $ \cos \angle {BHD} = $
$ \dfrac{B{H}^{2} + D{H}^{2} - B{D}^{2}}{{2BH} \cdot  {DH}} = \dfrac{{\left( \dfrac{\sqrt{6}}{3}a\right) }^{2} + {\left( \dfrac{\sqrt{6}}{3}a\right) }^{2} - {\left( \sqrt{2}a\right) }^{2}}{2 \times  \dfrac{\sqrt{6}}{3}a \times  \dfrac{\sqrt{6}}{3}a} = $
$ - \dfrac{1}{2} $ ,又 $ 0 < \angle {BHD} < \pi $ ,所以 $ \angle {BHD} = \dfrac{2\pi }{3} $ ,故二面角
$ B - {PC} - D $ 的大小是 $ \dfrac{2\pi }{3} $ .

\end{jieda}



\item 如图，已知 $ \operatorname{Rt} \bigtriangleup  {ABC} $ 的直角边 $ {AB} = 4 $ ， $ {AC} = 2 $ ， $ {PA}\bot $ 平面 $ {ABC} $ ， $ {PA} = 2 $ ，求二面角 $ P - {BC} - A $ 的大小.

\begin{center}

    \includegraphics[width=0.25\linewidth]{bodyxb//src/asuhf'aouhfasfa.png}
\end{center}

\begin{jieda}
如图,过点 $ A $ 作 $ {AD}\bot {BC} $ 于 $ D $ ,连接 $ {PD} $ .

  \begin{center}
  
      \includegraphics[width=0.25\linewidth]{bodyxb//src/hf'ahf'ahf;haosfhasfa.png}
  \end{center}
  $ \because {PA} \bot $ 平面 $ {ABC},\therefore {PA} \bot  {BC} $ ,
又 $ {AD} \bot  {BC},\therefore {BC} \bot $ 平面 $ {PAD},\therefore {BC} \bot  {PD} $ ,
$ \therefore \angle {PDA} $ 是二面角 $ P - {BC} - A $ 的平面角.
$ \because {AB} = 4,{AC} = 2,\therefore {BC} = \sqrt{{4}^{2} + {2}^{2}} = 2\sqrt{5},{AD} = \dfrac{4 \times  2}{2\sqrt{5}} = \dfrac{4\sqrt{5}}{5} $ ,
$ \therefore $ 在 Rt $ \bigtriangleup  {PAD} $ 中, $ \tan \angle {PDA} = \dfrac{PA}{AD} = \dfrac{\sqrt{5}}{2} $ ,
$ \therefore \angle {PDA} = \arctan \dfrac{\sqrt{5}}{2} $ ,即二面角 $ P - {BC} - A $ 的大小为 $ \arctan \dfrac{\sqrt{5}}{2} $ .

\end{jieda}



\item 已知直角三角形 $ {ABC} $ 的斜边 $ {AB} $ 在平面 $ a $ 内， $ {AC} $ 、 $ {BC} $ 与平面 $ a $ 所成的角分别为 $ {30}^{ \circ  } $ 和 $ {45}^{ \circ  } $ ,求 $ \bigtriangleup {ABC} $ 所在平面与 $ a $ 所成的二面角 (锐角) 的大小.

\begin{jieda}
    过点 $ C $ 作 $ {CO}\bot \alpha $ ,交 $ \alpha $ 于 $ O $ ,连结 $ {AO},{BO} $ ,作 $ {CD}\bot l $ ,交 $ l $ 于点 $ D $ ,连结 $ {OD} $ ,
则 $ \angle {CDO} $ 是二面角 $ \alpha  - l - \beta $ 的平面角,
$ \because {{Rt}}\angle {{ABC}} $ 在平面 $ \beta $ 内, $ \angle {ACB} = {90}^{ \circ  } $ ,斜边 $ {AB} $ 在二面角 $ \alpha  - l - \beta $ 的棱 $ l $ 上,
$ {AC} $ 与平面 $ \alpha $ 所成角为45 $ {}^{ \circ  } $ , $ {BC} $ 与平面 $ \alpha $ 所成角为30 $ {}^{ \circ  } $ ,
$ \therefore \angle {CAO} = {45}^{ \circ  },\angle {CBO} = {30}^{ \circ  } $ ,
设 $ {CO} = 1 $ ,则 $ {AO} = 1 $ , $ {AC} = \sqrt{2} $ , $ {BC} = 2 $ , $ {BO} = \sqrt{3} $ , $ {AB} = \sqrt{6} $ ,
$ {CD} = \dfrac{{AC} \times  {BC}}{AB} = \dfrac{\sqrt{2} \times  2}{\sqrt{6}} = \dfrac{2}{\sqrt{3}},{AD} = \sqrt{2 - \dfrac{4}{3}} = \sqrt{\dfrac{2}{3}},{OD} = \sqrt{1 - \dfrac{2}{3}} = \sqrt{\dfrac{1}{3}}, $
$ \therefore \cos \angle {CDO} = \dfrac{C{D}^{2} + O{D}^{2} - C{O}^{2}}{2 \times  {CD} \times  {OD}} = \dfrac{\dfrac{4}{3} + \dfrac{1}{3} - 1}{2 \times  \sqrt{\dfrac{4}{3}} \times  \sqrt{\dfrac{1}{3}}} = \dfrac{1}{2} $ ,
$ \therefore \angle {CDO} = {60}^{ \circ  } $ ,
$ \therefore $ 二面角 $ \alpha  - l - \beta $ 的平面角大小为 $ {60}^{ \circ  } $ .

\begin{center}

    \includegraphics[width=0.25\linewidth]{bodyxb//src/iisgjijgoiaejgofaerfa.png}
\end{center}

\end{jieda}


\item 如图，四棱锥 $ P - {ABCD} $ 的底面是边长为 $ a $ 的菱形， $ \angle {BCD} = {120}^{ \circ  } $ ， $ {PC} \bot $ 平面 $ {ABCD},{PC} = a,E $ 为 $ {PA} $ 的中点， $ O $ 为底面对角线的交点；

\begin{enumerate}
    \item 求证： $ {AC} \bot $ 平面 $ {EDB} $ ；
    \item 求二面角 $ P - {AB} - C $ 的正切值.
\end{enumerate}

\begin{center}

    \includegraphics[width=0.25\linewidth]{bodyxb//src/fa'jshf'ahsfasdf.png}
\end{center}


\begin{jieda}
    \begin{enumerate}
        \item 如图,连接 $ {EO} $ , $ \because E $ 为 $ {PA} $ 的中点, $ O $ 为菱形对角线的交点, $ \therefore O $ 为 $ {AC} $ 中点, $ \therefore {EO} //  {PC} $ ,又 $ \because {PC} \bot $ 平面 $ {ABCD} $
$ \therefore {PC} \bot  {AC} $ , $ \therefore {EO} \bot  {AC} $ ,又 $ \because $ 菱形的对角线互相垂直平分, $ \therefore {AC} \bot  {BD} $ ,又 $ {EO} \cap  {BD} = O $ , $ \therefore {AC} \bot $ 平面 $ {EDB} $ ；

\begin{center}

    \includegraphics[width=0.25\linewidth]{bodyxb//src/dsghjsoijgaoiefaf.png}
\end{center}

        \item 如图,设 $ {AB} $ 中点为 $ F $ ,由底面 $ {ABCD} $ 为菱形可知, $ \bigtriangleup  {ABC} $ 为等边三角形, $ \therefore {CF} \bot  {AB} $ ,
又 $ \because {PC} \bot $ 平面 $ {ABCD} $ , $ \therefore {PC} \bot  {AB} $ , $ \therefore {AB} \bot $ 平面 $ {PCF} $ , $ \therefore {AB} \bot  {PF} $ , $ \therefore \angle {PFC} $ 为二面角 $ P - {AB} - C $ 的平面角, $ {CF} = \dfrac{\sqrt{3}}{2}a - {PC} = a $
,故 $ \tan \angle {PFC} = \dfrac{PC}{CF} = \dfrac{2\sqrt{3}}{3} $.

\begin{center}

    \includegraphics[width=0.25\linewidth]{bodyxb//src/gijsdghsodhgaouhgf.png}
\end{center}

    \end{enumerate}
\end{jieda}

\end{enumerate}



C 组:

\begin{enumerate}

\item 如图，在正方体 $ \mathbf{{ABCD} - {A}_{1}{B}_{1}{C}_{1}{D}_{1}} $ 中，动点 $ \mathbf{M} $ 在线段 $ \mathbf{{A}_{1}C} $ 上，异面直线 $ \mathbf{A}{\mathbf{D}}_{\mathbf{1}} $ 和 $ \mathbf{B}\mathbf{M} $ 所成的角为 $ \mathbf{\theta } $ ，则 $ \mathbf{\theta } $ 的取值范围是\blkda{$[\dfrac{\pi}{6},\dfrac{\pi}{3}]$}.（用区间表示）

\begin{center}

    \includegraphics[width=0.25\linewidth]{bodyxb//src/a'sfh'aouhfadfasfsfa.png}
\end{center}



\item 已知正方体 $ {ABCD} - {A}_{1}{B}_{1}{C}_{1}{D}_{1} $ 的棱长为 $ 4,M $ 为棱 $ {DC} $ 的中点， $ N $ 为侧面 $ B{C}_{1} $ 的中心，过点 $ M $ 的平面 $ \alpha $ 垂直于 $ {DN} $ ，则平面 $ \alpha $ 截正方体 $ A{C}_{1} $ 所得的截面面积为\blkda{$4\sqrt{6}$}.




\end{enumerate}





\subsection*{平面与平面的位置关系（3）}


A 组:

\begin{enumerate}




\item 下列命题正确的是\dotfill{(\kh{B})}

\xx{直线 $ a// $ 平面 $ a $ ,直线 $ b\bot a,b \subset $ 平面 $ \beta $ ,则 $ a \bot  \beta $}
{若直线 $ a \bot $ 直线 $ b,a \bot $ 平面 $ a,b \bot $ 平面 $ \beta $ ，则 $ a \bot  \beta $}
{过平面外的一条直线有且只有一个平面与已知平面垂直}
{过平面外一点有且只有一个平面与已知平面垂直}

\item 若平面 $ \alpha $ 、 $ \beta $ 互相垂直，则 \dotfill{(\kh{D})}

\xx{$ \alpha $ 中的任意一条直线垂直于 $ \beta $} 
{$ \alpha $ 中有且只有一条直线垂直于 $ \beta $}
{平行于 $ \alpha $ 的直线垂直于 $ \beta $} 
{$ \alpha $ 内垂直于交线的直线必垂直于 $ \beta $}


\item 已知三个平面 $ a $ 、 $ \beta $ 、 $ \gamma $ ，则下列命题中正确的个数是 \dotfill{(\kh{A})}
\begin{enumerate}

\item 若 $ a \bot  \beta ,\gamma  \bot  \beta $ ，则 $ a \bot  \gamma $ ；
\item 若 $ a \bot  \beta ,\gamma  \bot  \beta , $ ，则 $ a//\gamma $ ；
\item 若 $ a//\beta ,\beta  \bot  \gamma $ ，则 $ a \bot  \gamma $ ；
    
\end{enumerate}
\xx{1}{2}{3}{0} 

\item 已知两个平面互相垂直，那么下列命题中正确的个数是 \dotfill{(\kh{C})}
\begin{enumerate}

\item 一个平面内的直线必垂直于另一个平面内的无数条直线；

\item 一个平面内且垂直于这两个平面交线的直线必垂直于另一个平面内的任意一条直线;

\item 一个平面内的任何一条直线必垂直于另一个平面；

\item 过一个平面内的任意一点作交线的垂线，则此直线必垂直于另一个平面。
\end{enumerate}
\xx{4}{3}{2}{1}


\item 在空间两两垂直的平面最多有\dotfill{(\kh{B})}

\xx{两个} {三个} {四个} {无数个}



\end{enumerate}


B 组:

\begin{enumerate}


\item 在空间四边形 ABCD 中，已知二面角 A一BCD,A一CD=B,A一BD一C 都相等，则 A 点在平面 BCD 上的射影是 $ \bigtriangleup  {BCD} $ 的\dotfill{(\kh{A})}

\xx{内心}{外心} {垂心} {重心}



\item 若一个二面角的两个半平面分别垂直与另一个二面角的两个半平面，则这两个二面角的关系是\dotfill{(\kh{D})}

\xx{相等} {互补} {相等或互补} {不能确定}

\item 从二面角 P-MN—Q 内一点 A 分别作 AB⊥平面 P，AC⊥平面 Q，其中 B、C 为垂足， $ \angle {BAC} = {60}^{ \circ  } $ ，那么二面角 P-MN—Q 的大小是\blkda{$\dfrac{2\pi}{3}$}.

\item 已知 $ {A} $ 、 $ {B} $ 是二面角 $ a - l - \beta $ 的棱 $ l $ 上两点，P 是面 $ \beta $ 内一点，PB $ \bot  l $ 于点 B， PA 和 $ l $ 所成的角为 $ {45}^{ \circ  } $ ，PA 和面 $ a $ 所成的角为 $ {30}^{ \circ  } $ ，则二面角 $ a - l - \beta $ 的大小为\blkda{$\dfrac{\pi}{4}$或$\dfrac{3\pi}{4}$}.

\item 已知点 $ {A} $ 、 $ {B} $ 分别是二面角 $ {A} - {{CD}} - {B} $ 的各面上的一点， $ {{AC}} \bot  {{CD}} $ ， $ {{BD}} \bot  {{CD}} $ ，C、D 分别为垂足，若 CD=3,AC= $ \sqrt{2} $ ,BD=2,AB= $ \sqrt{10} $ ，则二面角 A—CD—B 的度数是\blkda{$arccos\dfrac{5\sqrt{2}}{8}$}.



\item 如图，矩形 $ABEF$ 和正方形 $ABCD$ 有公共边 $AB$，它们所在的半平面成 $ {60}^{ \circ  } $ 的二面角,且 $ {{AB}} = {2a},{{BF}} = a $ ,则 $ {{FC}} = $\blkda{$\sqrt{7}a$}


\begin{center}

    \includegraphics[width=0.25\linewidth]{bodyxb//src/f'aohf'ajhfasfa.png}
\end{center}



\item 在正方体 $ {{ABCD}} - {{A}}_{1}{{B}}_{1}{{C}}_{1}{{D}}_{1} $ 中，平面 $ {{{AA}}}_{1}{{C}}_{1}{C} $ 和平面 $ {{A}}_{1}{{{BCD}}}_{1} $ 所成的二面角（锐角）大小为\blkda{$\dfrac{\pi}{3}$}.


\item 在矩形 $ {ABCD} $ 中， $ {AB} = 3,{AD} = 4,{PA} \bot $ 平面 $ {ABCD},{PA} = \sqrt{3} $ ，则平面 $ {PBD} $ 与平面 $ {ABCD} $ 的夹角的正切值为\blkda{$\dfrac{5\sqrt{3}}{12}$}.


\item 如图，在空间四边形 $ {ABCD} $ 中， $ \bigtriangleup  {BCD} $ 是正三角形， $ \bigtriangleup {ABD} $ 是等腰直角三角形，且 $ \angle {BAD} = {90}^{ \circ  } $ ，又二面角 $ A - {BD} - C $ 为直二面角,则二面角 $ A - {CD} - B $ 的正切值为\blkda{$\dfrac{2\sqrt{3}}{3}$}.
\begin{center}

    \includegraphics[width=0.25\linewidth]{bodyxb//src/hf'ihfuu'shdfui'sdfh.png}
\end{center}



\item 在二面角 $ \alpha  - {AB} - \beta $ 的一个平面 $ \alpha $ 内，有一直线 $ {AC} $ ，它与棱 $ {AB} $ 成 $ {45}^{ \circ  } $ 角， $ {AC} $ 与平面 $ \beta $ 成 $ {30}^{ \circ  } $ 角,二面角 $ \alpha  - {AB} - \beta $ 的大小为\blkda{$\dfrac{\pi}{4}$或$\dfrac{3\pi}{4}$}





\item 已知 $ {P} $ 为 $ {120}^{ \circ  } $ 的二面角 $ a - a - \beta $ 内一点，且点 $ {P} $ 到平面 $ a $ 和 $ \beta $ 的距离分别为 6 和8，求点 $ {P} $ 到棱 $ a $ 的距离.
\begin{jieda}
    设长度为x，过P作$\alpha$和$\beta$的垂线，交于A、B两点，过P作棱的垂线，交于M，则有$sin\angle{PAM}=\dfrac{6}{x}$,$sin\angle{PBM}=sin(120^{\circ}-\angle{PAM})=\dfrac{8}{x}$,从而求得$x=\dfrac{4\sqrt{39}}{3}$.
\end{jieda}



\item 如图，在三棱锥 $ P - {ABC} $ 中， $ {AB}\bot {BC} $ ，

$ {AB} = {PB} = {PC} = 2,{PA} = {BC} = 2\sqrt{2} $ ，棱 $ {PB},{PA},{AC},{BC} $ 的中点分别为 $ D,E,F,O $ .

\begin{enumerate}
    \item 证明： $ {EF}// $ 平面 $ {ADO} $ ；
    \item 证明：平面 $ {BEF} \bot $ 平面 $ {PAO} $ .
\end{enumerate}

\begin{center}

    \includegraphics[width=0.25\linewidth]{bodyxb//src/foushf'oushf'ouhsdfsdfs.png}
\end{center}

\begin{jieda}
    \begin{enumerate}
        \item 因为棱 $ {PB},{PA},{AC},{BC} $ 的中点分别为 $ D,E,F,O $ ,
所以 $ {EF}//{PC},{DO}//{PC} $ ,所以 $ {EF}//{DO} $ ,
又因为 $ {DO} \subset $ 平面 $ {ADO} $ , $ {EF} \not \subset$ 平面 $ {ADO} $ ,所以 $ {EF}// $ 平面 $ {ADO} $ .
        \item 因为 $ {PB} = {PC} = 2,{BC} = 2\sqrt{2} $ , $ O $ 是 $ {BC} $ 的中点,
所以 $ {PO}\bot {BC} $ ,且 $ {PO} = \sqrt{2} $ .
因为 $ {AB}\bot {BC} $ ,所以 $ {AO} = \sqrt{A{B}^{2} + B{O}^{2}} = \sqrt{6} $ ,
所以 $ P{O}^{2} + A{O}^{2} = P{A}^{2} $ ,所以 $ {PO}\bot {AO} $ .
又因为 $ {AO} \cap  {BC} = O,{AO},{BC} \subset $ 平面 $ {ABC} $ ,所以 $ {PO}\bot $ 平面 $ {ABC} $ .
又 $ {BF} \subset $ 平面 $ {ABC} $ ,所以 $ {PO} \bot  {BF} $ .
如图,连接 $ {OF} $ ,则 $ {OF} //  {AB} $ ,从而 $ {OF}\bot {BC} $ ,且 $ {OF} = \dfrac{1}{2}{AB} = 1 $ ,所以 $ \dfrac{OB}{AB} = \dfrac{OF}{OB} = \dfrac{\sqrt{2}}{2} $ ,得 $ \operatorname{Rt} \bigtriangleup  {ABO} \backsim  \operatorname{Rt} \bigtriangleup  {BOF} $ ,
所以 $ \angle {FBO} + \angle {AOB} = {90}^{ \circ  } $ ,所以 $ {AO}\bot {BF} $ .
由 $ {PO}\bot {BF},{AO}\bot {BF},{AO} \cap  {PO} = O,{AO},{PO} \subset $ 平面 $ {POA} $ ,得 $ {BF}\bot $ 平面 $ {POA} $ ,
又 $ {BF} \subset $ 平面 $ {BEF} $ ,所以平面 $ {BEF} \bot $ 平面 $ {POA} $ .

\begin{center}

    \includegraphics[width=0.25\linewidth]{bodyxb//src/fhs'ofhaoighaouhfaohfa.png}
\end{center}

    \end{enumerate}
\end{jieda}




\item 如图，在四棱锥 $ M - {ABCD} $ 中，底面 $ {ABCD} $ 是边长为 2 的菱形， $ {MA} \bot $ 平面 $ {ABCD} $ ，平面 $ {MAB} \cap $ 平面 $ {MCD} = l $ ，平面 $ {MAB} \bot $ 平面 $ {MBC} $ .
\begin{enumerate}

\item 证明： $ l //{CD} $ ；

\item 证明： $ {AD} \bot $ 平面 $ {MAB} $ ；

\item 当 $ {MA} $ 为何值时，二面角 $ B - {MC} - D $ 的余弦值为 $ - \dfrac{1}{3} $ .
    
\end{enumerate}

\begin{center}

    \includegraphics[width=0.25\linewidth]{bodyxb//src/glahgd'ush'fouhsdgfa.png}
\end{center}

\begin{jieda}
    \begin{enumerate}
        \item 菱形 $ {ABCD} $ 中, $ {CD}//{AB},{AB} \subset $ 平面 $ {MAB},{CD} \not \subset$ 平面 $ {MAB} $ ,
则 $ {CD}// $ 平面 $ {MAB} $ ,
而平面 $ {MAB} \cap $ 平面 $ {MCD} = l,{CD} \subset $ 平面 $ {MCD} $ ,
所以 $ l//{CD} $ .

        \item 在平面 $ {MAB} $ 内过点 $ {A作AE}\bot {MB} $ 于 $ E $ ,平面 $ {MAB}\bot $ 平面 $ {MBC} $ ,平面 $ {MAB} \cap $ 平面 $ {MBC} = {MB} $ ,
则 $ {AE} \bot $ 平面 $ {MBC} $ ,而 $ {BC} \subset $ 平面 $ {MBC} $ ,于是 $ {AE} \bot  {BC} $ ,
又 $ {MA}\bot $ 平面 $ {ABCD},{BC} \subset $ 平面 $ {ABCD} $ ,
则 $ {MA}\bot {BC} $ ,而 $ {MA} \cap  {AE} = A,{MA} \subset $ 平面 $ {MAB} $ ,
因此 $ {BC}\bot $ 平面 $ {MAB} $ ,又 $ {BC}//{AD} $ ,
所以 $ {AD} \bot $ 平面 $ {MAB} $ .

\begin{center}

    \includegraphics[width=0.25\linewidth]{bodyxb//src/gsrhg'uahgifa'fhaw.png}
\end{center}

        \item 由(2)知 $ {AD} \bot $ 平面 $ {MAB} $ , $ {AB} \subset $ 平面 $ {MAB} $ ,则 $ {AB} \bot  {AD} $ ,菱形 $ {ABCD} $ 为正方形,
由 $ {BC} \bot $ 平面 $ {MAB},{MB} \subset $ 平面 $ {MAB} $ ,得 $ {BC} \bot  {MB} $ ,
过 $ B $ 作 $ {BO}\bot {MC} $ 于 $ O $ ,连接 $ {BD},{DO} $ ,
$ {MB} = \sqrt{A{B}^{2} + M{A}^{2}} = \sqrt{A{D}^{2} + M{A}^{2}} = {MD} $ ,而 $ {BC} = {DC},{MC} = {MC} $ ,
则 $ {\Delta MBC} \cong  {\Delta MDC} $ ,有 $ \angle {BMO} = \angle {DMO} $ ,于是 $ {\Delta BMO} \cong  {\Delta DMO} $ ,则 $ \angle {DOM} = \angle {BOM} = {90}^{ \circ  } $ ,
即 $ {DO}\bot {MC},\angle {BOD} $ 是二面角 $ B - {MC} - D $ 的平面角,
令 $ {MA} = t,{MC} = \sqrt{M{B}^{2} + B{C}^{2}} = \sqrt{M{A}^{2} + A{B}^{2} + B{C}^{2}} = \sqrt{{t}^{2} + 8} $ ,
$ {DO} = {BO} = \dfrac{{MB} \cdot  {BC}}{MC} = \dfrac{2\sqrt{{t}^{2} + 4}}{\sqrt{{t}^{2} + 8}} $ ,而 $ {BD} = 2\sqrt{2} $ ,在 $ {\Delta BOD} $ 中,由余弦定理得：
$ \cos \angle {BOD} = \dfrac{2{\left( \dfrac{2\sqrt{{t}^{2} + 4}}{\sqrt{{t}^{2} + 8}}\right) }^{2}}{2{\left( \dfrac{2\sqrt{{t}^{2} + 4}}{\sqrt{{t}^{2} + 8}}\right) }^{2}} =  - \dfrac{4}{{t}^{2} + 4} =  - \dfrac{1}{3} $ ,解得 $ t = 2\sqrt{2} $ ,
所以当 $ {MA} $ 的值为 $ 2\sqrt{2} $ 时,二面角 $ B - {MC} - D $ 的余弦值为 $ - \dfrac{1}{3} $ .

    \end{enumerate}
\end{jieda}

    
\end{enumerate}

C 组:
\begin{enumerate}

\item 在直三棱柱 $ {ABC} - {A}_{1}{B}_{1}{C}_{1} $ 中， $ {A{A}_{1}} = 1,{{A}_{1}C} = 2,{\angle {BAC}} = \dfrac{\pi }{6},E,F $ 为线段 $ {A}_{1}C $ 的三等分点,点 $ D $ 在线段 $ {EF} $ 上 (包括端点) 运动,则二面角 $ D - {AB} - C $ 的正弦值的取值范围为\blkda{$[\dfrac{1}{2},\dfrac{4\sqrt{19}}{19}]$}.

\item 已知在矩形 $ {ABCD} $ 中， $ {AB} = 2 $ ， $ {BC} = \sqrt{2} $ ， $ P $ 为 $ {AB} $ 的中点，将 $ \bigtriangleup  {ADP} $ 沿 $ {DP} $ 翻折，得到四棱锥 $ {A}_{1} - {BCDP} $ ，则二面角 $ {A}_{1} - {DC} - B $ 的余弦值最小是\blkda{$\dfrac{\sqrt{2}}{2}$}.



\end{enumerate}


\clearpage

\subsection*{平面与平面的位置关系（4）}

\begin{enumerate}

\item 已知直线 $ l $ ， $ m $ 和平面 $ a $ ， $ \beta $ ，且 $ l \bot  a $ ， $ m \subset  \beta $ ，则下列命题中正确的是\blkda{(1)(4)}. ①若 $ a//\beta $ ，则 $ l \bot  m $ ②若 $ l \bot  m $ ，则 $ a//\beta $ ③若 $ a \bot  \beta $ ，则 $ l//m $ ④若 $ l//m $ ，则 $ a \bot  \beta $

\item 在长方体 $ {ABCD} - {A}_{1}{B}_{1}{C}_{1}{D}_{1} $ 中，过直线 $ A{C}_{1} $ 的平面交直线 $ B{B}_{1} $ 于点 $ E $ ，交直线 $ D{D}_{1} $ 于点 $ F $ ，则四边形 $ {AE}{C}_{1}F $ 的形状为\blkda{平行四边形}.

\item 已知三棱锥 $ D - {ABC} $ 的三个侧面与底面全等，且 $ {AB} = {AC} = \sqrt{3},{BC} = 2 $ ，则二面角 $ D - {BC} - A $ 的余弦值为\blkda{$0$}.


\item 如图，在长方体 $ {ABCD} - {A}_{1}{B}_{1}{C}_{1}{D}_{1} $ 中， $ {AB} = {BC} = \dfrac{\sqrt{6}}{3}{C{C}_{1}} $ ，则二面角 $ {C}_{1} - {DB} - C $ 的大小为\blkda{$\dfrac{\pi}{3}$}.

\begin{center}

    \includegraphics[width=0.25\linewidth]{bodyxb//src/ogh'sdofhg'sdgsda.png}
\end{center}

\item 直二面角 $ a - l - \beta $ 的棱 $ l $ 上有一点 $ A $ ，在平面 $ a,\beta $ 内各有一条射线 $ {AB},{AC} $ 与 $ l $ 均成 $ {45}^{ \circ  } $ ，则 $ \angle {BAC} = $ \blkda{$\dfrac{\pi}{3}$或$\dfrac{2\pi}{3}$}.

\item 若等边三角形 $ {{ABC}} $ 的边长为 1，BC 上的高为 AD，沿高 AD 折成直二面角， 则 A 到 BC 的距离是\blkda{$\dfrac{\sqrt{14}}{4}$}.


\item 已知 $ \bigtriangleup  {{ABC}} $ 和平面 $ \alpha ,\angle {A} = {30}^{ \circ  },\angle {B} = {60}^{ \circ  },{{AB}} = 2,{{AB}} \subset  \alpha $ ， 且平面 $ {ABC} $ 与 $ \alpha $ 所成角为 $ {30}^{ \circ  } $ ，则点 $ C $ 到平面 $ \alpha $ 的距离为\blkda{$\dfrac{\sqrt{3}}{4}$}.


\item 河堤斜面与水平面所成的二面角为 $ {60}^{ \circ  } $ ，堤面上有一条直道 CD，它与坡脚的水平线的夹角为 $ {30}^{ \circ  } $ ，沿这条直道从坡脚向上行走 $ {10}\sqrt{3}m $ 时人升高了\blkda{$\dfrac{15}{2}$m}

\item 已知 $ a $ ， $ \beta $ 是不重合的平面， $ m $ ， $ n $ 是不重合的直线，下列命题中正确的是\dotfill{(\kh{D})}

\xx{若 $ m//a,m//\beta $ ，则 $ a//\beta $}
{若 $ m//n,m//a $ ，则 $ n//a $}
{若 $ m \bot  a,m \bot  \beta $ ，则 $ a \bot  \beta $} 
{若 $ m \bot  a,m//\beta $ ，则 $ a \bot  \beta $}


\item 已知 $ l $ 为空间的一条直线， $ a $ ， $ \beta $ 为空间的两个不同平面，则下列命题错误的是
\dotfill{(\kh{A})}

\xx{若 $ l//a,l \subset  \beta $ ，则 $ a//\beta $}
{若 $ l \bot  a,l \subset  \beta $ ,则 $ a \bot  \beta $}
{若 $ a//\beta ,l \subset  a $ ,则 $ l//\beta $}
{若 $ a \bot  \beta ,a \cap  \beta  = b,l \bot  b,l \subset  a $ ,则 $ l \bot  \beta $}



\item 如图，在四棱锥 $ P - {ABCD} $ 中，底面 $ {ABCD} $ 是直角梯形, $ {AD}//{BC},{AB} \bot  {BC},{AB} = {BC} = 1,{AD} = 2 $ , $ {PA} \bot $ 平面 $ {ABCD},{PA} = 2 $ .

\begin{enumerate}
    \item 求证： $ {CD}\bot $ 平面 $ {PAC} $ ；
    \item 求二面角 $ P - {CD} - A $ 的正弦值.

\end{enumerate}

\begin{center}

    \includegraphics[width=0.25\linewidth]{bodyxb//src/fahsfhasfasfasg.png}
\end{center}

\begin{jieda}
    
\begin{enumerate}
    \item 因为 $ {AB} = {BC} = 1 $ , $ {AB}\bot {BC} $ ,所以 $ {AC} = \sqrt{2} $ , $ \angle {BAC} = {45}^{ \circ  } $ .
又 $ {AB}\bot {BC},{AD} //  {BC} $ ,所以 $ \angle {DAC} = {45}^{ \circ  } $ .
所以 $ C{D}^{2} = C{A}^{2} + D{A}^{2} - {2CA} \cdot  {DA} \cdot  \cos \angle {DAC} = 2 + 4 - 2 \times  \sqrt{2} \times  2 \times  \dfrac{\sqrt{2}}{2} = 2 $ .
所以 $ {CD} = \sqrt{2} $ .
因为 $ {2}^{2} = {\left( \sqrt{2}\right) }^{2} + {\left( \sqrt{2}\right) }^{2} $ ,即 $ A{D}^{2} = A{C}^{2} + C{D}^{2} $ ,
所以 $ \bigtriangleup  {ACD} $ 为直角三角形,且 $ {CD} \bot  {CA} $ .
又 $ {PA}\bot $ 平面 $ {ABCD},{CD} \subset $ 平面 $ {ABCD} $ ,所以 $ {PA}\bot {CD} $ .
$ {PA},{CA} \subset $ 平面 $ {PAC},{PA} \cap  {CA} = A $ ,所以 $ {CD} \bot $ 平面 $ {PAC} $ .
    \item 因为 $ {CD}\bot $ 平面 $ {PAC},{CA},{CP} \subset $ 平面 $ {PAC} $ ,
所以 $ {CD}\bot {CP},{CD}\bot {CA} $ .
所以 $ \angle {PCA} $ 即为二面角 $ P - {CD} - A $ 的平面角.
在Rt $ {\Delta PAC} $ 中, $ {PA}\bot {AC},{PA} = 2,{AC} = \sqrt{2} $ ,所以 $ {PC} = \sqrt{{PA}^{2} + {AC}^{2}} = \sqrt{6} $ ,
所以 $ \sin \angle {PCA} = \dfrac{PA}{PC} = \dfrac{2}{\sqrt{6}} = \dfrac{\sqrt{6}}{3} $ .
即二面角 $ P - {CD} - A $ 的正弦值为 $ \dfrac{\sqrt{6}}{3} $ .


\end{enumerate}
\end{jieda}

\item 过空间一点P的三条直线 PA、PB、PC 不共面，且两两成 $ {60}^{ \circ  } $ 角，求二面角 $ {A} - {{PB}} - {C} $ 的大小.

\begin{jieda}
    根据题意, $ \angle {APB} = \angle {BPC} = \angle {CPA} = {60}^{ \circ  }, $ ,
设 $ {PA} = {PB} = {PC} = a $ ,利用余弦定理计算 $ {AB} $ 、
$ {BC}\text{、}{AC} $ 的长度,得到 $ {AB} = {BC} = {AC} = a $ 。
取 $ {BP} $ 中点 $ E $ ,因为 $ \bigtriangleup  {ABP} $ , 和 $ \bigtriangleup  {BCP} $ , 是等边三角形
,所以 $ {AE} = {CE} = \dfrac{\sqrt{3}}{2}a, $ 。
在 $ \bigtriangleup  {AEC}, $ 中,根据余弦定理计算 $ \cos \angle {AEC}, $ 的值,
得到 $ \cos \angle {AEC} = \dfrac{1}{3} $ 。
答：二面角 $ A - {PB} - C $ 的大小为 $ \arccos \dfrac{1}{3}, $ 。
\end{jieda}



\item 已知直三棱柱 $ {ABC} - {A}_{1}{B}_{1}{C}_{1} $ (侧棱和地面垂直，上下底面是全等的三角形), 所有棱长均为 2,求二面角 $ A - {BC} - {A}_{1} $ 的余弦值大小.

\begin{jieda}
    解：取BC中点O,连结AO、
$ {\mathrm{A}}_{1}\mathrm{O} $ ,
$ \because $ 直三棱柱ABC- $ {\mathrm{A}}_{1}{\mathrm{B}}_{1}{\mathrm{C}}_{1} $ ,所有
棱长均为 2,
$ \therefore \mathrm{{AO}}\bot \mathrm{{BC}},{\mathrm{A}}_{1}\mathrm{O}\bot \mathrm{{BC}} $ ,
$ \therefore \angle {A}_{1}{OA} $ 是二面角 $ A - {BC} - {A}_{1} $ 的
平面角,
$ {AO} = \sqrt{4 - 1} = \sqrt{3},{A}_{1}O = \sqrt{3 + 4} = $
$ \sqrt{7} $ ,
$ \therefore \cos \angle {A}_{1}{OA} = \dfrac{AO}{{A}_{1}O} = \dfrac{\sqrt{3}}{\sqrt{7}} = \dfrac{\sqrt{21}}{7} $ .
$ \therefore $ 二面角 $ \mathrm{A} - \mathrm{{BC}} - {\mathrm{A}}_{1} $ 的余弦值为$\dfrac{\sqrt{21}}{7}$.

\begin{center}
    \includegraphics[width=0.3\textwidth]{bodyxb//src/hoijupijuoih'iugiyguftoyrdst.png}
\end{center}

\end{jieda}

\item 如图，在四棱锥 $ P - {ABCD} $ 中，底面 $ {ABCD} $ 是正方形，侧面 $ {PAD} \bot $ 底面 $ {ABCD} $ ， 且 $ {PA} = {PD} = \dfrac{\sqrt{2}}{2}{AD} $ ，若 $ E $ 、 $ F $ 分别为 $ {PC} $ 、 $ {BD} $ 的中点.

\begin{enumerate}
    \item 求证： $ {EF}// $ 平面 $ {PAD} $ ；
    \item 求证：平面 $ {ABP}\bot $ 平面 $ {CDP} $ .
\end{enumerate}

\begin{center}

    \includegraphics[width=0.25\linewidth]{bodyxb//src/gijsdohg;dgas.png}
\end{center}
\begin{jieda}
    \begin{enumerate}
        \item 证明：连接AC,
        
        \begin{center}
    \includegraphics[width=0.3\textwidth]{bodyxb//src/sdigjh'aouhrfaoiehFAI.png}
\end{center}

        
        因为四边形 $ {ABCD} $ 为正方形,且 $ F $ 为 $ {BD} $ 的中点,所以 $ F $ 为 $ {AC} $ 的中点,
又因为 $ E $ 为 $ {PC} $ 的中点,则 $ {EF}//{PA} $ ,
$ \because {EF} \not \subset$ 平面 $ {PAD},{PA} \subset $ 平面 $ {PAD},\therefore {EF}// $ 平面 $ {PAD} $ .
        \item 证明：因为四边形 $ {ABCD} $ 为正方形,则 $ {CD}\bot {AD} $ ,
因为平面 $ {PAD} \bot $ 平面 $ {ABCD} $ ,平面 $ {PAD} \cap $ 平面 $ {ABCD} = {AD},{CD} \subset $ 平面 $ {ABCD} $ ,
$ \therefore {CD}\bot $ 平面 $ {PAD} $ ,
$ \because {PA} \subset $ 平面 $ {PAD},\therefore {PA} \bot  {CD} $ ,
$ \because {PA} = {PD} = \dfrac{\sqrt{2}}{2}A{D}^{\prime } $ 所以 $ P{A}^{2} + P{D}^{2} = A{D}^{2} $ ,则 $ {PA} \bot  {PD} $ ,
$ \because {CD} \cap  {PD} = D,{CD} \cap  {PD} \subset $ 平面 $ {CDP} $ ,
$ \therefore {PA}\bot $ 平面 $ {CDP} $ ,又 $ {PA} \subset $ 平面 $ {ABP} $ ,
$ \therefore $ 平面 $ {ABP} \bot $ 平面 $ {CDP} $ .
    \end{enumerate}
\end{jieda}



\item 如图所示，在四棱锥 $ P - {ABCD} $ 中，底面 $ {ABCD} $ 是 $ \angle {DAB} = {60}^{ \circ  } $ 且边长为 $ a $ 的菱形，侧面 $ {PAD} $ 为正三角形，其所在的平面垂直于底面 $ {ABCD} $ . 
\begin{enumerate}
    \item 若 $ G $ 为 $ {AD} $ 边的中点，求证： $ {BG}\bot $ 平面 $ {PAD} $ ； 
    \item 求证： $ {AD}\bot {PB} $ ;
    \item 求二面角 $ A - {BC} - P $ 的大小；\item 若 $ E $ 为 $ {BC} $ 边的中点，能否在棱 $ {PC} $ 上找一点 $ F $ ,使得平面 $ {DEF}\bot $ 平面 $ {ABCD} $ ? 并证明你的结论.
\end{enumerate}

\begin{center}

    \includegraphics[width=0.25\linewidth]{bodyxb//src/'doigjsdigjh'gadg.png}
\end{center}

\begin{jieda}
    \begin{enumerate}
        \item 在底面菱形 $ {ABCD} $ 中, $ \angle {DAB} = {60}^{ \circ  } $ , $ G $ 为 $ {AD} $ 边的中点,所以 $ {BG}\bot {AD} $ ,
又平面 $ {PAD} \bot $ 平面 $ {ABCD} $ ,平面 $ {PAD} \cap $ 平面 $ {ABCD} = {AD} $ ,
所以 $ {BG} \bot $ 平面 $ {PAD} $ .
        \item 连接 $ {PG} $ ,因为 $ \bigtriangleup {PAD} $ 为正三角形, $ G $ 为 $ {AD} $ 边的中点,
得 $ {PG}\bot {AD} $ ,由(1)知 $ {BG}\bot {AD} $ ,
$ {PG} \subset $ 平面 $ {PGB},{BG} \subset $ 平面 $ {PGB},{PG} \cap  {BG} = G $ ,
所以 $ {AD} \bot $ 平面 $ {PGB} $ ,因为 $ {PB} \subset $ 平面 $ {PGB} $ .
所以 $ {AD}\bot {PB} $ .
        \item 由(2)可得 $ {PB}\bot {AD},{BG}\bot {AD} $ ,
$ \because {AD}// {BC} $ ,所以 $ {PB}\bot {BC},{BG}\bot {BC} $ ,
所以 $ \angle {PBG} $ 为二面角 $ A - {BC} - P $ 的平面角
因为 $ {PG} = {BG} = \dfrac{\sqrt{3}}{2}a $ ,所以 $ \angle {PBG} = {45}^{ \circ  } $ ；
        \item 当 $ F $ 为 $ {PC} $ 边的中点时,满足平面 $ {DEF}\bot $ 平面 $ {ABCD} $ ,证明如下：
取 $ {PC} $ 的中点 $ F $ ,连接 $ {DE} $ 、 $ {EF} $ 、 $ {DF} $ ,
在 $ \bigtriangleup  {PBC} $ 中, $ {FE}\left\|  {PB}\right\|  ,{FE} \not \subset$ 平面 $ {PGB},{PB} \subset $ 平面 $ {PGB} $
$ \therefore {FE} // $ 平面 $ {PGB} $
在菱形 $ {ABCD} $ 中, $ {DG}// {BE} $ 且 $ {DG} = {BE} $
$ {BEDG} $ 为平行四边形,则 $ {DE}// {BG},{DE} \not \subset$ 平面 $ {PGB},{BG} \subset $ 平面 $ {PGB} $
$ \therefore {DE} // $ 平面 $ {PGB} $
$ {EF} \cap  {DE} = E $ ,所以平面 $ {DEF}\square $ 平面 $ {PGB} $ ,
因为 $ {BG}\bot $ 平面 $ {PAD} $ ,所以 $ {BG}\bot {PG} $ ,
又因为 $ {PG}\bot {AD},{AD} \cap  {BG} = G $ ,
$ \therefore {PG} \bot $ 平面 $ {ABCD} $ ,而 $ {PG} \subset $ 平面 $ {PGB} $ ,
所以平面 $ {PGB} \bot $ 平面 $ {ABCD} $ ,
所以平面 $ {DEF}\bot $ 平面 $ {ABCD} $ .
    \end{enumerate}
\end{jieda}


\end{enumerate}

\section*{异面直线间的距离}

A 组:

\begin{enumerate}

\item 空间和两条异面直线同时都垂直且相交的直线\dotfill{(\kh{B})}

\xx{不一定存在} {有且只有 1 条} {有 1 条或不存在} {有无数条}

\item 在正方体 $ {ABCD} - {A}_{1}{B}_{1}{C}_{1}{D}_{1} $ 中，棱长为 $ a $ ，所在直线与 $ {AD} $ 成异面直线且距离为 $ a $ 的棱有\blkda{4}条.

\item 长方体 $ {{ABCD}} - {{A}}_{1}{{B}}_{1}{{C}}_{1}{{D}}_{1} $ 中， $ {{AB}} $ 和 $ {D{D}_{1}} $ 的公垂线段是\blkda{AD}. $ {A}{{A}}_{1} $ 和 $ {B}{{C}}_{1} $ 的公垂线段是\blkda{AB}.

\item 已知正方体 $ {ABCD} - {A}_{1}{B}_{1}{C}_{1}{D}_{1} $ 的棱长为 $ a $ ，异面直线 $ {DB} $ 与 $ C{C}_{1} $ 之间的距离为\blkda{$\dfrac{\sqrt{2}a}{2}$}.

\item 正方体 $ {{ABCD}} - {{A}}_{1}{{B}}_{1}{{C}}_{1}{{D}}_{1} $ 中，边长为 4，则异面直线 $ {{A}}_{1}{D} $ 与 $ {{B}}_{1}{{D}}_{1} $ 的距离为\blkda{$\dfrac{4\sqrt{3}}{3}$}.

    
\end{enumerate}


B 组:

\begin{enumerate}

\item 空间四边形 $ {{ABCD}} $ 中， $ {{AB}} = {{BD}} = {{AD}} = 2,{{BC}} = {{CD}} = {3} $ ,延长 $ {{BC}} $ 到 $ {E} $ ,使得 $ {{BC}} = {{CE}},{F} $ 为 $ {{BD}} $ 中点,则异面直线 $ {{AF}} $ 和 $ {{DE}} $ 的距离为\blkda{1}.

\item 边长为 1 的两个正方形 $ {ABCD} $ 和 $ {CDEF} $ 构成大小为 $ \dfrac{2}{3}\pi $ 的二面角，则异面直线 $ {AE} $ 和 $ {CD} $ 之间的距离为\blkda{$\dfrac{1}{2}$}.

\item 四面体 $ {SABC} $ 中， $ {SA} = {BC} = {13} $ ， $ {SB} = {AC} = {14} $ ， $ {SC} = {AB} = {15} $ ， 则异面直线 $ {A}{S} $ 与 $ {B}{C} $ 的距离为\blkda{$3\sqrt{14}$}.

\begin{jieda}
    【补形】将四面体 $ {SABC} $ 补成长方体 $ {SDBG} - {EAFC} $ ,连接 $ {DE} $ 交 $ {AS} $ 于点 $ M $ ,连接 $ {FG} $ 交 $ {BC} $ 于点 $ N $ ,连接 $ {MN} $ ,
则 $ M $ 、 $ N $ 分别为 $ {DE} $ 、 $ {BC} $ 的中点,


    \begin{center}
    \includegraphics[width=0.3\textwidth]{bodyxb//src/auhf'ahfafsfafa.png}
\end{center}
    由已知可得 $ \left\{  \begin{array}{l} S{A}^{2} = S{D}^{2} + S{E}^{2} = {169} \\  S{B}^{2} = S{D}^{2} + S{G}^{2} = {196} \\  S{C}^{2} = S{E}^{2} + S{G}^{2} = {225} \end{array}\right. $ ,可得 $ {SG} = 3\sqrt{14} $ ,
因为 $ {BD}//{CE} $ 且 $ {BD} = {CE} $ ,故四边形 $ {BDE}C $ 为平行四边形,则 $ {BC}//{DE} $ 且 $ {BC} = {DE} $ ,
又因为 $ M $ 、 $ N $ 分别为 $ {DE} $ 、 $ {BC} $ 的中点,所以, $ {DM}//{BN} $ 且 $ {DM} = {BN} $ ,
故四边形 $ {BDMN} $ 为平行四边形,故 $ {MN}//{BD} $ 且 $ {MN} = {BD} = {SG} = 3\sqrt{14} $ ,
$ \because {BD} \bot $ 平面 $ {SDAE},{AS} \subset $ 平面 $ {SDAE},\therefore {BD} \bot  {AS} $ ,即 $ {MN} \bot  {AS} $ ,
同理可得 $ {MN}\bot {BC} $ ,故异面直线 $ {AS} $ 与 $ {BC} $ 的距离为 $ {MN} = 3\sqrt{14} $ .
故答案为: $ 3\sqrt{14} $ .

\end{jieda}


\item 如图，在棱长为 2 的正方体 $ {ABCD} - {A}_{1}{B}_{1}{C}_{1}{D}_{1} $ 中， $ E $ 为 $ {BC} $ 的中点,点 $ P $ 在线段 $ {D}_{1}E $ 上,点 $ P $ 到直线 $ C{C}_{1} $ 的距离的最小值为\blkda{$\dfrac{2\sqrt{5}}{5}$}.

\begin{center}

    \includegraphics[width=0.25\linewidth]{bodyxb//src/gsddlihjgsoigo'aedgs.png}
\end{center}



\item 已知菱形 $ {ABCD} $ 中， $ \angle {BAD} = \dfrac{\pi }{3} $ ，沿对角线 $ {BD} $ 折起，使二面角 $ A - {BD} - C $ 的平面角为 $ \theta $ ，若异面直线 $ {AC} $ 与 $ {BD} $ 的距离是菱形边长的 $ \dfrac{3}{4} $ ，则 $ \theta  = $ \blkda{$\dfrac{\pi}{3}$}.

\begin{jieda}

如图,设菱形的边长为 $ {2a} $ ,连接两条对角线 $ {AC} \cap  {BD} = O $
易得 $ {AC} \bot  {BD},{AO} = {OC} = \sqrt{3}a $

    \begin{center}
    \includegraphics[width=0.3\textwidth]{bodyxb//src/aihfahfoaofhafa.png}
\end{center}


    菱形 $ {ABCD} $ 沿对角线 $ {BD} $ 折起,连接 $ {AC} $ ,得到三棱锥 $ A - {BCD} $
在菱形 $ {ABCD} $ 中, $ {AC} \bot  {BD} $ ,翻着后垂直不变,即 $ \left\{  \begin{array}{l} {AO} \bot  {BD} \\  {CO} \bot  {BD} \end{array}\right. $
所以 $ \angle {AOC} $ 是二面角 $ A - {BD} - C $ 的平面角,即 $ \angle {AOC} = \theta $
又因为 $ \left\{  \begin{array}{l} {AO} \bot  {BD} \\  {CO} \bot  {BD} \\  {AO} \cap  {CO} = O \\  {AO},{CO} \subset  \text{ 面 }{AOC} \end{array}\right. $
所以 $ {BD} \bot $ 平面 $ {AOC} $ ,取 $ {AC} $ 中点 $ M $ ,连接 $ {OM} $
又因为 $ {OM} \subset $ 平面 $ {AOC} $
所以 $ {OM} \bot  {BD} $
在 $ \bigtriangleup  {AOC} $ 中, $ {AO} = {OC} = \sqrt{3}a $ ,并且 $ M $ 为 $ {AC} $ 的中点,
所以 $ {OM} \bot  {AC} $
故 $ {OM} $ 是异面直线 $ {AC} $ 与 $ {BD} $ 的距离
又因为异面直线 $ {AC} $ 与 $ {BD} $ 的距离是菱形边长的 $ \dfrac{3}{4} $
所以 $ {OM} = \dfrac{3}{4} \times  {2a} = \dfrac{3}{2}a $
在 $ \mathrm{R}{t\Delta AOM} $ 中, $ \angle {AOM} = \dfrac{\theta }{2} $
所以 $ \cos \dfrac{\theta }{2} = \dfrac{OM}{AO} = \dfrac{\dfrac{3}{2}a}{\sqrt{3}a} = \dfrac{\sqrt{3}}{2} $ ,又因为 $ \dfrac{\theta }{2} \in  \left( {0,\dfrac{\pi }{2}}\right) $
所以 $ \dfrac{\theta }{2} = \dfrac{\pi }{6}\therefore \theta  = \dfrac{\pi }{3} $

\end{jieda}


\item 下列命题其中是真命题的有 \blkda{(2)(3)}

\begin{enumerate}
    \item 两条异面直线的公垂线有无数条;
    \item 异面直线之间的距离就是两条异面直线上点之间距离的最小值;
    \item 过两条异面直线中的一条有且只有一个平面与另一条直线平行.
\end{enumerate}



\item 正方体 $ {{ABCD}} - {{A}}_{1}{{B}}_{1}{{C}}_{1}{{D}}_{1} $ 中，边长为 1，与直线 $ {{A}}_{1}{D} $ 异面且距离为 $ \dfrac{\sqrt{3}}{3} $ 的面对角线有 \blkda{4}条.



\item 若直线 $ a $ 和已知直线 $ b $ 满足：① $ a $ 、 $ b $ 是异面直线，② $ a $ 、 $ b $ 之间的距离是定值，③a与 $ b $ 所成的角也是定值，则直线 $ a $ \dotfill{(\kh{D})}

\xx{仅有一条}{有两条} {有四条}{有无数条}

\item 空间四边形 $ {{ABCD}} $ 中， $ {{AB}} = {{CD}} = {4},{{BC}} = {{AD}} = {3},{{AD}} \bot  {{DC}} $ , $ {{AD}} \bot  {{BC}},{{AB}} \bot  {{BC}} $ ,求异面直线 $ {{AD}} $ 和 $ {{BC}} $ 的距离.


\begin{jieda}
    【鳖臑，补形】$\sqrt{7}$

    \begin{center}
    \includegraphics[width=0.3\textwidth]{bodyxb//src/gsdijgoahdfgohahf.png}
\end{center}
\end{jieda}


\item 所有棱长都为 1 的四面体 $ {{ABCD}} $ 中，找到异面直线 $ {{AB}} $ 和 $ {{CD}} $ 的公垂线， 求出 $ {{AB}} $ 和 $ {{CD}} $ 的距离.
\begin{jieda}

如图,取 $ {AB} $ 中点 $ E,{CD} $ 中点 $ F $ ,连接 $ {BF},{AF},{EF} $ ,
由已知 $ {AF} = {BF} = \sqrt{{1}^{2} - {\left( \dfrac{1}{2}\right) }^{2}} = \dfrac{\sqrt{3}}{2}，\therefore {EF} \bot  {AB} $ ,同理 $ {EF} \bot  {CD} $ ,
所以 $ {EF} $ 是 $ {AB} $ 与 $ {CD} $ 的公垂线, $ {AB} $ 与 $ {CD} $ 的距离即为线段 $ {EF} $ 的长.
且 $ {EF} = \sqrt{A{F}^{2} - A{E}^{2}} = \sqrt{\dfrac{3}{4} - \dfrac{1}{4}} = \dfrac{\sqrt{2}}{2} $ ,
所以 $ {AB} $ 与 $ {CD} $ 的距离即为 $ \dfrac{\sqrt{2}}{2} $ .

    \begin{center}
    \includegraphics[width=0.3\textwidth]{bodyxb//src/fa'hfaouhfiahfouha.png}
\end{center}
\end{jieda}

\item 正方体 $ {{ABCD}} - {{A}}_{1}{{B}}_{1}{{C}}_{1}{{D}}_{1} $ 中，边长为 2，求异面直线 $ {{CD}} $ 与 $ {A}{{C}}_{1} $ 的距离.

\begin{jieda}
    异面直线的距离即为BC到面$ABC_1D_1$的距离，答案为$\sqrt{2}$。
\end{jieda}

\item 已知 $ {S} $ 是矩形 $ {{ABCD}} $ 所在平面外一点， $ {{SA}} \bot  {{BC}},{{BS}} \bot  {{CD}},{{SA}} $ 与 $ {{CD}} $ 所成角的大小为 $ \dfrac{\pi }{3},{{SD}} $ 与 $ {{BC}} $ 所成角大小为 $ \dfrac{\pi }{6},{{SA}} = {1} $ ,分别求直线 $ {{SA}} $ 与 $ {{CD}} $ 的距离及 $ {{SB}} $ 与 $ {{AD}} $ 的距离.

\begin{jieda}

$ \because {AD}//{BC},{AB}//{CD},{SA}\bot {BC},{BS}\bot {CD}\therefore {SA}\bot {AD},{SB}\bot {AB} $
因为 $ {SD} $ 与 $ {BC} $ 所成角大小为 $ \dfrac{\pi }{6} $ ,而 $ {BC}//{AD} $ ,则 $ \angle {SDA} = \dfrac{\pi }{6} $ ,
因为 $ {SA} $ 与 $ {CD} $ 所成角大小为 $ \dfrac{\pi }{3} $ ,而 $ {AB}//{CD} $ ,则 $ \angle {SAB} = \dfrac{\pi }{3} $ ,
$ {SA} = 1 $ ,则 $ {SD} = 2,{AD} = \sqrt{3},{AB} = \dfrac{1}{2} $ ,
又 $ {ABCD} $ 是矩形,
所以线段 $ {AD} $ 是直线 $ {SA} $ 与 $ {CD} $ 的公垂线段,线段 $ {AB} $ 是 $ {SB} $ 与 $ {AD} $ 的公垂线线段,
所以直线 $ {SA} $ 与 $ {CD} $ 的距离是 $ \sqrt{3},{SB} $ 与 $ {AD} $ 的距离是 $ \dfrac{1}{2} $ .
    \begin{center}
    \includegraphics[width=0.3\textwidth]{bodyxb//src/sdoufhauhfOAuhwOFa.png}
\end{center}
\end{jieda}


\item 四棱锥 $ P - {ABCD} $ 中， $ {PA}\bot $ 平面 $ {ABCD} $ ，底面 $ {ABCD} $ 是平行四边形,且 $ \angle {BAD} = \dfrac{2\pi }{3},{AB} = 2,{PA} = 3,{AD} = 4,M $ 是 $ {BC} $ 的中点.

\begin{center}

    \includegraphics[width=0.25\linewidth]{bodyxb//src/dig'sodihgsohdgagfa.png}
\end{center}

\begin{enumerate}
    \item 求二面角 $ A - {PM} - B $ 的余弦值；
    \item 求异面直线 $ {AD} $ 和 $ {PM} $ 之间的距离.
\end{enumerate}


\begin{jieda}

\begin{center}
    \includegraphics[width=0.3\textwidth]{bodyxb//src/lfhajhfAOHFwfa.png}
\end{center}
    \begin{enumerate}
        \item 取 $ {BM} $ 中点 $ E $ ,连接 $ {AE},{PE} $ ,作 $ {AH}\bot {PE} $ ,垂足为 $ H $ ,
再过点 $ {A作AG} \bot  {PM} $ ,连接 $ {HG} $ ,
根据题意可知 $ \bigtriangleup  {ABM} $ 为正三角形,
则 $ {AE} \bot  {BM},{AE} = \sqrt{3},{AM} = 2,\therefore {PM} = {PB} = \sqrt{13},{PE} \bot  {BM} $ ,
又 $ {PE} \cap  {AE} = E,{PE}\text{、}{AE} \subset $ 平面 $ {PAE} $ ,则 $ {BM}\bot $ 平面 $ {PAE} $ ,
因为 $ {AH} \subset $ 平面 $ {PAE} $ ,则 $ {BM} \bot  {AH} $ ,
又 $ {BM} \cap  {PE} = H,{BM}\text{、}{PE} \subset $ 平面 $ {PBM} $ ,
所以 $ {AH} \bot $ 平面 $ {PBM} $ ,
因为 $ {PM},{HG} \subset $ 平面 $ {PBM} $ ,所以 $ {AH} \bot  {PM},{AH} \bot  {HG} $ ,
因为 $ {AG} \cap  {AH} = A,{AG}\text{、}{AH} \subset $ 平面 $ {AGH} $ ,
所以 $ {PM} \bot $ 平面 $ {AGH} $ ,
因为 $ {HG} \subset $ 平面 $ {AGH} $ ,所以 $ {HG} \bot  {PM} $ ,
所以 $ \angle {AGH} $ 为二面角 $ A - {PM} - B $ 的平面角,
在Rt $ \bigtriangleup  {APM} $ 中, $ {AG} = \dfrac{{AM} \cdot  {AP}}{PM} = \dfrac{6}{\sqrt{13}} $ ,
在Rt $ \bigtriangleup {APE} $ 中, $ {AH} = \dfrac{{AE} \cdot  {AP}}{PE} = \dfrac{3\sqrt{3}}{2\sqrt{3}} = \dfrac{3}{2} $ ,
在Rt $ {\Delta AGH} $ 中, $ \cos \angle {AGH} = \dfrac{GH}{AG} = \dfrac{\sqrt{A{G}^{2} - A{H}^{2}}}{AG} = \dfrac{\sqrt{\dfrac{36}{13} - \dfrac{9}{4}}}{\dfrac{6}{\sqrt{13}}} = \dfrac{\sqrt{3}}{4} $ ,
所以二面角 $ A - {PM} - B $ 的余弦值为 $ \dfrac{\sqrt{3}}{4} $ .

\item 根据底面 $ {ABCD} $ 是平行四边形,所以 $ {AD}//{BC} $ ,
因为 $ {BC} \subset $ 平面 $ {PBC},{AD} \not \subset$ 平面 $ {PBC} $ ,
故 $ {AD}// $ 平面 $ {PBC} $ ,
所以线段 $ {AH} $ 的长度即为直线 $ {AD} $ 与平面 $ {PBC} $ 间的距离,
也即异面直线 $ {AD} $ 和 $ {PM} $ 之间的距离.
由上可知 $ {AH} = \dfrac{3}{2} $ ,所以异面直线 $ {AD} $ 和 $ {PM} $ 之间的距离为 $ \dfrac{3}{2} $ .
    \end{enumerate}

    
\end{jieda}

\end{enumerate}


C 组:
\begin{enumerate}

\item 已知在正方体 $ {ABCD} - {A}_{1}{B}_{1}{C}_{1}{D}_{1} $ 中， $ E,F,G $ 分别是棱 $ {D}_{1}{C}_{1},{AB},{BC} $ 的中点， $ H $ 是棱 $ {CD} $ 上一点，则下列命题中正确的个数为\blkda{3}
\begin{enumerate}
    \item 异面直线 $ {D}_{1}H $ 与 $ {AB} $ 之间的距离为定值；

    \item 平面 $ {EFG}// $ 平面 $ A{D}_{1}H $ ；

    \item 设平面 $ {A}_{1}{AH} \cap $ 平面 $ {A}_{1}{B}_{1}{C}_{1}{D}_{1} = l $ ，则 $ {AH}//l $ ；

    \item 直线 $ {A}_{1}{C}_{1} $ 与平面 $ {AB}{C}_{1}{D}_{1} $ 所成的角为 $ {30}^{ \circ  } $ .
\end{enumerate}

\begin{jieda}
    \begin{center}
    \includegraphics[width=0.3\textwidth]{bodyxb//src/alhfaouhfaohwfawef.png}
\end{center}
    对于①,过 $ H $ 点作 $ {AD} $ 的平行线,与 $ {AB} $ 相交于点 $ M $ ,则 $ {HM} $ 为异面直线 $ {D}_{1}H $ 与 $ {AB} $ 的公垂线段,且长度等于正方体棱长,所以异面直线 $ {D}_{1}H $ 与
$ {AB} $ 之间的距离为定值,①正确；
对于②,平面 $ {ABCD} \cap $ 平面 $ {EFG} = {FG} $ ,平面 $ {ABCD} \cap $ 平面 $ A{D}_{1}H = {AH} $ ,若平面 $ {EFG} // $ 平面 $ A{D}_{1}H $ ,则有 $ {FG}//{AH} $ ,而 $ {FG}//{AH} $ 不一定
成立,②错误；
对于③,平面 $ {ABCD}// $ 平面 $ {A}_{1}{B}_{1}{C}_{1}{D}_{1} $ , $ {AH} \subset $ 平面 $ {ABCD} $ ,所以 $ {AH}// $ 平面 $ {A}_{1}{B}_{1}{C}_{1}{D}_{1} $ ,
$ {AH} \subset $ 平面 $ {A}_{1}{AH} $ ,平面 $ {A}_{1}{AH} \cap $ 平面 $ {A}_{1}{B}_{1}{C}_{1}{D}_{1} = l $ ,则 $ {AH}// l $ ,③正确；
对于④, $ N $ 为 $ A{D}_{1} $ 的中点,连接 $ {A}_{1}N,{C}_{1}N $
平面 $ {AB}{C}_{1}{D}_{1} \bot $ 平面 $ {AD}{D}_{1}{A}_{1} $ ,平面 $ {AB}{C}_{1}{D}_{1} \cap $ 平面 $ {AD}{D}_{1}{A}_{1} = A{D}_{1},{A}_{1}N \subset $ 平面 $ {AD}{D}_{1}{A}_{1},{A}_{1}N \bot  A{D}_{1} $ ,
所以 $ {A}_{1}N \bot $ 平面 $ {AB}{C}_{1}{D}_{1} $ ,直线 $ {A}_{1}{C}_{1} $ 与平面 $ {AB}{C}_{1}{D}_{1} $ 所成的角为 $ \angle {A}_{1}{C}_{1}N $ ,
在Rt $ \Delta {A}_{1}{C}_{1}N $ 中, $ {A}_{1}N = \dfrac{1}{2}{A}_{1}{C}_{1},\angle {A}_{1}{C}_{1}N = {30}^{ \circ  } $ ,直线 $ {A}_{1}{C}_{1} $ 与平面 $ {AB}{C}_{1}{D}_{1} $ 所成的角为 $ {30}^{ \circ  } $ ,④正确.

\end{jieda}


\item 已知长方体底面边长分别为 2,1 , 正四面体的四个顶点分别在长方体的顶点及长度为 1 的棱所在的直线上, 求长方体的高.

\begin{jieda}

因为任何一个三棱锥均存在唯一的外接平行六面体,只要分别过三组对棱作三组平行平面,
即能围成所要的平行六面体.

    \begin{center}
    \includegraphics[width=0.3\textwidth]{bodyxb//src/alfhaouhfaefae.png}
\end{center}


    反之,任意一个平行六面体也一定内接一个四面体,且长方体对应的四面体一定是对棱相等
的四面体,正方体对应的一定是正四面体.
由于一组对棱 $ \left( {C{B}^{\prime }}\right. $ 与 $ A{D}^{\prime } $ )距离是2,即 $ {AB} = 2 $ ,故另一组对棱 $ \left( {AC}\right. $ 与 $ {B}^{\prime }{D}^{\prime } $ )也为2,
即高 $ B{B}^{\prime } = 2 $ .
 
\end{jieda}



\end{enumerate}